% unit-id: d90.orig.v1.tr03.unit
% target-edition-id: d90.orig.v1.tr03.edition.id-ID
% authorship: independent coursebook completion layer
% license: CC BY-SA 4.0
\chapter{Asesmen, Laboratorium, dan Proyek Penutup}
\label{orig03:chapter:course-closure}

Bab ini menyediakan lapisan penutup kursus yang dapat dipakai untuk belajar
mandiri: peta asesmen kumulatif, diagnostik prasyarat, enam set soal, rubrik
pembuktian, ujian tengah dan akhir, dua laboratorium komputasi, serta satu
proyek kapstone. Setiap soal mempunyai petunjuk bertahap, jawaban singkat, dan
solusi acuan lengkap. Kode laboratorium dan proyek menghasilkan sertifikat
deterministik yang dapat diperiksa tanpa perangkat lunak berpemilik.

Seluruh uraian, soal, solusi, rubrik, kode, dan instans sintetik dalam bab ini
ditulis mandiri dan tersedia berdasarkan CC BY-SA 4.0. Bab ini tidak menyalin
materi O018 tentang program linear atau bilangan bulat, simpleks, sensitivitas
LP, ataupun optimisasi jaringan. Pendamping MIT OCW, Royer, dan Penn tetap
merupakan karya terpisah dengan hak dan atribusinya sendiri.

% ORIG03-SEGMENT-CODE: OR03-S001
% ORIG03-SEGMENT-ID: d90.orig.v1.tr03.seg0001
% ORIG03-MODULE-ID: map.assessment.0001
% SPDX-License-Identifier: CC-BY-SA-4.0
% Karya asli berbahasa Indonesia; tidak menyalin butir dari sumber donor.

\section{Peta asesmen dan kontrak topologi}
\label{orig03:section:peta-asesmen}

Bagian ini memetakan asesmen Original-03 tanpa mengulang bunyi soal. Seluruh
teks baru pada tranche ini ditulis secara independen dan dilisensikan di bawah
CC BY-SA 4.0. Notasi mengikuti spine Habring serta suplemen dan tranche yang
sudah diterima, tetapi tidak ada butir donor yang disalin. Materi program linear,
integer, simpleks, pemodelan riset operasi, dan wilayah O018 lainnya sengaja
tidak dimasukkan.

Kontrak yang dapat diurai mesin adalah sebagai berikut. Setiap butir induk
memiliki ID \texttt{*.problem.NNNN}. Setiap subbutir memiliki ID
\texttt{*.prompt.NNNN} dan tepat satu pasangan
\texttt{*.hint1.NNNN}/\texttt{*.hint2.NNNN}, satu
\texttt{*.answer.NNNN}, serta satu \texttt{*.solution.NNNN}. Komentar
\texttt{ORIG03-ASSESSMENT-MAP-ID} pada peta ini hanya merupakan rujukan
nondeklaratif. Deklarasi otoritatif \texttt{ORIG03-STABLE-ID} dan label LaTeX
yang berpadanan berada tepat pada butir soal di modul 01--10. Tujuh label
\texttt{orig03:rubric:proof:0001} sampai
\texttt{orig03:rubric:proof:0007} disediakan sebagai rujukan logis bagi rubrik
bukti yang dihimpun oleh agregator.

\subsection{Lima puluh empat prompt baru}

\paragraph{Diagnostik prasyarat: 20 prompt.}
\begin{itemize}
% ORIG03-ASSESSMENT-MAP-ID: diag.prompt.0001
\item \texttt{diag.prompt.0001}: norma dasar dan aritmetika vektor.
% ORIG03-ASSESSMENT-MAP-ID: diag.prompt.0002
\item \texttt{diag.prompt.0002}: rantai ketaksamaan norma.
% ORIG03-ASSESSMENT-MAP-ID: diag.prompt.0003
\item \texttt{diag.prompt.0003}: proyeksi pada hiperbidang.
% ORIG03-ASSESSMENT-MAP-ID: diag.prompt.0004
\item \texttt{diag.prompt.0004}: karakterisasi proyeksi konveks.
% ORIG03-ASSESSMENT-MAP-ID: diag.prompt.0005
\item \texttt{diag.prompt.0005}: koordinat kombinasi konveks.
% ORIG03-ASSESSMENT-MAP-ID: diag.prompt.0006
\item \texttt{diag.prompt.0006}: kestabilan kekonveksan terhadap irisan.
% ORIG03-ASSESSMENT-MAP-ID: diag.prompt.0007
\item \texttt{diag.prompt.0007}: gradien dan Hessian kuadratik.
% ORIG03-ASSESSMENT-MAP-ID: diag.prompt.0008
\item \texttt{diag.prompt.0008}: konstanta kuat-konveks dan mulus.
% ORIG03-ASSESSMENT-MAP-ID: diag.prompt.0009
\item \texttt{diag.prompt.0009}: penyangga afin orde pertama.
% ORIG03-ASSESSMENT-MAP-ID: diag.prompt.0010
\item \texttt{diag.prompt.0010}: stasioneritas dan optimalitas global.
% ORIG03-ASSESSMENT-MAP-ID: diag.prompt.0011
\item \texttt{diag.prompt.0011}: subdiferensial nilai mutlak.
% ORIG03-ASSESSMENT-MAP-ID: diag.prompt.0012
\item \texttt{diag.prompt.0012}: ambang lunak skalar.
% ORIG03-ASSESSMENT-MAP-ID: diag.prompt.0013
\item \texttt{diag.prompt.0013}: ekspektasi dan varians diskret.
% ORIG03-ASSESSMENT-MAP-ID: diag.prompt.0014
\item \texttt{diag.prompt.0014}: rata-rata sampel independen.
% ORIG03-ASSESSMENT-MAP-ID: diag.prompt.0015
\item \texttt{diag.prompt.0015}: penyelesaian rekurensi kontraktif.
% ORIG03-ASSESSMENT-MAP-ID: diag.prompt.0016
\item \texttt{diag.prompt.0016}: kompleksitas iterasi geometrik.
% ORIG03-ASSESSMENT-MAP-ID: diag.prompt.0017
\item \texttt{diag.prompt.0017}: kerucut normal dan optimalitas terkendala.
% ORIG03-ASSESSMENT-MAP-ID: diag.prompt.0018
\item \texttt{diag.prompt.0018}: bentuk titik tetap proyeksi.
% ORIG03-ASSESSMENT-MAP-ID: diag.prompt.0019
\item \texttt{diag.prompt.0019}: ortogonalitas dalam Cauchy--Schwarz.
% ORIG03-ASSESSMENT-MAP-ID: diag.prompt.0020
\item \texttt{diag.prompt.0020}: ketaksamaan Young berbobot.
\end{itemize}

\paragraph{Set soal dasar konveks: 12 prompt.}
\begin{itemize}
% ORIG03-ASSESSMENT-MAP-ID: ps01.prompt.0001
\item \texttt{ps01.prompt.0001}: maksimum fungsi afin.
% ORIG03-ASSESSMENT-MAP-ID: ps01.prompt.0002
\item \texttt{ps01.prompt.0002}: subdiferensial fungsi nilai mutlak.
% ORIG03-ASSESSMENT-MAP-ID: ps01.prompt.0003
\item \texttt{ps01.prompt.0003}: minimisasi komposit satu dimensi.
% ORIG03-ASSESSMENT-MAP-ID: ps01.prompt.0004
\item \texttt{ps01.prompt.0004}: spektrum kuadratik konveks.
% ORIG03-ASSESSMENT-MAP-ID: ps01.prompt.0005
\item \texttt{ps01.prompt.0005}: kontraksi eksak langkah gradien.
% ORIG03-ASSESSMENT-MAP-ID: ps01.prompt.0006
\item \texttt{ps01.prompt.0006}: evaluasi numerik Jensen.
% ORIG03-ASSESSMENT-MAP-ID: ps01.prompt.0007
\item \texttt{ps01.prompt.0007}: pembuktian Jensen dua titik.
% ORIG03-ASSESSMENT-MAP-ID: ps01.prompt.0008
\item \texttt{ps01.prompt.0008}: konjugat kuadratik.
% ORIG03-ASSESSMENT-MAP-ID: ps01.prompt.0009
\item \texttt{ps01.prompt.0009}: bikonjugat kuadratik.
% ORIG03-ASSESSMENT-MAP-ID: ps01.prompt.0010
\item \texttt{ps01.prompt.0010}: celah Fenchel--Young.
% ORIG03-ASSESSMENT-MAP-ID: ps01.prompt.0011
\item \texttt{ps01.prompt.0011}: proyeksi pada kotak.
% ORIG03-ASSESSMENT-MAP-ID: ps01.prompt.0012
\item \texttt{ps01.prompt.0012}: bukti pemotongan koordinat.
\end{itemize}

\paragraph{Set soal metode proksimal: 11 prompt.}
\begin{itemize}
% ORIG03-ASSESSMENT-MAP-ID: ps02.prompt.0001
\item \texttt{ps02.prompt.0001}: penurunan rumus ambang lunak.
% ORIG03-ASSESSMENT-MAP-ID: ps02.prompt.0002
\item \texttt{ps02.prompt.0002}: proks nilai mutlak berdimensi tiga.
% ORIG03-ASSESSMENT-MAP-ID: ps02.prompt.0003
\item \texttt{ps02.prompt.0003}: operator iterasi gradien proksimal.
% ORIG03-ASSESSMENT-MAP-ID: ps02.prompt.0004
\item \texttt{ps02.prompt.0004}: dua iterasi dan solusi eksak.
% ORIG03-ASSESSMENT-MAP-ID: ps02.prompt.0005
\item \texttt{ps02.prompt.0005}: lema penurunan untuk gradien Lipschitz.
% ORIG03-ASSESSMENT-MAP-ID: ps02.prompt.0006
\item \texttt{ps02.prompt.0006}: langkah kuadratik tajam.
% ORIG03-ASSESSMENT-MAP-ID: ps02.prompt.0007
\item \texttt{ps02.prompt.0007}: bentuk operator maju--mundur.
% ORIG03-ASSESSMENT-MAP-ID: ps02.prompt.0008
\item \texttt{ps02.prompt.0008}: titik tetap dan minimizer.
% ORIG03-ASSESSMENT-MAP-ID: ps02.prompt.0009
\item \texttt{ps02.prompt.0009}: kontraksi operator maju--mundur.
% ORIG03-ASSESSMENT-MAP-ID: ps02.prompt.0010
\item \texttt{ps02.prompt.0010}: resolven operator linear monoton kuat.
% ORIG03-ASSESSMENT-MAP-ID: ps02.prompt.0011
\item \texttt{ps02.prompt.0011}: laju titik proksimal.
\end{itemize}

\paragraph{Set soal dualitas dan KKT: 11 prompt.}
\begin{itemize}
% ORIG03-ASSESSMENT-MAP-ID: ps03.prompt.0001
\item \texttt{ps03.prompt.0001}: fungsi dual masalah kuadratik berekualitas.
% ORIG03-ASSESSMENT-MAP-ID: ps03.prompt.0002
\item \texttt{ps03.prompt.0002}: sistem KKT dan celah nol.
% ORIG03-ASSESSMENT-MAP-ID: ps03.prompt.0003
\item \texttt{ps03.prompt.0003}: pengali KKT kendala pertidaksamaan.
% ORIG03-ASSESSMENT-MAP-ID: ps03.prompt.0004
\item \texttt{ps03.prompt.0004}: Slater dan dual skalar.
% ORIG03-ASSESSMENT-MAP-ID: ps03.prompt.0005
\item \texttt{ps03.prompt.0005}: KKT proyeksi pada hiperbidang afin.
% ORIG03-ASSESSMENT-MAP-ID: ps03.prompt.0006
\item \texttt{ps03.prompt.0006}: nilai primal dan dual hiperbidang.
% ORIG03-ASSESSMENT-MAP-ID: ps03.prompt.0007
\item \texttt{ps03.prompt.0007}: keunikan dan kecukupan KKT.
% ORIG03-ASSESSMENT-MAP-ID: ps03.prompt.0008
\item \texttt{ps03.prompt.0008}: pengali sebagai sensitivitas nilai.
% ORIG03-ASSESSMENT-MAP-ID: ps03.prompt.0009
\item \texttt{ps03.prompt.0009}: ekspansi eksak fungsi nilai.
% ORIG03-ASSESSMENT-MAP-ID: ps03.prompt.0010
\item \texttt{ps03.prompt.0010}: KKT proyeksi pada bola.
% ORIG03-ASSESSMENT-MAP-ID: ps03.prompt.0011
\item \texttt{ps03.prompt.0011}: Slater dan kecukupan pada bola.
\end{itemize}

\subsection{Empat belas latihan terselesaikan yang sudah diterima}

Daftar berikut hanya memetakan ID dan peran kurikuler; bunyi latihan, petunjuk,
dan solusi tetap berada pada modul asalnya.
\begin{itemize}
% ORIG03-EXISTING-ASSESSMENT-ID: d90.becker.98ed693.b03.exercise.0001
\item \texttt{d90.becker.98ed693.b03.exercise.0001}: ketakbiasan dalam
  reduksi varians.
% ORIG03-EXISTING-ASSESSMENT-ID: d90.becker.98ed693.b03.exercise.0002
\item \texttt{d90.becker.98ed693.b03.exercise.0002}: identitas varians dan
  langkah kuadratik.
% ORIG03-EXISTING-ASSESSMENT-ID: d90.orig.v1.tr01.exercise.0001
\item \texttt{d90.orig.v1.tr01.exercise.0001}: operator proksimal Lasso.
% ORIG03-EXISTING-ASSESSMENT-ID: d90.orig.v1.tr01.exercise.0002
\item \texttt{d90.orig.v1.tr01.exercise.0002}: ketaksamaan satu langkah.
% ORIG03-EXISTING-ASSESSMENT-ID: d90.orig.v1.tr01.exercise.0003
\item \texttt{d90.orig.v1.tr01.exercise.0003}: sampling tanpa penggantian.
% ORIG03-EXISTING-ASSESSMENT-ID: d90.orig.v1.tr01.exercise.0004
\item \texttt{d90.orig.v1.tr01.exercise.0004}: cermin entropi.
% ORIG03-EXISTING-ASSESSMENT-ID: d90.orig.v1.tr01.exercise.0005
\item \texttt{d90.orig.v1.tr01.exercise.0005}: anggaran oracle minibatch.
% ORIG03-EXISTING-ASSESSMENT-ID: d90.orig.v1.tr01.exercise.0006
\item \texttt{d90.orig.v1.tr01.exercise.0006}: varians Prox-SAGA.
% ORIG03-EXISTING-ASSESSMENT-ID: d90.orig.v1.tr02.exercise.0001
\item \texttt{d90.orig.v1.tr02.exercise.0001}: VI dan kerucut normal.
% ORIG03-EXISTING-ASSESSMENT-ID: d90.orig.v1.tr02.exercise.0002
\item \texttt{d90.orig.v1.tr02.exercise.0002}: maksimalitas.
% ORIG03-EXISTING-ASSESSMENT-ID: d90.orig.v1.tr02.exercise.0003
\item \texttt{d90.orig.v1.tr02.exercise.0003}: resolven.
% ORIG03-EXISTING-ASSESSMENT-ID: d90.orig.v1.tr02.exercise.0004
\item \texttt{d90.orig.v1.tr02.exercise.0004}: operator skew.
% ORIG03-EXISTING-ASSESSMENT-ID: d90.orig.v1.tr02.exercise.0005
\item \texttt{d90.orig.v1.tr02.exercise.0005}: rentang langkah.
% ORIG03-EXISTING-ASSESSMENT-ID: d90.orig.v1.tr02.exercise.0006
\item \texttt{d90.orig.v1.tr02.exercise.0006}: bayangan Douglas--Rachford.
\end{itemize}

Dengan demikian, peta awal ini merujuk 54 prompt pada modul 01--04 dan 14
latihan terselesaikan yang sudah diterima tanpa mendeklarasikan ulang ID-nya.
Asesmen kumulatif lanjutan, rubrik bukti, dua laboratorium komputasi, dan
proyek kapstone kini lengkap pada modul 05--13 yang mengikuti peta ini.

% ORIG03-SEGMENT-CODE: OR03-S002
% ORIG03-SEGMENT-ID: d90.orig.v1.tr03.seg0002
% ORIG03-MODULE-ID: diag.module.0001
% SPDX-License-Identifier: CC-BY-SA-4.0
% Karya asli berbahasa Indonesia; semua butir disusun independen.

\section{Diagnostik prasyarat}
\label{orig03:section:diagnostik}

Sepuluh butir berikut menguji prasyarat yang dipakai di seluruh buku. Setiap
subbutir merupakan prompt tersendiri dengan dua tingkat petunjuk, jawaban
ringkas, dan solusi lengkap. Materi ini tidak memuat pemodelan atau algoritma
yang menjadi wilayah O018.

% ORIG03-STABLE-ID: diag.problem.0001
\begin{exercise}[Norma dan ketaksamaan dasar]
\label{orig03:diag:problem:0001}
\begin{enumerate}
% ORIG03-STABLE-ID: diag.prompt.0001
\item\label{orig03:diag:prompt:0001}
Untuk $x=(3,-4)\in\mathbb{R}^2$, hitung $\lVert x\rVert_1$,
$\lVert x\rVert_2$, dan $\lVert x\rVert_\infty$.
% ORIG03-STABLE-ID: diag.prompt.0002
\item\label{orig03:diag:prompt:0002}
Buktikan bahwa untuk setiap $z\in\mathbb{R}^n$,
\[
  \lVert z\rVert_\infty\leq\lVert z\rVert_2
  \leq\lVert z\rVert_1\leq\sqrt n\,\lVert z\rVert_2.
\]
\end{enumerate}
\end{exercise}

% ORIG03-STABLE-ID: diag.hint1.0001
\par\noindent\textbf{Petunjuk tahap 1 (diag.prompt.0001).}
Gunakan langsung tiga definisi norma.
\label{orig03:diag:hint1:0001}
% ORIG03-STABLE-ID: diag.hint2.0001
\par\noindent\textbf{Petunjuk tahap 2 (diag.prompt.0001).}
Nilai mutlak koordinatnya adalah $3$ dan $4$.
\label{orig03:diag:hint2:0001}
% ORIG03-STABLE-ID: diag.answer.0001
\par\noindent\textbf{Jawaban ringkas (diag.prompt.0001).}
$\lVert x\rVert_1=7$, $\lVert x\rVert_2=5$, dan
$\lVert x\rVert_\infty=4$.
\label{orig03:diag:answer:0001}
% ORIG03-STABLE-ID: diag.solution.0001
\par\noindent\textbf{Solusi lengkap (diag.prompt.0001).}
Definisi memberi
$\lVert x\rVert_1=|3|+|-4|=7$,
$\lVert x\rVert_2=\sqrt{3^2+(-4)^2}=5$, dan
$\lVert x\rVert_\infty=\max\{3,4\}=4$.
\label{orig03:diag:solution:0001}

% ORIG03-STABLE-ID: diag.hint1.0002
\par\noindent\textbf{Petunjuk tahap 1 (diag.prompt.0002).}
Bandingkan setiap koordinat dengan jumlah kuadrat, lalu gunakan
Cauchy--Schwarz untuk ketaksamaan terakhir.
\label{orig03:diag:hint1:0002}
% ORIG03-STABLE-ID: diag.hint2.0002
\par\noindent\textbf{Petunjuk tahap 2 (diag.prompt.0002).}
Kembangkan $(\sum_i|z_i|)^2$ dan gunakan
$\sum_i|z_i|\leq\sqrt{\sum_i1^2}\sqrt{\sum_i z_i^2}$.
\label{orig03:diag:hint2:0002}
% ORIG03-STABLE-ID: diag.answer.0002
\par\noindent\textbf{Jawaban ringkas (diag.prompt.0002).}
Rantai tersebut berlaku untuk semua $z$; konstanta $\sqrt n$ tajam pada
vektor yang semua koordinat mutlaknya sama.
\label{orig03:diag:answer:0002}
% ORIG03-STABLE-ID: diag.solution.0002
\par\noindent\textbf{Solusi lengkap (diag.prompt.0002).}
Untuk setiap $i$, $|z_i|^2\leq\sum_j|z_j|^2$, sehingga
$\lVert z\rVert_\infty\leq\lVert z\rVert_2$. Selanjutnya,
$\sum_i|z_i|^2\leq(\sum_i|z_i|)^2$, sehingga
$\lVert z\rVert_2\leq\lVert z\rVert_1$. Akhirnya Cauchy--Schwarz memberi
$\sum_i|z_i|\leq(\sum_i1)^{1/2}(\sum_i|z_i|^2)^{1/2}$.
Ini menghasilkan ketaksamaan terakhir. Jika $|z_1|=\cdots=|z_n|>0$,
ketaksamaan terakhir menjadi kesamaan.
\emph{Rubrik bukti: \ref{orig03:rubric:proof:0001}.}
\label{orig03:diag:solution:0002}

% ORIG03-STABLE-ID: diag.problem.0002
\begin{exercise}[Proyeksi dan ortogonalitas]
\label{orig03:diag:problem:0002}
\begin{enumerate}
% ORIG03-STABLE-ID: diag.prompt.0003
\item\label{orig03:diag:prompt:0003}
Proyeksikan $y=(2,-1)$ pada hiperbidang
$H=\{x\in\mathbb{R}^2:x_1+x_2=0\}$.
% ORIG03-STABLE-ID: diag.prompt.0004
\item\label{orig03:diag:prompt:0004}
Untuk himpunan konveks tertutup takkosong $C$, buktikan bahwa
$p=P_C(y)$ jika dan hanya jika
$\langle y-p,z-p\rangle\leq0$ untuk setiap $z\in C$.
\end{enumerate}
\end{exercise}

% ORIG03-STABLE-ID: diag.hint1.0003
\par\noindent\textbf{Petunjuk tahap 1 (diag.prompt.0003).}
Gunakan normal $a=(1,1)$ dan rumus proyeksi pada $a^Tx=0$.
\label{orig03:diag:hint1:0003}
% ORIG03-STABLE-ID: diag.hint2.0003
\par\noindent\textbf{Petunjuk tahap 2 (diag.prompt.0003).}
Kurangkan $((a^Ty)/(a^Ta))a$ dari $y$.
\label{orig03:diag:hint2:0003}
% ORIG03-STABLE-ID: diag.answer.0003
\par\noindent\textbf{Jawaban ringkas (diag.prompt.0003).}
$P_H(y)=(3/2,-3/2)$.
\label{orig03:diag:answer:0003}
% ORIG03-STABLE-ID: diag.solution.0003
\par\noindent\textbf{Solusi lengkap (diag.prompt.0003).}
Karena $a^Ty=1$ dan $a^Ta=2$,
\[
P_H(y)=y-\frac{a^Ty}{a^Ta}a=(2,-1)-\frac12(1,1)
=\left(\frac32,-\frac32\right).
\]
Jumlah koordinat hasilnya nol, sedangkan selisih $y-P_H(y)$ sejajar normal
$a$, jadi hasil tersebut memang proyeksi ortogonal.
\label{orig03:diag:solution:0003}

% ORIG03-STABLE-ID: diag.hint1.0004
\par\noindent\textbf{Petunjuk tahap 1 (diag.prompt.0004).}
Untuk arah maju gunakan titik $p+t(z-p)$; untuk arah balik kembangkan
$\lVert y-z\rVert^2$.
\label{orig03:diag:hint1:0004}
% ORIG03-STABLE-ID: diag.hint2.0004
\par\noindent\textbf{Petunjuk tahap 2 (diag.prompt.0004).}
Minimalitas fungsi kuadrat dalam $t\in[0,1]$ memberi turunan kanan taknegatif
di $t=0$.
\label{orig03:diag:hint2:0004}
% ORIG03-STABLE-ID: diag.answer.0004
\par\noindent\textbf{Jawaban ringkas (diag.prompt.0004).}
Kondisi hasil kali dalam tersebut tepat merupakan kondisi optimalitas orde
pertama bagi minimisasi jarak kuadrat pada $C$.
\label{orig03:diag:answer:0004}
% ORIG03-STABLE-ID: diag.solution.0004
\par\noindent\textbf{Solusi lengkap (diag.prompt.0004).}
Jika $p=P_C(y)$, maka untuk $z\in C$ dan $t\in[0,1]$, kekonveksan memberi
$p+t(z-p)\in C$. Fungsi
$q(t)=\lVert y-p-t(z-p)\rVert^2$ minimum di $t=0$, sehingga
$q'(0)=-2\langle y-p,z-p\rangle\geq0$. Sebaliknya, bila kondisi hasil kali
dalam berlaku, maka untuk setiap $z\in C$,
\[
\lVert y-z\rVert^2
=\lVert y-p\rVert^2+\lVert z-p\rVert^2
-2\langle y-p,z-p\rangle\geq\lVert y-p\rVert^2.
\]
Jadi $p$ meminimalkan jarak. Keunikan mengikuti kekonveksan ketat jarak
kuadrat.
\emph{Rubrik bukti: \ref{orig03:rubric:proof:0004}.}
\label{orig03:diag:solution:0004}

% ORIG03-STABLE-ID: diag.problem.0003
\begin{exercise}[Kombinasi dan irisan konveks]
\label{orig03:diag:problem:0003}
\begin{enumerate}
% ORIG03-STABLE-ID: diag.prompt.0005
\item\label{orig03:diag:prompt:0005}
Tentukan $t\in[0,1]$ sehingga
$(1,2)=t(0,1)+(1-t)(2,3)$.
% ORIG03-STABLE-ID: diag.prompt.0006
\item\label{orig03:diag:prompt:0006}
Buktikan bahwa irisan sebarang keluarga himpunan konveks adalah konveks.
\end{enumerate}
\end{exercise}

% ORIG03-STABLE-ID: diag.hint1.0005
\par\noindent\textbf{Petunjuk tahap 1 (diag.prompt.0005).}
Samakan koordinat pertama.
\label{orig03:diag:hint1:0005}
% ORIG03-STABLE-ID: diag.hint2.0005
\par\noindent\textbf{Petunjuk tahap 2 (diag.prompt.0005).}
Persamaan $2(1-t)=1$ memberi kandidat tunggal; periksa koordinat kedua.
\label{orig03:diag:hint2:0005}
% ORIG03-STABLE-ID: diag.answer.0005
\par\noindent\textbf{Jawaban ringkas (diag.prompt.0005).}
$t=1/2$.
\label{orig03:diag:answer:0005}
% ORIG03-STABLE-ID: diag.solution.0005
\par\noindent\textbf{Solusi lengkap (diag.prompt.0005).}
Koordinat pertama memberi $2(1-t)=1$, jadi $t=1/2$. Koordinat kedua lalu
bernilai $t+3(1-t)=1/2+3/2=2$. Karena $1/2\in[0,1]$, ini kombinasi konveks
yang sah.
\label{orig03:diag:solution:0005}

% ORIG03-STABLE-ID: diag.hint1.0006
\par\noindent\textbf{Petunjuk tahap 1 (diag.prompt.0006).}
Ambil dua titik dalam irisan dan uji kombinasi konveksnya pada setiap anggota
keluarga.
\label{orig03:diag:hint1:0006}
% ORIG03-STABLE-ID: diag.hint2.0006
\par\noindent\textbf{Petunjuk tahap 2 (diag.prompt.0006).}
Keanggotaan pada irisan berarti keanggotaan serentak pada semua himpunan.
\label{orig03:diag:hint2:0006}
% ORIG03-STABLE-ID: diag.answer.0006
\par\noindent\textbf{Jawaban ringkas (diag.prompt.0006).}
Irisan tetap konveks; irisan kosong juga konveks secara vakum.
\label{orig03:diag:answer:0006}
% ORIG03-STABLE-ID: diag.solution.0006
\par\noindent\textbf{Solusi lengkap (diag.prompt.0006).}
Tuliskan $C=\bigcap_{i\in I}C_i$ dengan setiap $C_i$ konveks. Jika
$x,y\in C$ dan $t\in[0,1]$, maka $x,y\in C_i$ untuk setiap $i$. Karena
$C_i$ konveks, $tx+(1-t)y\in C_i$ untuk setiap $i$, sehingga titik itu berada
di $C$. Bila $C$ kosong, implikasi dalam definisi konveks benar secara vakum.
\emph{Rubrik bukti: \ref{orig03:rubric:proof:0001}.}
\label{orig03:diag:solution:0006}

% ORIG03-STABLE-ID: diag.problem.0004
\begin{exercise}[Kuadratik, gradien, dan spektrum]
\label{orig03:diag:problem:0004}
\begin{enumerate}
% ORIG03-STABLE-ID: diag.prompt.0007
\item\label{orig03:diag:prompt:0007}
Untuk $f(x)=\tfrac12x^TQx+b^Tx$ dengan
$Q=\begin{psmallmatrix}4&1\\1&2\end{psmallmatrix}$ dan $b=(-1,3)$,
hitung $\nabla f(1,-1)$ dan $\nabla^2f$.
% ORIG03-STABLE-ID: diag.prompt.0008
\item\label{orig03:diag:prompt:0008}
Tentukan konstanta kuat-konveks terbesar $\mu$ dan konstanta Lipschitz gradien
terkecil $L$ bagi $f$.
\end{enumerate}
\end{exercise}

% ORIG03-STABLE-ID: diag.hint1.0007
\par\noindent\textbf{Petunjuk tahap 1 (diag.prompt.0007).}
Karena $Q$ simetris, $\nabla f(x)=Qx+b$.
\label{orig03:diag:hint1:0007}
% ORIG03-STABLE-ID: diag.hint2.0007
\par\noindent\textbf{Petunjuk tahap 2 (diag.prompt.0007).}
Hitung $Q(1,-1)^T=(3,-1)^T$.
\label{orig03:diag:hint2:0007}
% ORIG03-STABLE-ID: diag.answer.0007
\par\noindent\textbf{Jawaban ringkas (diag.prompt.0007).}
$\nabla f(1,-1)=(2,2)$ dan $\nabla^2f=Q$.
\label{orig03:diag:answer:0007}
% ORIG03-STABLE-ID: diag.solution.0007
\par\noindent\textbf{Solusi lengkap (diag.prompt.0007).}
Diferensiasi bentuk kuadratik simetris memberi $\nabla f(x)=Qx+b$ dan
$\nabla^2f=Q$. Pada $x=(1,-1)$,
$Qx=(4-1,1-2)=(3,-1)$, sehingga $Qx+b=(2,2)$.
\label{orig03:diag:solution:0007}

% ORIG03-STABLE-ID: diag.hint1.0008
\par\noindent\textbf{Petunjuk tahap 1 (diag.prompt.0008).}
Cari kedua nilai eigen matriks simetris $Q$.
\label{orig03:diag:hint1:0008}
% ORIG03-STABLE-ID: diag.hint2.0008
\par\noindent\textbf{Petunjuk tahap 2 (diag.prompt.0008).}
Persamaan karakteristiknya $\lambda^2-6\lambda+7=0$.
\label{orig03:diag:hint2:0008}
% ORIG03-STABLE-ID: diag.answer.0008
\par\noindent\textbf{Jawaban ringkas (diag.prompt.0008).}
$\mu=3-\sqrt2$ dan $L=3+\sqrt2$.
\label{orig03:diag:answer:0008}
% ORIG03-STABLE-ID: diag.solution.0008
\par\noindent\textbf{Solusi lengkap (diag.prompt.0008).}
Nilai eigen $Q$ adalah akar
$\lambda^2-6\lambda+7=0$, yaitu $3\pm\sqrt2$. Untuk kuadratik dengan Hessian
konstan simetris positif definit, konstanta kuat-konveks terbesar adalah nilai
eigen minimum dan konstanta Lipschitz gradien terkecil adalah nilai eigen
maksimum. Jadi hasilnya $\mu=3-\sqrt2$ dan $L=3+\sqrt2$.
\emph{Rubrik bukti: \ref{orig03:rubric:proof:0002}.}
\label{orig03:diag:solution:0008}

% ORIG03-STABLE-ID: diag.problem.0005
\begin{exercise}[Kondisi orde pertama]
\label{orig03:diag:problem:0005}
\begin{enumerate}
% ORIG03-STABLE-ID: diag.prompt.0009
\item\label{orig03:diag:prompt:0009}
Untuk $f(x)=(x-2)^2+1$, tuliskan penyangga afin orde pertama di $x_0=0$
dan verifikasi bahwa penyangga itu menyentuh $f$ di $x_0$.
% ORIG03-STABLE-ID: diag.prompt.0010
\item\label{orig03:diag:prompt:0010}
Buktikan bahwa jika $f$ konveks dan terdiferensialkan serta
$\nabla f(x^*)=0$, maka $x^*$ adalah minimizer global.
\end{enumerate}
\end{exercise}

% ORIG03-STABLE-ID: diag.hint1.0009
\par\noindent\textbf{Petunjuk tahap 1 (diag.prompt.0009).}
Gunakan $f(x_0)+f'(x_0)(x-x_0)$.
\label{orig03:diag:hint1:0009}
% ORIG03-STABLE-ID: diag.hint2.0009
\par\noindent\textbf{Petunjuk tahap 2 (diag.prompt.0009).}
$f(0)=5$ dan $f'(0)=-4$.
\label{orig03:diag:hint2:0009}
% ORIG03-STABLE-ID: diag.answer.0009
\par\noindent\textbf{Jawaban ringkas (diag.prompt.0009).}
Penyangga afinnya $\ell(x)=5-4x$, dan $\ell(0)=f(0)=5$.
\label{orig03:diag:answer:0009}
% ORIG03-STABLE-ID: diag.solution.0009
\par\noindent\textbf{Solusi lengkap (diag.prompt.0009).}
Turunan $f'(x)=2(x-2)$, jadi $f'(0)=-4$. Karena itu
$\ell(x)=f(0)+f'(0)x=5-4x$. Selisihnya
$f(x)-\ell(x)=x^2\geq0$, dan selisih nol di $x=0$.
\label{orig03:diag:solution:0009}

% ORIG03-STABLE-ID: diag.hint1.0010
\par\noindent\textbf{Petunjuk tahap 1 (diag.prompt.0010).}
Gunakan ketaksamaan orde pertama bagi fungsi konveks.
\label{orig03:diag:hint1:0010}
% ORIG03-STABLE-ID: diag.hint2.0010
\par\noindent\textbf{Petunjuk tahap 2 (diag.prompt.0010).}
Substitusikan gradien nol dalam
$f(y)\geq f(x^*)+\langle\nabla f(x^*),y-x^*\rangle$.
\label{orig03:diag:hint2:0010}
% ORIG03-STABLE-ID: diag.answer.0010
\par\noindent\textbf{Jawaban ringkas (diag.prompt.0010).}
$f(y)\geq f(x^*)$ untuk setiap $y$.
\label{orig03:diag:answer:0010}
% ORIG03-STABLE-ID: diag.solution.0010
\par\noindent\textbf{Solusi lengkap (diag.prompt.0010).}
Kekonveksan dan keterdiferensialan memberi, untuk setiap $y$,
\[
f(y)\geq f(x^*)+\langle\nabla f(x^*),y-x^*\rangle=f(x^*).
\]
Jadi tidak ada titik dengan nilai lebih kecil daripada $x^*$.
\emph{Rubrik bukti: \ref{orig03:rubric:proof:0001}.}
\label{orig03:diag:solution:0010}

% ORIG03-STABLE-ID: diag.problem.0006
\begin{exercise}[Subgradien dan operator proksimal]
\label{orig03:diag:problem:0006}
\begin{enumerate}
% ORIG03-STABLE-ID: diag.prompt.0011
\item\label{orig03:diag:prompt:0011}
Tentukan $\partial|x|$ pada $x=-2$, $x=0$, dan $x=3$.
% ORIG03-STABLE-ID: diag.prompt.0012
\item\label{orig03:diag:prompt:0012}
Hitung $\operatorname{prox}_{|\cdot|}(-3)$.
\end{enumerate}
\end{exercise}

% ORIG03-STABLE-ID: diag.hint1.0011
\par\noindent\textbf{Petunjuk tahap 1 (diag.prompt.0011).}
Di luar nol gunakan turunan; di nol gunakan definisi subgradien.
\label{orig03:diag:hint1:0011}
% ORIG03-STABLE-ID: diag.hint2.0011
\par\noindent\textbf{Petunjuk tahap 2 (diag.prompt.0011).}
Pada nol, $g$ harus memenuhi $|y|\geq gy$ untuk semua $y$.
\label{orig03:diag:hint2:0011}
% ORIG03-STABLE-ID: diag.answer.0011
\par\noindent\textbf{Jawaban ringkas (diag.prompt.0011).}
$\partial|-2|=\{-1\}$, $\partial|0|=[-1,1]$, dan
$\partial|3|=\{1\}$.
\label{orig03:diag:answer:0011}
% ORIG03-STABLE-ID: diag.solution.0011
\par\noindent\textbf{Solusi lengkap (diag.prompt.0011).}
Fungsi $|x|$ terdiferensialkan dengan turunan $-1$ untuk $x<0$ dan $1$
untuk $x>0$. Pada nol, syarat subgradien adalah $|y|\geq gy$ untuk setiap
$y$. Pemilihan $y>0$ memberi $g\leq1$, sedangkan $y<0$ memberi $g\geq-1$.
Jadi subdiferensial di nol adalah $[-1,1]$.
\label{orig03:diag:solution:0011}

% ORIG03-STABLE-ID: diag.hint1.0012
\par\noindent\textbf{Petunjuk tahap 1 (diag.prompt.0012).}
Gunakan rumus ambang lunak dengan ambang satu.
\label{orig03:diag:hint1:0012}
% ORIG03-STABLE-ID: diag.hint2.0012
\par\noindent\textbf{Petunjuk tahap 2 (diag.prompt.0012).}
$S_1(v)=\operatorname{sign}(v)\max\{|v|-1,0\}$.
\label{orig03:diag:hint2:0012}
% ORIG03-STABLE-ID: diag.answer.0012
\par\noindent\textbf{Jawaban ringkas (diag.prompt.0012).}
$\operatorname{prox}_{|\cdot|}(-3)=-2$.
\label{orig03:diag:answer:0012}
% ORIG03-STABLE-ID: diag.solution.0012
\par\noindent\textbf{Solusi lengkap (diag.prompt.0012).}
Operator proksimal meminimalkan $|x|+\tfrac12(x+3)^2$. Karena minimizer
negatif, turunan pada cabang $x<0$ adalah $-1+x+3=x+2$, yang nol pada
$x=-2$. Titik itu konsisten dengan cabangnya dan memberi hasil yang sama
dengan $S_1(-3)=-2$.
\label{orig03:diag:solution:0012}

% ORIG03-STABLE-ID: diag.problem.0007
\begin{exercise}[Ekspektasi dan rata-rata sampel]
\label{orig03:diag:problem:0007}
\begin{enumerate}
% ORIG03-STABLE-ID: diag.prompt.0013
\item\label{orig03:diag:prompt:0013}
Peubah acak $X$ bernilai $-1$ dengan peluang $2/3$ dan $2$ dengan peluang
$1/3$. Hitung $\mathbb{E}X$ dan $\operatorname{Var}(X)$.
% ORIG03-STABLE-ID: diag.prompt.0014
\item\label{orig03:diag:prompt:0014}
Jika $X_1,\dots,X_m$ saling independen dan berdistribusi seperti $X$, tentukan
ekspektasi dan varians $\bar X_m=m^{-1}\sum_{i=1}^mX_i$.
\end{enumerate}
\end{exercise}

% ORIG03-STABLE-ID: diag.hint1.0013
\par\noindent\textbf{Petunjuk tahap 1 (diag.prompt.0013).}
Hitung momen pertama dan kedua.
\label{orig03:diag:hint1:0013}
% ORIG03-STABLE-ID: diag.hint2.0013
\par\noindent\textbf{Petunjuk tahap 2 (diag.prompt.0013).}
$\mathbb{E}X^2=(2/3)1+(1/3)4$.
\label{orig03:diag:hint2:0013}
% ORIG03-STABLE-ID: diag.answer.0013
\par\noindent\textbf{Jawaban ringkas (diag.prompt.0013).}
$\mathbb{E}X=0$ dan $\operatorname{Var}(X)=2$.
\label{orig03:diag:answer:0013}
% ORIG03-STABLE-ID: diag.solution.0013
\par\noindent\textbf{Solusi lengkap (diag.prompt.0013).}
Kita memperoleh
$\mathbb{E}X=(2/3)(-1)+(1/3)2=0$ dan
$\mathbb{E}X^2=(2/3)1+(1/3)4=2$. Maka
$\operatorname{Var}(X)=\mathbb{E}X^2-(\mathbb{E}X)^2=2$.
\label{orig03:diag:solution:0013}

% ORIG03-STABLE-ID: diag.hint1.0014
\par\noindent\textbf{Petunjuk tahap 1 (diag.prompt.0014).}
Gunakan linearitas ekspektasi dan aditivitas varians bagi peubah independen.
\label{orig03:diag:hint1:0014}
% ORIG03-STABLE-ID: diag.hint2.0014
\par\noindent\textbf{Petunjuk tahap 2 (diag.prompt.0014).}
Faktor $1/m$ dikuadratkan ketika menghitung varians.
\label{orig03:diag:hint2:0014}
% ORIG03-STABLE-ID: diag.answer.0014
\par\noindent\textbf{Jawaban ringkas (diag.prompt.0014).}
$\mathbb{E}\bar X_m=0$ dan $\operatorname{Var}(\bar X_m)=2/m$.
\label{orig03:diag:answer:0014}
% ORIG03-STABLE-ID: diag.solution.0014
\par\noindent\textbf{Solusi lengkap (diag.prompt.0014).}
Linearitas memberi
$\mathbb{E}\bar X_m=m^{-1}\sum_i\mathbb{E}X_i=0$. Independensi memberi
\[
\operatorname{Var}(\bar X_m)=\frac1{m^2}\sum_{i=1}^m
\operatorname{Var}(X_i)=\frac1{m^2}(2m)=\frac2m.
\]
\emph{Rubrik bukti: \ref{orig03:rubric:proof:0003}.}
\label{orig03:diag:solution:0014}

% ORIG03-STABLE-ID: diag.problem.0008
\begin{exercise}[Rekurensi kontraktif]
\label{orig03:diag:problem:0008}
\begin{enumerate}
% ORIG03-STABLE-ID: diag.prompt.0015
\item\label{orig03:diag:prompt:0015}
Selesaikan $x_{k+1}=\tfrac34x_k+1$ dengan $x_0=0$ dan tentukan titik
tetapnya.
% ORIG03-STABLE-ID: diag.prompt.0016
\item\label{orig03:diag:prompt:0016}
Tentukan bilangan bulat terkecil $k$ yang menjamin $|x_k-4|\leq10^{-2}$.
\end{enumerate}
\end{exercise}

% ORIG03-STABLE-ID: diag.hint1.0015
\par\noindent\textbf{Petunjuk tahap 1 (diag.prompt.0015).}
Kurangkan titik tetap dari kedua ruas.
\label{orig03:diag:hint1:0015}
% ORIG03-STABLE-ID: diag.hint2.0015
\par\noindent\textbf{Petunjuk tahap 2 (diag.prompt.0015).}
Titik tetap memenuhi $x=\tfrac34x+1$.
\label{orig03:diag:hint2:0015}
% ORIG03-STABLE-ID: diag.answer.0015
\par\noindent\textbf{Jawaban ringkas (diag.prompt.0015).}
Titik tetapnya $4$ dan $x_k=4(1-(3/4)^k)$.
\label{orig03:diag:answer:0015}
% ORIG03-STABLE-ID: diag.solution.0015
\par\noindent\textbf{Solusi lengkap (diag.prompt.0015).}
Persamaan titik tetap memberi $x^*=4$. Dengan $e_k=x_k-4$ diperoleh
$e_{k+1}=\tfrac34e_k$ dan $e_0=-4$. Jadi
$e_k=-4(3/4)^k$ dan $x_k=4(1-(3/4)^k)$.
\label{orig03:diag:solution:0015}

% ORIG03-STABLE-ID: diag.hint1.0016
\par\noindent\textbf{Petunjuk tahap 1 (diag.prompt.0016).}
Gunakan bentuk tertutup galat dari subbutir sebelumnya.
\label{orig03:diag:hint1:0016}
% ORIG03-STABLE-ID: diag.hint2.0016
\par\noindent\textbf{Petunjuk tahap 2 (diag.prompt.0016).}
Selesaikan $(3/4)^k\leq1/400$ dengan logaritma.
\label{orig03:diag:hint2:0016}
% ORIG03-STABLE-ID: diag.answer.0016
\par\noindent\textbf{Jawaban ringkas (diag.prompt.0016).}
$k=21$.
\label{orig03:diag:answer:0016}
% ORIG03-STABLE-ID: diag.solution.0016
\par\noindent\textbf{Solusi lengkap (diag.prompt.0016).}
Karena $|x_k-4|=4(3/4)^k$, syaratnya adalah
$(3/4)^k\leq1/400$. Karena $\log(3/4)<0$,
\[
k\geq\frac{\log(1/400)}{\log(3/4)}\approx20.827.
\]
Bilangan bulat terkecil yang memenuhi adalah $21$.
\label{orig03:diag:solution:0016}

% ORIG03-STABLE-ID: diag.problem.0009
\begin{exercise}[Kerucut normal dan titik tetap proyeksi]
\label{orig03:diag:problem:0009}
\begin{enumerate}
% ORIG03-STABLE-ID: diag.prompt.0017
\item\label{orig03:diag:prompt:0017}
Minimalkan $f(x)=(x-2)^2$ pada $C=[0,1]$ dan verifikasi
$0\in f'(x^*)+N_C(x^*)$ dengan konvensi
$N_C(x)=\{v:v(y-x)\leq0\ \forall y\in C\}$.
% ORIG03-STABLE-ID: diag.prompt.0018
\item\label{orig03:diag:prompt:0018}
Verifikasi bahwa untuk setiap $\gamma>0$, solusi tersebut memenuhi
$x^*=P_C(x^*-\gamma f'(x^*))$.
\end{enumerate}
\end{exercise}

% ORIG03-STABLE-ID: diag.hint1.0017
\par\noindent\textbf{Petunjuk tahap 1 (diag.prompt.0017).}
Minimizer tak terkendala berada di luar interval.
\label{orig03:diag:hint1:0017}
% ORIG03-STABLE-ID: diag.hint2.0017
\par\noindent\textbf{Petunjuk tahap 2 (diag.prompt.0017).}
Pada ujung kanan, $N_{[0,1]}(1)=[0,\infty)$.
\label{orig03:diag:hint2:0017}
% ORIG03-STABLE-ID: diag.answer.0017
\par\noindent\textbf{Jawaban ringkas (diag.prompt.0017).}
$x^*=1$, $f'(1)=-2$, dan $2\in N_C(1)$.
\label{orig03:diag:answer:0017}
% ORIG03-STABLE-ID: diag.solution.0017
\par\noindent\textbf{Solusi lengkap (diag.prompt.0017).}
Fungsi menurun pada $[0,1]$, jadi minimumnya tercapai di $x^*=1$. Turunannya
$f'(x)=2(x-2)$ memberi $f'(1)=-2$. Dari definisi, setiap $v\geq0$ memenuhi
$v(y-1)\leq0$ untuk $y\in[0,1]$, sehingga $N_C(1)=[0,\infty)$. Memilih
$v=2$ menghasilkan $0=-2+2\in f'(1)+N_C(1)$.
\emph{Rubrik bukti: \ref{orig03:rubric:proof:0004}.}
\label{orig03:diag:solution:0017}

% ORIG03-STABLE-ID: diag.hint1.0018
\par\noindent\textbf{Petunjuk tahap 1 (diag.prompt.0018).}
Substitusikan $x^*=1$ dan $f'(1)=-2$.
\label{orig03:diag:hint1:0018}
% ORIG03-STABLE-ID: diag.hint2.0018
\par\noindent\textbf{Petunjuk tahap 2 (diag.prompt.0018).}
Argumen proyeksinya menjadi $1+2\gamma>1$.
\label{orig03:diag:hint2:0018}
% ORIG03-STABLE-ID: diag.answer.0018
\par\noindent\textbf{Jawaban ringkas (diag.prompt.0018).}
$P_{[0,1]}(1+2\gamma)=1=x^*$.
\label{orig03:diag:answer:0018}
% ORIG03-STABLE-ID: diag.solution.0018
\par\noindent\textbf{Solusi lengkap (diag.prompt.0018).}
Untuk $\gamma>0$,
$x^*-\gamma f'(x^*)=1-\gamma(-2)=1+2\gamma$. Proyeksi setiap bilangan lebih
besar daripada satu pada $[0,1]$ adalah satu. Karena itu
$P_C(x^*-\gamma f'(x^*))=1=x^*$.
\label{orig03:diag:solution:0018}

% ORIG03-STABLE-ID: diag.problem.0010
\begin{exercise}[Cauchy--Schwarz dan Young]
\label{orig03:diag:problem:0010}
\begin{enumerate}
% ORIG03-STABLE-ID: diag.prompt.0019
\item\label{orig03:diag:prompt:0019}
Untuk $u=(1,2)$ dan $v=(2,-1)$, hitung hasil kali dalam dan kedua normanya,
lalu periksa ketaksamaan Cauchy--Schwarz.
% ORIG03-STABLE-ID: diag.prompt.0020
\item\label{orig03:diag:prompt:0020}
Buktikan bahwa untuk $a,b\in\mathbb{R}$ dan $\eta>0$,
$2ab\leq\eta a^2+\eta^{-1}b^2$, serta tentukan syarat kesamaan.
\end{enumerate}
\end{exercise}

% ORIG03-STABLE-ID: diag.hint1.0019
\par\noindent\textbf{Petunjuk tahap 1 (diag.prompt.0019).}
Hitung $u^Tv$ sebelum normanya.
\label{orig03:diag:hint1:0019}
% ORIG03-STABLE-ID: diag.hint2.0019
\par\noindent\textbf{Petunjuk tahap 2 (diag.prompt.0019).}
Kedua norma sama dengan $\sqrt5$.
\label{orig03:diag:hint2:0019}
% ORIG03-STABLE-ID: diag.answer.0019
\par\noindent\textbf{Jawaban ringkas (diag.prompt.0019).}
$\langle u,v\rangle=0$ dan $0\leq5$.
\label{orig03:diag:answer:0019}
% ORIG03-STABLE-ID: diag.solution.0019
\par\noindent\textbf{Solusi lengkap (diag.prompt.0019).}
$\langle u,v\rangle=1\cdot2+2\cdot(-1)=0$, sedangkan
$\lVert u\rVert_2=\lVert v\rVert_2=\sqrt5$. Jadi
$|\langle u,v\rangle|=0\leq\lVert u\rVert_2\lVert v\rVert_2=5$.
Vektor-vektor tersebut ortogonal, bukan sejajar, sehingga ketaksamaan tidak
menjadi kesamaan pada batas atas yang positif.
\label{orig03:diag:solution:0019}

% ORIG03-STABLE-ID: diag.hint1.0020
\par\noindent\textbf{Petunjuk tahap 1 (diag.prompt.0020).}
Mulai dari kuadrat yang taknegatif.
\label{orig03:diag:hint1:0020}
% ORIG03-STABLE-ID: diag.hint2.0020
\par\noindent\textbf{Petunjuk tahap 2 (diag.prompt.0020).}
Kembangkan $(\sqrt\eta\,a-b/\sqrt\eta)^2$.
\label{orig03:diag:hint2:0020}
% ORIG03-STABLE-ID: diag.answer.0020
\par\noindent\textbf{Jawaban ringkas (diag.prompt.0020).}
Ketaksamaan selalu berlaku dan kesamaan terjadi tepat ketika $b=\eta a$.
\label{orig03:diag:answer:0020}
% ORIG03-STABLE-ID: diag.solution.0020
\par\noindent\textbf{Solusi lengkap (diag.prompt.0020).}
Karena
\[
0\leq\left(\sqrt\eta\,a-\frac{b}{\sqrt\eta}\right)^2
=\eta a^2-2ab+\eta^{-1}b^2,
\]
pemindahan suku memberi ketaksamaan yang diminta. Kuadrat itu nol tepat jika
$\sqrt\eta\,a=b/\sqrt\eta$, yaitu $b=\eta a$.
\emph{Rubrik bukti: \ref{orig03:rubric:proof:0001}.}
\label{orig03:diag:solution:0020}

% ORIG03-SEGMENT-CODE: OR03-S003
% ORIG03-SEGMENT-ID: d90.orig.v1.tr03.seg0003
% ORIG03-MODULE-ID: ps01.module.0001
% SPDX-License-Identifier: CC-BY-SA-4.0
% Karya asli berbahasa Indonesia; semua butir disusun independen.

\section{Set soal I: dasar konveks}
\label{orig03:section:set-soal-dasar-konveks}

Lima masalah berikut melatih definisi konveksitas, kelengkungan kuadratik,
ketaksamaan Jensen, konjugasi Fenchel, dan proyeksi. Setiap subbutir adalah
prompt tersendiri dengan dua petunjuk, jawaban ringkas, dan solusi lengkap.

% ORIG03-STABLE-ID: ps01.problem.0001
\begin{exercise}[Maksimum afin, subdiferensial, dan optimalitas]
\label{orig03:ps01:problem:0001}
\begin{enumerate}
% ORIG03-STABLE-ID: ps01.prompt.0001
\item\label{orig03:ps01:prompt:0001}
Untuk $a_1,a_2\in\mathbb{R}^n$ dan $b_1,b_2\in\mathbb{R}$, buktikan bahwa
\[
  f(x)=\max\{a_1^\top x+b_1,\,a_2^\top x+b_2\}
\]
adalah fungsi konveks.
% ORIG03-STABLE-ID: ps01.prompt.0002
\item\label{orig03:ps01:prompt:0002}
Tentukan $\partial |x|$ untuk setiap $x\in\mathbb{R}$.
% ORIG03-STABLE-ID: ps01.prompt.0003
\item\label{orig03:ps01:prompt:0003}
Tentukan peminimum dan nilai minimum
\[
  \phi(x)=|x|+\frac12(x-1)^2,
\]
serta verifikasi syarat optimalitas subgradien di titik tersebut.
\end{enumerate}
\end{exercise}

% ORIG03-STABLE-ID: ps01.hint1.0001
\par\noindent\textbf{Petunjuk tahap 1 (ps01.prompt.0001).}
Tuliskan kedua fungsi afin sebagai $\ell_i(x)=a_i^\top x+b_i$.
\label{orig03:ps01:hint1:0001}
% ORIG03-STABLE-ID: ps01.hint2.0001
\par\noindent\textbf{Petunjuk tahap 2 (ps01.prompt.0001).}
Untuk $t\in[0,1]$, batasi setiap
$t\ell_i(x)+(1-t)\ell_i(y)$ dengan kombinasi maksimum di $x$ dan $y$.
\label{orig03:ps01:hint2:0001}
% ORIG03-STABLE-ID: ps01.answer.0001
\par\noindent\textbf{Jawaban ringkas (ps01.prompt.0001).}
$f$ konveks karena maksimum berhingga fungsi-fungsi afin adalah konveks.
\label{orig03:ps01:answer:0001}
% ORIG03-STABLE-ID: ps01.solution.0001
\par\noindent\textbf{Solusi lengkap (ps01.prompt.0001).}
Ambil $x,y\in\mathbb{R}^n$ dan $t\in[0,1]$. Untuk setiap $i\in\{1,2\}$,
\[
 \ell_i(tx+(1-t)y)
 =t\ell_i(x)+(1-t)\ell_i(y)
 \leq t f(x)+(1-t)f(y).
\]
Mengambil maksimum terhadap $i$ pada ruas kiri memberi
$f(tx+(1-t)y)\leq t f(x)+(1-t)f(y)$, yakni definisi konveksitas.
\emph{Rubrik bukti: \ref{orig03:rubric:proof:0001}.}
\label{orig03:ps01:solution:0001}

% ORIG03-STABLE-ID: ps01.hint1.0002
\par\noindent\textbf{Petunjuk tahap 1 (ps01.prompt.0002).}
Di luar nol, $|x|$ terdiferensialkan.
\label{orig03:ps01:hint1:0002}
% ORIG03-STABLE-ID: ps01.hint2.0002
\par\noindent\textbf{Petunjuk tahap 2 (ps01.prompt.0002).}
Di $x=0$, gunakan definisi $g\in\partial |\cdot|(0)$, yaitu
$|y|\geq gy$ untuk semua $y$.
\label{orig03:ps01:hint2:0002}
% ORIG03-STABLE-ID: ps01.answer.0002
\par\noindent\textbf{Jawaban ringkas (ps01.prompt.0002).}
\[
 \partial |x|=
 \begin{cases}
   \{-1\},&x<0,\\
   [-1,1],&x=0,\\
   \{1\},&x>0.
 \end{cases}
\]
\label{orig03:ps01:answer:0002}
% ORIG03-STABLE-ID: ps01.solution.0002
\par\noindent\textbf{Solusi lengkap (ps01.prompt.0002).}
Untuk $x>0$ dan $x<0$, turunan $|x|$ masing-masing adalah $1$ dan $-1$;
subdiferensial fungsi konveks yang terdiferensialkan adalah singleton yang
berisi turunannya. Di nol, syarat subgradien adalah $|y|\geq gy$ untuk semua
$y$. Memilih $y>0$ memberi $g\leq1$, sedangkan $y<0$ memberi $g\geq-1$.
Sebaliknya, setiap $g\in[-1,1]$ memenuhi $gy\leq|y|$, sehingga interval
tersebut tepat merupakan subdiferensial di nol.
\label{orig03:ps01:solution:0002}

% ORIG03-STABLE-ID: ps01.hint1.0003
\par\noindent\textbf{Petunjuk tahap 1 (ps01.prompt.0003).}
Gunakan $0\in\partial\phi(x)$.
\label{orig03:ps01:hint1:0003}
% ORIG03-STABLE-ID: ps01.hint2.0003
\par\noindent\textbf{Petunjuk tahap 2 (ps01.prompt.0003).}
Uji titik $x=0$ dengan
$\partial\phi(0)=[-1,1]+(0-1)$.
\label{orig03:ps01:hint2:0003}
% ORIG03-STABLE-ID: ps01.answer.0003
\par\noindent\textbf{Jawaban ringkas (ps01.prompt.0003).}
Peminimum uniknya adalah $x^\star=0$ dan nilai minimumnya $\phi(x^\star)=1/2$.
\label{orig03:ps01:answer:0003}
% ORIG03-STABLE-ID: ps01.solution.0003
\par\noindent\textbf{Solusi lengkap (ps01.prompt.0003).}
Karena suku kuadratik bersifat konveks kuat, $\phi$ mempunyai paling banyak
satu peminimum. Di nol,
\[
 \partial\phi(0)=\partial|\cdot|(0)+(0-1)=[-1,1]-1=[-2,0],
\]
yang memuat nol. Maka syarat optimalitas perlu dan cukup untuk fungsi konveks
terpenuhi, sehingga $x^\star=0$. Substitusi memberi
$\phi(0)=0+\tfrac12=\tfrac12$.
\label{orig03:ps01:solution:0003}

% ORIG03-STABLE-ID: ps01.problem.0002
\begin{exercise}[Parameter kelengkungan dan langkah gradien]
\label{orig03:ps01:problem:0002}
Untuk
\[
 Q=\begin{pmatrix}1&0\\0&4\end{pmatrix},
 \qquad f(x)=\frac12x^\top Qx,
\]
jawab pertanyaan berikut.
\begin{enumerate}
% ORIG03-STABLE-ID: ps01.prompt.0004
\item\label{orig03:ps01:prompt:0004}
Tentukan konstanta konveksitas kuat terbesar $\mu$ dan konstanta Lipschitz
terkecil $L$ untuk gradien $f$ dalam norma Euclid.
% ORIG03-STABLE-ID: ps01.prompt.0005
\item\label{orig03:ps01:prompt:0005}
Untuk langkah gradien $x^+=x-\tfrac14\nabla f(x)$, hitung $x^+$ dan buktikan
$f(x^+)\leq\tfrac9{16}f(x)$. Tentukan kapan batas itu tajam.
\end{enumerate}
\end{exercise}

% ORIG03-STABLE-ID: ps01.hint1.0004
\par\noindent\textbf{Petunjuk tahap 1 (ps01.prompt.0004).}
Gunakan nilai eigen ekstrem Hessian.
\label{orig03:ps01:hint1:0004}
% ORIG03-STABLE-ID: ps01.hint2.0004
\par\noindent\textbf{Petunjuk tahap 2 (ps01.prompt.0004).}
Hessian $\nabla^2f=Q$ mempunyai nilai eigen $1$ dan $4$.
\label{orig03:ps01:hint2:0004}
% ORIG03-STABLE-ID: ps01.answer.0004
\par\noindent\textbf{Jawaban ringkas (ps01.prompt.0004).}
$\mu=1$ dan $L=4$.
\label{orig03:ps01:answer:0004}
% ORIG03-STABLE-ID: ps01.solution.0004
\par\noindent\textbf{Solusi lengkap (ps01.prompt.0004).}
Untuk kuadratik simetris, konstanta konveksitas kuat terbesar adalah
$\lambda_{\min}(Q)$ dan konstanta Lipschitz gradien terkecil adalah
$\lVert Q\rVert_2=\lambda_{\max}(Q)$. Karena spektrum $Q$ adalah $\{1,4\}$,
diperoleh $\mu=1$ dan $L=4$. Keduanya tajam pada masing-masing arah eigen.
\label{orig03:ps01:solution:0004}

% ORIG03-STABLE-ID: ps01.hint1.0005
\par\noindent\textbf{Petunjuk tahap 1 (ps01.prompt.0005).}
Hitung $(I-\tfrac14Q)x$.
\label{orig03:ps01:hint1:0005}
% ORIG03-STABLE-ID: ps01.hint2.0005
\par\noindent\textbf{Petunjuk tahap 2 (ps01.prompt.0005).}
Bandingkan $\tfrac12(\tfrac34x_1)^2$ dengan
$\tfrac9{16}\,\tfrac12(x_1^2+4x_2^2)$.
\label{orig03:ps01:hint2:0005}
% ORIG03-STABLE-ID: ps01.answer.0005
\par\noindent\textbf{Jawaban ringkas (ps01.prompt.0005).}
$x^+=(\tfrac34x_1,0)$; faktor $9/16$ tajam tepat pada arah sumbu pertama
(selain kasus nol, $x_2=0$).
\label{orig03:ps01:answer:0005}
% ORIG03-STABLE-ID: ps01.solution.0005
\par\noindent\textbf{Solusi lengkap (ps01.prompt.0005).}
Karena $\nabla f(x)=(x_1,4x_2)$,
\[
 x^+=x-\frac14\nabla f(x)=\left(\frac34x_1,0\right).
\]
Jadi
\[
 f(x^+)=\frac12\frac9{16}x_1^2
 \leq \frac9{16}\frac12(x_1^2+4x_2^2)=\frac9{16}f(x).
\]
Selisih ruas kanan dan kiri adalah $\tfrac98x_2^2$, sehingga kesamaan
terjadi jika dan hanya jika $x_2=0$.
\label{orig03:ps01:solution:0005}

% ORIG03-STABLE-ID: ps01.problem.0003
\begin{exercise}[Ketaksamaan Jensen]
\label{orig03:ps01:problem:0003}
Ambil $f(x)=x^2$, titik $x_1=-1$, $x_2=3$, dan bobot
$\theta_1=1/4$, $\theta_2=3/4$.
\begin{enumerate}
% ORIG03-STABLE-ID: ps01.prompt.0006
\item\label{orig03:ps01:prompt:0006}
Hitung kedua ruas ketaksamaan Jensen untuk data tersebut.
% ORIG03-STABLE-ID: ps01.prompt.0007
\item\label{orig03:ps01:prompt:0007}
Buktikan ketaksamaan Jensen dua titik langsung dari definisi fungsi konveks.
\end{enumerate}
\end{exercise}

% ORIG03-STABLE-ID: ps01.hint1.0006
\par\noindent\textbf{Petunjuk tahap 1 (ps01.prompt.0006).}
Hitung dahulu rata-rata berbobot kedua titik.
\label{orig03:ps01:hint1:0006}
% ORIG03-STABLE-ID: ps01.hint2.0006
\par\noindent\textbf{Petunjuk tahap 2 (ps01.prompt.0006).}
$\tfrac14(-1)+\tfrac34(3)=2$.
\label{orig03:ps01:hint2:0006}
% ORIG03-STABLE-ID: ps01.answer.0006
\par\noindent\textbf{Jawaban ringkas (ps01.prompt.0006).}
Ruas kiri adalah $4$ dan ruas kanan adalah $7$, jadi $4\leq7$.
\label{orig03:ps01:answer:0006}
% ORIG03-STABLE-ID: ps01.solution.0006
\par\noindent\textbf{Solusi lengkap (ps01.prompt.0006).}
Rata-rata berbobotnya $\bar x=2$, sehingga $f(\bar x)=2^2=4$.
Sementara itu,
\[
 \frac14f(-1)+\frac34f(3)=\frac14(1)+\frac34(9)=7.
\]
Karena $4\leq7$, data tersebut memenuhi ketaksamaan Jensen.
\label{orig03:ps01:solution:0006}

% ORIG03-STABLE-ID: ps01.hint1.0007
\par\noindent\textbf{Petunjuk tahap 1 (ps01.prompt.0007).}
Tuliskan definisi konveksitas dengan parameter $t\in[0,1]$.
\label{orig03:ps01:hint1:0007}
% ORIG03-STABLE-ID: ps01.hint2.0007
\par\noindent\textbf{Petunjuk tahap 2 (ps01.prompt.0007).}
Ganti $t$ dengan $\theta_1$ dan $1-t$ dengan $\theta_2$.
\label{orig03:ps01:hint2:0007}
% ORIG03-STABLE-ID: ps01.answer.0007
\par\noindent\textbf{Jawaban ringkas (ps01.prompt.0007).}
Untuk $\theta_1,\theta_2\geq0$ dan $\theta_1+\theta_2=1$,
$f(\theta_1x_1+\theta_2x_2)\leq\theta_1f(x_1)+\theta_2f(x_2)$.
\label{orig03:ps01:answer:0007}
% ORIG03-STABLE-ID: ps01.solution.0007
\par\noindent\textbf{Solusi lengkap (ps01.prompt.0007).}
Definisi konveksitas menyatakan bahwa, untuk semua $x_1,x_2$ dan
$t\in[0,1]$,
\[
 f(tx_1+(1-t)x_2)\leq t f(x_1)+(1-t)f(x_2).
\]
Setiap pasangan bobot nonnegatif dengan jumlah satu dapat ditulis sebagai
$t=\theta_1$ dan $1-t=\theta_2$. Substitusi memberi persis ketaksamaan Jensen
dua titik.
\emph{Rubrik bukti: \ref{orig03:rubric:proof:0001}.}
\label{orig03:ps01:solution:0007}

% ORIG03-STABLE-ID: ps01.problem.0004
\begin{exercise}[Konjugat kuadratik dan celah Fenchel--Young]
\label{orig03:ps01:problem:0004}
Untuk $a>0$, definisikan $f(x)=\tfrac a2x^2$ pada $\mathbb{R}$.
\begin{enumerate}
% ORIG03-STABLE-ID: ps01.prompt.0008
\item\label{orig03:ps01:prompt:0008}
Hitung konjugat Fenchel $f^*(y)$.
% ORIG03-STABLE-ID: ps01.prompt.0009
\item\label{orig03:ps01:prompt:0009}
Hitung $f^{**}(x)$ secara langsung dan bandingkan dengan $f(x)$.
% ORIG03-STABLE-ID: ps01.prompt.0010
\item\label{orig03:ps01:prompt:0010}
Nyatakan celah Fenchel--Young $f(x)+f^*(y)-xy$ sebagai kuadrat dan tentukan
syarat kesamaan.
\end{enumerate}
\end{exercise}

% ORIG03-STABLE-ID: ps01.hint1.0008
\par\noindent\textbf{Petunjuk tahap 1 (ps01.prompt.0008).}
Maksimalkan $xy-\tfrac a2x^2$ terhadap $x$.
\label{orig03:ps01:hint1:0008}
% ORIG03-STABLE-ID: ps01.hint2.0008
\par\noindent\textbf{Petunjuk tahap 2 (ps01.prompt.0008).}
Titik stasionernya memenuhi $y-ax=0$.
\label{orig03:ps01:hint2:0008}
% ORIG03-STABLE-ID: ps01.answer.0008
\par\noindent\textbf{Jawaban ringkas (ps01.prompt.0008).}
$f^*(y)=y^2/(2a)$.
\label{orig03:ps01:answer:0008}
% ORIG03-STABLE-ID: ps01.solution.0008
\par\noindent\textbf{Solusi lengkap (ps01.prompt.0008).}
Menurut definisi,
\[
 f^*(y)=\sup_x\left\{xy-\frac a2x^2\right\}.
\]
Fungsi di dalam supremum cekung kuat terhadap $x$; turunannya nol di
$x=y/a$. Substitusi menghasilkan
$y(y/a)-\tfrac a2(y/a)^2=y^2/(2a)$.
\label{orig03:ps01:solution:0008}

% ORIG03-STABLE-ID: ps01.hint1.0009
\par\noindent\textbf{Petunjuk tahap 1 (ps01.prompt.0009).}
Terapkan kembali definisi konjugat pada $f^*(y)=y^2/(2a)$.
\label{orig03:ps01:hint1:0009}
% ORIG03-STABLE-ID: ps01.hint2.0009
\par\noindent\textbf{Petunjuk tahap 2 (ps01.prompt.0009).}
Maksimum $xy-y^2/(2a)$ tercapai di $y=ax$.
\label{orig03:ps01:hint2:0009}
% ORIG03-STABLE-ID: ps01.answer.0009
\par\noindent\textbf{Jawaban ringkas (ps01.prompt.0009).}
$f^{**}(x)=ax^2/2=f(x)$.
\label{orig03:ps01:answer:0009}
% ORIG03-STABLE-ID: ps01.solution.0009
\par\noindent\textbf{Solusi lengkap (ps01.prompt.0009).}
Kita memperoleh
\[
 f^{**}(x)=\sup_y\left\{xy-\frac{y^2}{2a}\right\}.
\]
Turunan terhadap $y$ adalah $x-y/a$, jadi pemaksimum uniknya $y=ax$.
Nilai supremumnya $a x^2-a x^2/2=a x^2/2=f(x)$.
\label{orig03:ps01:solution:0009}

% ORIG03-STABLE-ID: ps01.hint1.0010
\par\noindent\textbf{Petunjuk tahap 1 (ps01.prompt.0010).}
Samakan penyebut pada $\tfrac a2x^2+\tfrac1{2a}y^2-xy$.
\label{orig03:ps01:hint1:0010}
% ORIG03-STABLE-ID: ps01.hint2.0010
\par\noindent\textbf{Petunjuk tahap 2 (ps01.prompt.0010).}
Kenali kuadrat dari $ax-y$.
\label{orig03:ps01:hint2:0010}
% ORIG03-STABLE-ID: ps01.answer.0010
\par\noindent\textbf{Jawaban ringkas (ps01.prompt.0010).}
Celahnya $(ax-y)^2/(2a)\geq0$, dan kesamaan berlaku tepat ketika $y=ax$.
\label{orig03:ps01:answer:0010}
% ORIG03-STABLE-ID: ps01.solution.0010
\par\noindent\textbf{Solusi lengkap (ps01.prompt.0010).}
Perhitungan langsung memberi
\[
 f(x)+f^*(y)-xy
 =\frac a2x^2+\frac{y^2}{2a}-xy
 =\frac{(ax-y)^2}{2a}.
\]
Karena $a>0$, nilai ini nonnegatif. Nilainya nol jika dan hanya jika
$ax-y=0$, yakni $y=ax=\nabla f(x)$.
\emph{Rubrik bukti: \ref{orig03:rubric:proof:0002}.}
\label{orig03:ps01:solution:0010}

% ORIG03-STABLE-ID: ps01.problem.0005
\begin{exercise}[Proyeksi pada kotak]
\label{orig03:ps01:problem:0005}
Misalkan $C=[0,3]^3$ dan $v=(-1,2,4)$.
\begin{enumerate}
% ORIG03-STABLE-ID: ps01.prompt.0011
\item\label{orig03:ps01:prompt:0011}
Hitung proyeksi Euclid $P_C(v)$.
% ORIG03-STABLE-ID: ps01.prompt.0012
\item\label{orig03:ps01:prompt:0012}
Buktikan bahwa proyeksi pada kotak $[\ell,u]^n$ diperoleh dengan memotong
setiap koordinat ke interval $[\ell,u]$.
\end{enumerate}
\end{exercise}

% ORIG03-STABLE-ID: ps01.hint1.0011
\par\noindent\textbf{Petunjuk tahap 1 (ps01.prompt.0011).}
Potong setiap koordinat ke interval $[0,3]$.
\label{orig03:ps01:hint1:0011}
% ORIG03-STABLE-ID: ps01.hint2.0011
\par\noindent\textbf{Petunjuk tahap 2 (ps01.prompt.0011).}
$-1$ dipetakan ke $0$, $2$ tetap, dan $4$ dipetakan ke $3$.
\label{orig03:ps01:hint2:0011}
% ORIG03-STABLE-ID: ps01.answer.0011
\par\noindent\textbf{Jawaban ringkas (ps01.prompt.0011).}
$P_C(v)=(0,2,3)$.
\label{orig03:ps01:answer:0011}
% ORIG03-STABLE-ID: ps01.solution.0011
\par\noindent\textbf{Solusi lengkap (ps01.prompt.0011).}
Meminimumkan $\tfrac12\lVert x-v\rVert_2^2$ pada $C$ terpisah menjadi tiga
masalah skalar. Peminimum $\tfrac12(x_i-v_i)^2$ pada $[0,3]$ masing-masing
adalah $0$, $2$, dan $3$. Karena jumlah mencapai minimum tepat ketika setiap
sukunya minimum, $P_C(v)=(0,2,3)$.
\label{orig03:ps01:solution:0011}

% ORIG03-STABLE-ID: ps01.hint1.0012
\par\noindent\textbf{Petunjuk tahap 1 (ps01.prompt.0012).}
Tuliskan kuadrat jarak sebagai jumlah fungsi satu koordinat.
\label{orig03:ps01:hint1:0012}
% ORIG03-STABLE-ID: ps01.hint2.0012
\par\noindent\textbf{Petunjuk tahap 2 (ps01.prompt.0012).}
Peminimum skalar adalah $\ell$ jika $v_i<\ell$, $v_i$ jika
$\ell\leq v_i\leq u$, dan $u$ jika $v_i>u$.
\label{orig03:ps01:hint2:0012}
% ORIG03-STABLE-ID: ps01.answer.0012
\par\noindent\textbf{Jawaban ringkas (ps01.prompt.0012).}
\[
 [P_{[\ell,u]^n}(v)]_i=\min\{u,\max\{\ell,v_i\}\}.
\]
\label{orig03:ps01:answer:0012}
% ORIG03-STABLE-ID: ps01.solution.0012
\par\noindent\textbf{Solusi lengkap (ps01.prompt.0012).}
Masalah proyeksi adalah
\[
 \min_{\ell\leq x_i\leq u}\frac12\sum_{i=1}^n(x_i-v_i)^2.
\]
Daerah layak merupakan hasil kali interval dan objektif merupakan jumlah
suku yang masing-masing hanya bergantung pada satu koordinat. Karena itu,
setiap koordinat dapat diminimumkan secara independen. Fungsi skalar
$(s-v_i)^2/2$ berkurang hingga $s=v_i$ lalu meningkat; pembatasan pada
$[\ell,u]$ menghasilkan $s=\min\{u,\max\{\ell,v_i\}\}$. Menggabungkan semua
koordinat membuktikan rumus tersebut.
\emph{Rubrik bukti: \ref{orig03:rubric:proof:0004}.}
\label{orig03:ps01:solution:0012}

% ORIG03-SEGMENT-CODE: OR03-S004
% ORIG03-SEGMENT-ID: d90.orig.v1.tr03.seg0004
% ORIG03-MODULE-ID: ps02.module.0001
% SPDX-License-Identifier: CC-BY-SA-4.0
% Karya asli berbahasa Indonesia; semua butir disusun independen.

\section{Set soal II: metode proksimal}
\label{orig03:section:set-soal-metode-proksimal}

Lima masalah berikut menghubungkan operator proksimal, pemisahan
maju--mundur, lemma penurunan, resolven monoton, dan laju titik proksimal.

% ORIG03-STABLE-ID: ps02.problem.0001
\begin{exercise}[Ambang lunak sebagai operator proksimal]
\label{orig03:ps02:problem:0001}
Untuk $\lambda>0$, definisikan
$S_\lambda(v)=\operatorname{sign}(v)\max\{|v|-\lambda,0\}$.
\begin{enumerate}
% ORIG03-STABLE-ID: ps02.prompt.0001
\item\label{orig03:ps02:prompt:0001}
Buktikan bahwa $\operatorname{prox}_{\lambda|\cdot|}(v)=S_\lambda(v)$ untuk
setiap $v\in\mathbb{R}$.
% ORIG03-STABLE-ID: ps02.prompt.0002
\item\label{orig03:ps02:prompt:0002}
Dengan penerapan per koordinat, hitung
$\operatorname{prox}_{(3/4)\lVert\cdot\rVert_1}(3,-1/2,-2)$.
\end{enumerate}
\end{exercise}

% ORIG03-STABLE-ID: ps02.hint1.0001
\par\noindent\textbf{Petunjuk tahap 1 (ps02.prompt.0001).}
Tuliskan syarat optimalitas untuk
$\min_x\{\lambda|x|+\tfrac12(x-v)^2\}$.
\label{orig03:ps02:hint1:0001}
% ORIG03-STABLE-ID: ps02.hint2.0001
\par\noindent\textbf{Petunjuk tahap 2 (ps02.prompt.0001).}
Pisahkan kasus $x>0$, $x<0$, dan $x=0$ dalam
$0\in x-v+\lambda\partial|x|$.
\label{orig03:ps02:hint2:0001}
% ORIG03-STABLE-ID: ps02.answer.0001
\par\noindent\textbf{Jawaban ringkas (ps02.prompt.0001).}
\[
 \operatorname{prox}_{\lambda|\cdot|}(v)=
 \begin{cases}
 v-\lambda,&v>\lambda,\\
 0,&|v|\leq\lambda,\\
 v+\lambda,&v<-\lambda,
 \end{cases}
 =S_\lambda(v).
\]
\label{orig03:ps02:answer:0001}
% ORIG03-STABLE-ID: ps02.solution.0001
\par\noindent\textbf{Solusi lengkap (ps02.prompt.0001).}
Objektif bersifat konveks kuat, sehingga syarat
\[
 0\in x-v+\lambda\partial|x|
\]
menentukan peminimum unik. Jika $x>0$, syaratnya $x=v-\lambda$, konsisten
tepat ketika $v>\lambda$. Jika $x<0$, diperoleh $x=v+\lambda$, konsisten
tepat ketika $v<-\lambda$. Jika $x=0$, syarat menjadi
$v\in\lambda[-1,1]$, ekuivalen dengan $|v|\leq\lambda$. Ketiga kasus itu
persis rumus $S_\lambda(v)$.
\emph{Rubrik bukti: \ref{orig03:rubric:proof:0005}.}
\label{orig03:ps02:solution:0001}

% ORIG03-STABLE-ID: ps02.hint1.0002
\par\noindent\textbf{Petunjuk tahap 1 (ps02.prompt.0002).}
Norma satu dan kuadrat jarak keduanya terpisah menurut koordinat.
\label{orig03:ps02:hint1:0002}
% ORIG03-STABLE-ID: ps02.hint2.0002
\par\noindent\textbf{Petunjuk tahap 2 (ps02.prompt.0002).}
Terapkan $S_{3/4}$ pada ketiga koordinat.
\label{orig03:ps02:hint2:0002}
% ORIG03-STABLE-ID: ps02.answer.0002
\par\noindent\textbf{Jawaban ringkas (ps02.prompt.0002).}
$\operatorname{prox}_{(3/4)\lVert\cdot\rVert_1}(3,-1/2,-2)
=(9/4,0,-5/4)$.
\label{orig03:ps02:answer:0002}
% ORIG03-STABLE-ID: ps02.solution.0002
\par\noindent\textbf{Solusi lengkap (ps02.prompt.0002).}
Pemisahan koordinat memberi
\[
 S_{3/4}(3)=3-\frac34=\frac94,
 \qquad S_{3/4}(-1/2)=0,
 \qquad S_{3/4}(-2)=-2+\frac34=-\frac54.
\]
Menggabungkan ketiga peminimum skalar menghasilkan vektor yang dinyatakan.
\label{orig03:ps02:solution:0002}

% ORIG03-STABLE-ID: ps02.problem.0002
\begin{exercise}[Iterasi maju--mundur skalar]
\label{orig03:ps02:problem:0002}
Pertimbangkan
$F(x)=\tfrac12(x-3)^2+|x|$ dan langkah $\gamma=1/2$.
\begin{enumerate}
% ORIG03-STABLE-ID: ps02.prompt.0003
\item\label{orig03:ps02:prompt:0003}
Turunkan pemetaan iterasi maju--mundur $x^+=T(x)$.
% ORIG03-STABLE-ID: ps02.prompt.0004
\item\label{orig03:ps02:prompt:0004}
Dari $x_0=0$, hitung $x_1$ dan $x_2$, lalu tentukan peminimum eksak $F$.
\end{enumerate}
\end{exercise}

% ORIG03-STABLE-ID: ps02.hint1.0003
\par\noindent\textbf{Petunjuk tahap 1 (ps02.prompt.0003).}
Pisahkan bagian mulus $g(x)=\tfrac12(x-3)^2$ dan $h(x)=|x|$.
\label{orig03:ps02:hint1:0003}
% ORIG03-STABLE-ID: ps02.hint2.0003
\par\noindent\textbf{Petunjuk tahap 2 (ps02.prompt.0003).}
Gunakan $T(x)=\operatorname{prox}_{\gamma h}(x-\gamma g'(x))$.
\label{orig03:ps02:hint2:0003}
% ORIG03-STABLE-ID: ps02.answer.0003
\par\noindent\textbf{Jawaban ringkas (ps02.prompt.0003).}
$T(x)=S_{1/2}((x+3)/2)$.
\label{orig03:ps02:answer:0003}
% ORIG03-STABLE-ID: ps02.solution.0003
\par\noindent\textbf{Solusi lengkap (ps02.prompt.0003).}
Karena $g'(x)=x-3$ dan $\gamma=1/2$,
\[
 x-\gamma g'(x)=x-\frac12(x-3)=\frac{x+3}{2}.
\]
Langkah mundur adalah operator proksimal $\operatorname{prox}_{(1/2)|\cdot|}
=S_{1/2}$, sehingga $T(x)=S_{1/2}((x+3)/2)$.
\label{orig03:ps02:solution:0003}

% ORIG03-STABLE-ID: ps02.hint1.0004
\par\noindent\textbf{Petunjuk tahap 1 (ps02.prompt.0004).}
Substitusikan berturut-turut $x_0$ dan $x_1$ ke $T$.
\label{orig03:ps02:hint1:0004}
% ORIG03-STABLE-ID: ps02.hint2.0004
\par\noindent\textbf{Petunjuk tahap 2 (ps02.prompt.0004).}
Untuk peminimum, selesaikan $0\in x-3+\partial|x|$.
\label{orig03:ps02:hint2:0004}
% ORIG03-STABLE-ID: ps02.answer.0004
\par\noindent\textbf{Jawaban ringkas (ps02.prompt.0004).}
$x_1=1$, $x_2=3/2$, dan peminimum uniknya $x^\star=2$.
\label{orig03:ps02:answer:0004}
% ORIG03-STABLE-ID: ps02.solution.0004
\par\noindent\textbf{Solusi lengkap (ps02.prompt.0004).}
Kita memperoleh
\[
 x_1=S_{1/2}(3/2)=1,
 \qquad
 x_2=S_{1/2}(2)=\frac32.
\]
Syarat optimalitas adalah $0\in x-3+\partial|x|$. Pada cabang $x>0$,
syarat itu menjadi $0=x-2$, sehingga $x=2$, yang konsisten dengan cabang.
Cabang $x<0$ tidak menghasilkan titik yang konsisten, dan di nol interval
$-3+[-1,1]$ tidak memuat nol. Konveksitas kuat suku kuadratik menjamin
keunikan, jadi $x^\star=2$.
\label{orig03:ps02:solution:0004}

% ORIG03-STABLE-ID: ps02.problem.0003
\begin{exercise}[Lemma penurunan]
\label{orig03:ps02:problem:0003}
Misalkan $f:\mathbb{R}^n\to\mathbb{R}$ terdiferensialkan dan gradiennya
$L$-Lipschitz.
\begin{enumerate}
% ORIG03-STABLE-ID: ps02.prompt.0005
\item\label{orig03:ps02:prompt:0005}
Buktikan
\[
 f(y)\leq f(x)+\nabla f(x)^\top(y-x)+\frac L2\lVert y-x\rVert_2^2.
\]
% ORIG03-STABLE-ID: ps02.prompt.0006
\item\label{orig03:ps02:prompt:0006}
Untuk $f(x)=\tfrac L2x^2$ pada $\mathbb{R}$, terapkan langkah gradien
$x^+=x-L^{-1}f'(x)$ dan periksa ketajaman batas penurunan yang dihasilkan.
\end{enumerate}
\end{exercise}

% ORIG03-STABLE-ID: ps02.hint1.0005
\par\noindent\textbf{Petunjuk tahap 1 (ps02.prompt.0005).}
Integralkan gradien sepanjang ruas dari $x$ ke $y$.
\label{orig03:ps02:hint1:0005}
% ORIG03-STABLE-ID: ps02.hint2.0005
\par\noindent\textbf{Petunjuk tahap 2 (ps02.prompt.0005).}
Dengan $d=y-x$, batasi
$[\nabla f(x+td)-\nabla f(x)]^\top d$ oleh $Lt\lVert d\rVert_2^2$.
\label{orig03:ps02:hint2:0005}
% ORIG03-STABLE-ID: ps02.answer.0005
\par\noindent\textbf{Jawaban ringkas (ps02.prompt.0005).}
Batas mengikuti teorema dasar kalkulus pada $t\mapsto f(x+t(y-x))$ dan
integral $\int_0^1Lt\,dt=L/2$.
\label{orig03:ps02:answer:0005}
% ORIG03-STABLE-ID: ps02.solution.0005
\par\noindent\textbf{Solusi lengkap (ps02.prompt.0005).}
Tetapkan $d=y-x$. Teorema dasar kalkulus memberi
\[
 f(y)-f(x)=\int_0^1\nabla f(x+td)^\top d\,dt
 =\nabla f(x)^\top d+
 \int_0^1[\nabla f(x+td)-\nabla f(x)]^\top d\,dt.
\]
Cauchy--Schwarz dan sifat Lipschitz menghasilkan
\[
 [\nabla f(x+td)-\nabla f(x)]^\top d
 \leq \lVert\nabla f(x+td)-\nabla f(x)\rVert_2\lVert d\rVert_2
 \leq Lt\lVert d\rVert_2^2.
\]
Mengintegralkan batas terakhir dari nol sampai satu membuktikan klaim.
\emph{Rubrik bukti: \ref{orig03:rubric:proof:0002}.}
\label{orig03:ps02:solution:0005}

% ORIG03-STABLE-ID: ps02.hint1.0006
\par\noindent\textbf{Petunjuk tahap 1 (ps02.prompt.0006).}
Hitung $f'(x)=Lx$.
\label{orig03:ps02:hint1:0006}
% ORIG03-STABLE-ID: ps02.hint2.0006
\par\noindent\textbf{Petunjuk tahap 2 (ps02.prompt.0006).}
Bandingkan hasil dengan
$f(x^+)\leq f(x)-(2L)^{-1}|f'(x)|^2$.
\label{orig03:ps02:hint2:0006}
% ORIG03-STABLE-ID: ps02.answer.0006
\par\noindent\textbf{Jawaban ringkas (ps02.prompt.0006).}
$x^+=0$ dan
$f(x^+)=f(x)-(2L)^{-1}|f'(x)|^2=0$; batasnya mencapai kesamaan.
\label{orig03:ps02:answer:0006}
% ORIG03-STABLE-ID: ps02.solution.0006
\par\noindent\textbf{Solusi lengkap (ps02.prompt.0006).}
Karena $f'(x)=Lx$, langkah berukuran $1/L$ memberi $x^+=x-x=0$.
Lemma penurunan dengan $y=x-L^{-1}f'(x)$ memberi
\[
 f(x^+)\leq f(x)-\frac1{2L}|f'(x)|^2.
\]
Untuk kuadratik ini, ruas kanan sama dengan
$\tfrac L2x^2-\tfrac1{2L}L^2x^2=0=f(0)=f(x^+)$, sehingga batas tersebut
tajam untuk setiap $x$.
\label{orig03:ps02:solution:0006}

% ORIG03-STABLE-ID: ps02.problem.0004
\begin{exercise}[Pemisahan maju--mundur sebagai kontraksi]
\label{orig03:ps02:problem:0004}
Pada $\mathbb{R}$, ambil $A=\partial|\cdot|$, $B(x)=x-2$, dan $\gamma=1/2$.
\begin{enumerate}
% ORIG03-STABLE-ID: ps02.prompt.0007
\item\label{orig03:ps02:prompt:0007}
Hitung $T=J_{\gamma A}(I-\gamma B)$ secara eksplisit.
% ORIG03-STABLE-ID: ps02.prompt.0008
\item\label{orig03:ps02:prompt:0008}
Tentukan titik tetap $T$ dan hubungkan dengan peminimum
$|x|+\tfrac12(x-2)^2$.
% ORIG03-STABLE-ID: ps02.prompt.0009
\item\label{orig03:ps02:prompt:0009}
Buktikan bahwa $T$ adalah kontraksi dengan faktor paling besar $1/2$ dan
simpulkan keunikan titik tetapnya.
\end{enumerate}
\end{exercise}

% ORIG03-STABLE-ID: ps02.hint1.0007
\par\noindent\textbf{Petunjuk tahap 1 (ps02.prompt.0007).}
$J_{\gamma\partial|\cdot|}=\operatorname{prox}_{\gamma|\cdot|}$.
\label{orig03:ps02:hint1:0007}
% ORIG03-STABLE-ID: ps02.hint2.0007
\par\noindent\textbf{Petunjuk tahap 2 (ps02.prompt.0007).}
Hitung $x-\tfrac12(x-2)=x/2+1$.
\label{orig03:ps02:hint2:0007}
% ORIG03-STABLE-ID: ps02.answer.0007
\par\noindent\textbf{Jawaban ringkas (ps02.prompt.0007).}
$T(x)=S_{1/2}(x/2+1)$.
\label{orig03:ps02:answer:0007}
% ORIG03-STABLE-ID: ps02.solution.0007
\par\noindent\textbf{Solusi lengkap (ps02.prompt.0007).}
Definisi resolven dan identitas resolven subdiferensial memberi
\[
 T(x)=\operatorname{prox}_{(1/2)|\cdot|}
 \left(x-\frac12(x-2)\right)
 =S_{1/2}\left(\frac x2+1\right).
\]
\label{orig03:ps02:solution:0007}

% ORIG03-STABLE-ID: ps02.hint1.0008
\par\noindent\textbf{Petunjuk tahap 1 (ps02.prompt.0008).}
Coba $x=1$ dalam rumus eksplisit $T$.
\label{orig03:ps02:hint1:0008}
% ORIG03-STABLE-ID: ps02.hint2.0008
\par\noindent\textbf{Petunjuk tahap 2 (ps02.prompt.0008).}
Bandingkan persamaan titik tetap dengan
$0\in x-2+\partial|x|$.
\label{orig03:ps02:hint2:0008}
% ORIG03-STABLE-ID: ps02.answer.0008
\par\noindent\textbf{Jawaban ringkas (ps02.prompt.0008).}
Titik tetapnya $x^\star=1$, yang juga merupakan peminimum unik objektif.
\label{orig03:ps02:answer:0008}
% ORIG03-STABLE-ID: ps02.solution.0008
\par\noindent\textbf{Solusi lengkap (ps02.prompt.0008).}
Rumus eksplisit memberi $T(1)=S_{1/2}(3/2)=1$. Selain itu,
$x=T(x)$ ekuivalen dengan
$x=J_{\gamma A}(x-\gamma Bx)$, lalu dengan
$0\in Ax+Bx=\partial|x|+x-2$. Ini persis syarat optimalitas untuk
$|x|+\tfrac12(x-2)^2$. Pada $x=1>0$, syaratnya $0=1+1-2$.
Suku kuadratik konveks kuat, jadi peminimum tersebut unik.
\label{orig03:ps02:solution:0008}

% ORIG03-STABLE-ID: ps02.hint1.0009
\par\noindent\textbf{Petunjuk tahap 1 (ps02.prompt.0009).}
Operator proksimal bersifat non-ekspansif.
\label{orig03:ps02:hint1:0009}
% ORIG03-STABLE-ID: ps02.hint2.0009
\par\noindent\textbf{Petunjuk tahap 2 (ps02.prompt.0009).}
Bandingkan argumen ambang lunak untuk $x$ dan $y$.
\label{orig03:ps02:hint2:0009}
% ORIG03-STABLE-ID: ps02.answer.0009
\par\noindent\textbf{Jawaban ringkas (ps02.prompt.0009).}
$|T(x)-T(y)|\leq\tfrac12|x-y|$; karena itu titik tetapnya tunggal.
\label{orig03:ps02:answer:0009}
% ORIG03-STABLE-ID: ps02.solution.0009
\par\noindent\textbf{Solusi lengkap (ps02.prompt.0009).}
Non-ekspansivitas $S_{1/2}$ menghasilkan
\[
 |T(x)-T(y)|
 \leq\left|\left(\frac x2+1\right)-
              \left(\frac y2+1\right)\right|
 =\frac12|x-y|.
\]
Jadi $T$ adalah kontraksi. Jika $x$ dan $y$ keduanya titik tetap, maka
$|x-y|=|T(x)-T(y)|\leq|x-y|/2$, yang memaksa $x=y$.
\emph{Rubrik bukti: \ref{orig03:rubric:proof:0005}.}
\label{orig03:ps02:solution:0009}

% ORIG03-STABLE-ID: ps02.problem.0005
\begin{exercise}[Laju titik proksimal untuk operator linear]
\label{orig03:ps02:problem:0005}
Ambil $A(x)=mx$ dengan $m>0$ dan parameter $\gamma>0$.
\begin{enumerate}
% ORIG03-STABLE-ID: ps02.prompt.0010
\item\label{orig03:ps02:prompt:0010}
Hitung resolven $J_{\gamma A}=(I+\gamma A)^{-1}$.
% ORIG03-STABLE-ID: ps02.prompt.0011
\item\label{orig03:ps02:prompt:0011}
Untuk iterasi titik proksimal $x_{k+1}=J_{\gamma A}(x_k)$, turunkan rumus
eksak $x_k$ dan lajunya menuju nol.
\end{enumerate}
\end{exercise}

% ORIG03-STABLE-ID: ps02.hint1.0010
\par\noindent\textbf{Petunjuk tahap 1 (ps02.prompt.0010).}
Tetapkan $z=J_{\gamma A}(x)$.
\label{orig03:ps02:hint1:0010}
% ORIG03-STABLE-ID: ps02.hint2.0010
\par\noindent\textbf{Petunjuk tahap 2 (ps02.prompt.0010).}
Selesaikan $x=z+\gamma mz$ terhadap $z$.
\label{orig03:ps02:hint2:0010}
% ORIG03-STABLE-ID: ps02.answer.0010
\par\noindent\textbf{Jawaban ringkas (ps02.prompt.0010).}
$J_{\gamma A}(x)=x/(1+\gamma m)$.
\label{orig03:ps02:answer:0010}
% ORIG03-STABLE-ID: ps02.solution.0010
\par\noindent\textbf{Solusi lengkap (ps02.prompt.0010).}
Definisi resolven menyatakan
$z=J_{\gamma A}(x)$ jika dan hanya jika
$x\in z+\gamma A(z)$. Karena $A(z)=mz$, ini menjadi
$x=(1+\gamma m)z$. Faktor $1+\gamma m$ positif, sehingga
$z=x/(1+\gamma m)$.
\emph{Rubrik bukti: \ref{orig03:rubric:proof:0005}.}
\label{orig03:ps02:solution:0010}

% ORIG03-STABLE-ID: ps02.hint1.0011
\par\noindent\textbf{Petunjuk tahap 1 (ps02.prompt.0011).}
Substitusikan rumus resolven ke relasi iterasi.
\label{orig03:ps02:hint1:0011}
% ORIG03-STABLE-ID: ps02.hint2.0011
\par\noindent\textbf{Petunjuk tahap 2 (ps02.prompt.0011).}
Ulangi perkalian oleh konstanta $(1+\gamma m)^{-1}$.
\label{orig03:ps02:hint2:0011}
% ORIG03-STABLE-ID: ps02.answer.0011
\par\noindent\textbf{Jawaban ringkas (ps02.prompt.0011).}
$x_k=(1+\gamma m)^{-k}x_0$ dan
$|x_k|=(1+\gamma m)^{-k}|x_0|\to0$ secara linear.
\label{orig03:ps02:answer:0011}
% ORIG03-STABLE-ID: ps02.solution.0011
\par\noindent\textbf{Solusi lengkap (ps02.prompt.0011).}
Rumus resolven memberi
$x_{k+1}=(1+\gamma m)^{-1}x_k$. Induksi terhadap $k$ menghasilkan
$x_k=(1+\gamma m)^{-k}x_0$. Karena $m,\gamma>0$, faktor
$q=(1+\gamma m)^{-1}$ berada dalam $(0,1)$. Maka
$|x_k|=q^k|x_0|$ dan konvergensinya menuju nol bersifat linear dengan faktor
eksak $q$.
\label{orig03:ps02:solution:0011}

% ORIG03-SEGMENT-CODE: OR03-S005
% ORIG03-SEGMENT-ID: d90.orig.v1.tr03.seg0005
% ORIG03-MODULE-ID: ps03.module.0001
% SPDX-License-Identifier: CC-BY-SA-4.0
% Karya asli berbahasa Indonesia; semua butir disusun independen.

\section{Set soal III: dualitas dan kondisi KKT}
\label{orig03:section:set-soal-dualitas-kkt}

Lima masalah berikut melatih pembentukan fungsi dual, kondisi
Karush--Kuhn--Tucker (KKT), kondisi Slater, dan sensitivitas nilai optimal.

% ORIG03-STABLE-ID: ps03.problem.0001
\begin{exercise}[Kendala persamaan skalar]
\label{orig03:ps03:problem:0001}
Pertimbangkan masalah
\[
 \min_{x\in\mathbb{R}}\ \frac12x^2
 \qquad\text{dengan kendala}\qquad x=1.
\]
\begin{enumerate}
% ORIG03-STABLE-ID: ps03.prompt.0001
\item\label{orig03:ps03:prompt:0001}
Turunkan fungsi dual, lalu tentukan pengali dan nilai dual optimal.
% ORIG03-STABLE-ID: ps03.prompt.0002
\item\label{orig03:ps03:prompt:0002}
Verifikasi semua kondisi KKT dan kesenjangan dualitas nol.
\end{enumerate}
\end{exercise}

% ORIG03-STABLE-ID: ps03.hint1.0001
\par\noindent\textbf{Petunjuk tahap 1 (ps03.prompt.0001).}
Gunakan Lagrangian $L(x,\lambda)=\tfrac12x^2+\lambda(x-1)$.
\label{orig03:ps03:hint1:0001}
% ORIG03-STABLE-ID: ps03.hint2.0001
\par\noindent\textbf{Petunjuk tahap 2 (ps03.prompt.0001).}
Infimum terhadap $x$ tercapai di $x=-\lambda$.
\label{orig03:ps03:hint2:0001}
% ORIG03-STABLE-ID: ps03.answer.0001
\par\noindent\textbf{Jawaban ringkas (ps03.prompt.0001).}
$g(\lambda)=-\tfrac12\lambda^2-\lambda$,
$\lambda^\star=-1$, dan $d^\star=1/2$.
\label{orig03:ps03:answer:0001}
% ORIG03-STABLE-ID: ps03.solution.0001
\par\noindent\textbf{Solusi lengkap (ps03.prompt.0001).}
Meminimumkan Lagrangian terhadap $x$ memberi syarat $x+\lambda=0$, jadi
$x=-\lambda$. Karena Lagrangian konveks kuat terhadap $x$, titik ini memberi
infimum dan
\[
 g(\lambda)=L(-\lambda,\lambda)
 =-\frac12\lambda^2-\lambda.
\]
Masalah dual memaksimalkan kuadratik cekung ini pada
$\lambda\in\mathbb{R}$. Persamaan $g'(\lambda)=-\lambda-1=0$ memberi
$\lambda^\star=-1$ dan $g(-1)=1/2$.
\emph{Rubrik bukti: \ref{orig03:rubric:proof:0006}.}
\label{orig03:ps03:solution:0001}

% ORIG03-STABLE-ID: ps03.hint1.0002
\par\noindent\textbf{Petunjuk tahap 1 (ps03.prompt.0002).}
KKT untuk kendala persamaan terdiri atas kelayakan primal dan stasioneritas.
\label{orig03:ps03:hint1:0002}
% ORIG03-STABLE-ID: ps03.hint2.0002
\par\noindent\textbf{Petunjuk tahap 2 (ps03.prompt.0002).}
Substitusikan $x^\star=1$ dan $\lambda^\star=-1$.
\label{orig03:ps03:hint2:0002}
% ORIG03-STABLE-ID: ps03.answer.0002
\par\noindent\textbf{Jawaban ringkas (ps03.prompt.0002).}
$x^\star=1$ layak dan
$x^\star+\lambda^\star=0$; nilai primal dan dual sama-sama $1/2$.
\label{orig03:ps03:answer:0002}
% ORIG03-STABLE-ID: ps03.solution.0002
\par\noindent\textbf{Solusi lengkap (ps03.prompt.0002).}
Kelayakan primal memberi $x^\star-1=0$. Stasioneritas memberi
$\nabla_xL(x^\star,\lambda^\star)=x^\star+\lambda^\star=1-1=0$.
Tidak ada syarat tanda atau kelengkapan komplementer untuk pengali persamaan.
Nilai primal adalah $p^\star=\tfrac12(1)^2=1/2$, sedangkan hasil dual pada
subbutir sebelumnya memberi $d^\star=1/2$. Jadi
$p^\star-d^\star=0$.
\emph{Rubrik bukti: \ref{orig03:rubric:proof:0006}.}
\label{orig03:ps03:solution:0002}

% ORIG03-STABLE-ID: ps03.problem.0002
\begin{exercise}[Kendala ketaksamaan aktif]
\label{orig03:ps03:problem:0002}
Pertimbangkan
\[
 \min_{x\in\mathbb{R}}\ (x-2)^2
 \qquad\text{dengan kendala}\qquad x\leq1.
\]
\begin{enumerate}
% ORIG03-STABLE-ID: ps03.prompt.0003
\item\label{orig03:ps03:prompt:0003}
Gunakan kondisi KKT untuk menentukan solusi primal dan pengali optimal.
% ORIG03-STABLE-ID: ps03.prompt.0004
\item\label{orig03:ps03:prompt:0004}
Verifikasi kondisi Slater, turunkan fungsi dual, dan periksa nilainya di
pengali optimal.
\end{enumerate}
\end{exercise}

% ORIG03-STABLE-ID: ps03.hint1.0003
\par\noindent\textbf{Petunjuk tahap 1 (ps03.prompt.0003).}
Tuliskan kendala sebagai $x-1\leq0$ dengan pengali $\lambda\geq0$.
\label{orig03:ps03:hint1:0003}
% ORIG03-STABLE-ID: ps03.hint2.0003
\par\noindent\textbf{Petunjuk tahap 2 (ps03.prompt.0003).}
Gunakan stasioneritas $2(x-2)+\lambda=0$ dan
kelengkapan komplementer $\lambda(x-1)=0$.
\label{orig03:ps03:hint2:0003}
% ORIG03-STABLE-ID: ps03.answer.0003
\par\noindent\textbf{Jawaban ringkas (ps03.prompt.0003).}
$x^\star=1$ dan $\lambda^\star=2$.
\label{orig03:ps03:answer:0003}
% ORIG03-STABLE-ID: ps03.solution.0003
\par\noindent\textbf{Solusi lengkap (ps03.prompt.0003).}
Jika kendala tidak aktif, kelengkapan komplementer memberi $\lambda=0$ dan
stasioneritas memberi $x=2$, yang tidak layak. Maka kendala aktif:
$x^\star=1$. Stasioneritas kemudian memberi
$2(1-2)+\lambda^\star=0$, sehingga $\lambda^\star=2\geq0$.
Kelayakan primal, kelayakan dual, stasioneritas, dan kelengkapan komplementer
semuanya terpenuhi. Karena masalah konveks, titik KKT ini optimal.
\emph{Rubrik bukti: \ref{orig03:rubric:proof:0006}.}
\label{orig03:ps03:solution:0003}

% ORIG03-STABLE-ID: ps03.hint1.0004
\par\noindent\textbf{Petunjuk tahap 1 (ps03.prompt.0004).}
$x=0$ adalah kandidat titik Slater.
\label{orig03:ps03:hint1:0004}
% ORIG03-STABLE-ID: ps03.hint2.0004
\par\noindent\textbf{Petunjuk tahap 2 (ps03.prompt.0004).}
Dalam $L(x,\lambda)=(x-2)^2+\lambda(x-1)$, infimum terhadap $x$ tercapai
di $x=2-\lambda/2$.
\label{orig03:ps03:hint2:0004}
% ORIG03-STABLE-ID: ps03.answer.0004
\par\noindent\textbf{Jawaban ringkas (ps03.prompt.0004).}
Slater berlaku; $g(\lambda)=\lambda-\lambda^2/4$ untuk $\lambda\geq0$, dan
$g(2)=1=p^\star$.
\label{orig03:ps03:answer:0004}
% ORIG03-STABLE-ID: ps03.solution.0004
\par\noindent\textbf{Solusi lengkap (ps03.prompt.0004).}
Titik $x=0$ memenuhi $x-1=-1<0$, jadi kondisi Slater berlaku. Untuk
$\lambda\geq0$, peminimum Lagrangian terhadap $x$ adalah
$x(\lambda)=2-\lambda/2$. Substitusi memberi
\[
 g(\lambda)=\left(-\frac\lambda2\right)^2
 +\lambda\left(1-\frac\lambda2\right)
 =\lambda-\frac{\lambda^2}{4}.
\]
Kuadratik cekung ini maksimum di $\lambda=2$ dengan nilai $1$. Nilai primal
di $x^\star=1$ juga $(1-2)^2=1$, sesuai dualitas kuat yang dijamin Slater.
\emph{Rubrik bukti: \ref{orig03:rubric:proof:0006}.}
\label{orig03:ps03:solution:0004}

% ORIG03-STABLE-ID: ps03.problem.0003
\begin{exercise}[Proyeksi norma minimum pada hiperbidang]
\label{orig03:ps03:problem:0003}
Ambil $a\in\mathbb{R}^n\setminus\{0\}$ dan $b\in\mathbb{R}$, lalu
pertimbangkan
\[
 \min_x\ \frac12\lVert x\rVert_2^2
 \qquad\text{dengan kendala}\qquad a^\top x=b.
\]
\begin{enumerate}
% ORIG03-STABLE-ID: ps03.prompt.0005
\item\label{orig03:ps03:prompt:0005}
Selesaikan sistem KKT untuk memperoleh $x^\star$ dan $\lambda^\star$.
% ORIG03-STABLE-ID: ps03.prompt.0006
\item\label{orig03:ps03:prompt:0006}
Turunkan fungsi dual dan verifikasi bahwa nilai primal dan dual optimal sama.
% ORIG03-STABLE-ID: ps03.prompt.0007
\item\label{orig03:ps03:prompt:0007}
Buktikan bahwa solusi primal itu unik dan bahwa kondisi KKT cukup untuk
optimalitas pada masalah ini.
\end{enumerate}
\end{exercise}

% ORIG03-STABLE-ID: ps03.hint1.0005
\par\noindent\textbf{Petunjuk tahap 1 (ps03.prompt.0005).}
Gunakan $L(x,\lambda)=\tfrac12\lVert x\rVert_2^2+
\lambda(a^\top x-b)$.
\label{orig03:ps03:hint1:0005}
% ORIG03-STABLE-ID: ps03.hint2.0005
\par\noindent\textbf{Petunjuk tahap 2 (ps03.prompt.0005).}
Stasioneritas memberi $x=-\lambda a$; substitusikan ke kendala.
\label{orig03:ps03:hint2:0005}
% ORIG03-STABLE-ID: ps03.answer.0005
\par\noindent\textbf{Jawaban ringkas (ps03.prompt.0005).}
\[
 x^\star=\frac{b}{\lVert a\rVert_2^2}a,
 \qquad
 \lambda^\star=-\frac{b}{\lVert a\rVert_2^2}.
\]
\label{orig03:ps03:answer:0005}
% ORIG03-STABLE-ID: ps03.solution.0005
\par\noindent\textbf{Solusi lengkap (ps03.prompt.0005).}
Stasioneritas adalah $x+\lambda a=0$, sehingga $x=-\lambda a$.
Kelayakan persamaan lalu memberi
$-\lambda\lVert a\rVert_2^2=b$. Karena $a\neq0$, penyebutnya positif dan
\[
 \lambda^\star=-\frac{b}{\lVert a\rVert_2^2},
 \qquad
 x^\star=-\lambda^\star a=
 \frac{b}{\lVert a\rVert_2^2}a.
\]
\emph{Rubrik bukti: \ref{orig03:rubric:proof:0006}.}
\label{orig03:ps03:solution:0005}

% ORIG03-STABLE-ID: ps03.hint1.0006
\par\noindent\textbf{Petunjuk tahap 1 (ps03.prompt.0006).}
Ambil infimum Lagrangian di $x=-\lambda a$.
\label{orig03:ps03:hint1:0006}
% ORIG03-STABLE-ID: ps03.hint2.0006
\par\noindent\textbf{Petunjuk tahap 2 (ps03.prompt.0006).}
Maksimalkan kuadratik cekung yang diperoleh terhadap $\lambda\in\mathbb{R}$.
\label{orig03:ps03:hint2:0006}
% ORIG03-STABLE-ID: ps03.answer.0006
\par\noindent\textbf{Jawaban ringkas (ps03.prompt.0006).}
$g(\lambda)=-\tfrac12\lambda^2\lVert a\rVert_2^2-\lambda b$ dan
$p^\star=d^\star=b^2/(2\lVert a\rVert_2^2)$.
\label{orig03:ps03:answer:0006}
% ORIG03-STABLE-ID: ps03.solution.0006
\par\noindent\textbf{Solusi lengkap (ps03.prompt.0006).}
Substitusi $x=-\lambda a$ ke Lagrangian memberi
\[
 g(\lambda)=-\frac12\lambda^2\lVert a\rVert_2^2-\lambda b.
\]
Turunannya nol pada
$\lambda^\star=-b/\lVert a\rVert_2^2$, dan substitusi memberi
$d^\star=b^2/(2\lVert a\rVert_2^2)$. Di sisi primal,
\[
 \frac12\lVert x^\star\rVert_2^2
 =\frac12\frac{b^2\lVert a\rVert_2^2}{\lVert a\rVert_2^4}
 =\frac{b^2}{2\lVert a\rVert_2^2}=d^\star.
\]
\label{orig03:ps03:solution:0006}

% ORIG03-STABLE-ID: ps03.hint1.0007
\par\noindent\textbf{Petunjuk tahap 1 (ps03.prompt.0007).}
Gunakan konveksitas kuat objektif untuk keunikan.
\label{orig03:ps03:hint1:0007}
% ORIG03-STABLE-ID: ps03.hint2.0007
\par\noindent\textbf{Petunjuk tahap 2 (ps03.prompt.0007).}
Gunakan ketaksamaan konveksitas orde pertama pada titik KKT dan hilangkan
suku linear dengan kelayakan persamaan.
\label{orig03:ps03:hint2:0007}
% ORIG03-STABLE-ID: ps03.answer.0007
\par\noindent\textbf{Jawaban ringkas (ps03.prompt.0007).}
Objektif konveks kuat pada himpunan layak konveks, sehingga peminimum tunggal;
stasioneritas dan kelayakan KKT menjamin optimalitas global.
\label{orig03:ps03:answer:0007}
% ORIG03-STABLE-ID: ps03.solution.0007
\par\noindent\textbf{Solusi lengkap (ps03.prompt.0007).}
Fungsi $f(x)=\tfrac12\lVert x\rVert_2^2$ bersifat konveks kuat dengan
parameter satu. Pembatasannya pada hiperbidang afin tetap konveks kuat, jadi
tidak mungkin mempunyai dua peminimum berbeda. Untuk kecukupan, ambil
pasangan KKT $(x^\star,\lambda^\star)$. Stasioneritas memberi
$\nabla f(x^\star)=-\lambda^\star a$. Untuk setiap $x$ layak, konveksitas
memberi
\[
 f(x)\geq f(x^\star)+\nabla f(x^\star)^\top(x-x^\star)
 =f(x^\star)-\lambda^\star a^\top(x-x^\star)=f(x^\star),
\]
karena $a^\top x=a^\top x^\star=b$. Maka titik KKT optimal global, dan
konveksitas kuat membuatnya unik.
\emph{Rubrik bukti: \ref{orig03:rubric:proof:0006}.}
\label{orig03:ps03:solution:0007}

% ORIG03-STABLE-ID: ps03.problem.0004
\begin{exercise}[Sensitivitas ruas kanan]
\label{orig03:ps03:problem:0004}
Definisikan fungsi nilai
\[
 p(u)=\min_x\left\{\frac12x^2:x=u\right\}.
\]
\begin{enumerate}
% ORIG03-STABLE-ID: ps03.prompt.0008
\item\label{orig03:ps03:prompt:0008}
Hitung $p(u)$, pengali optimal $\lambda^\star(u)$ untuk kendala $x-u=0$,
dan hubungan antara $p'(u)$ dan pengali itu.
% ORIG03-STABLE-ID: ps03.prompt.0009
\item\label{orig03:ps03:prompt:0009}
Untuk perturbasi $\delta$, turunkan ekspansi eksak $p(u+\delta)$ dan
identifikasi suku sensitivitas orde pertama serta galat orde keduanya.
\end{enumerate}
\end{exercise}

% ORIG03-STABLE-ID: ps03.hint1.0008
\par\noindent\textbf{Petunjuk tahap 1 (ps03.prompt.0008).}
Kendala langsung memaksa $x=u$.
\label{orig03:ps03:hint1:0008}
% ORIG03-STABLE-ID: ps03.hint2.0008
\par\noindent\textbf{Petunjuk tahap 2 (ps03.prompt.0008).}
Stasioneritas Lagrangian
$\tfrac12x^2+\lambda(x-u)$ adalah $x+\lambda=0$.
\label{orig03:ps03:hint2:0008}
% ORIG03-STABLE-ID: ps03.answer.0008
\par\noindent\textbf{Jawaban ringkas (ps03.prompt.0008).}
$p(u)=u^2/2$, $\lambda^\star(u)=-u$, dan
$p'(u)=u=-\lambda^\star(u)$.
\label{orig03:ps03:answer:0008}
% ORIG03-STABLE-ID: ps03.solution.0008
\par\noindent\textbf{Solusi lengkap (ps03.prompt.0008).}
Satu-satunya titik layak adalah $x^\star(u)=u$, sehingga
$p(u)=u^2/2$. Stasioneritas memberi
$x^\star(u)+\lambda^\star(u)=0$, jadi $\lambda^\star(u)=-u$.
Mendiferensialkan fungsi nilai menghasilkan
$p'(u)=u=-\lambda^\star(u)$. Tanda minus muncul karena kendala ditulis
sebagai $x-u=0$.
\emph{Rubrik bukti: \ref{orig03:rubric:proof:0007}.}
\label{orig03:ps03:solution:0008}

% ORIG03-STABLE-ID: ps03.hint1.0009
\par\noindent\textbf{Petunjuk tahap 1 (ps03.prompt.0009).}
Kembangkan $(u+\delta)^2/2$.
\label{orig03:ps03:hint1:0009}
% ORIG03-STABLE-ID: ps03.hint2.0009
\par\noindent\textbf{Petunjuk tahap 2 (ps03.prompt.0009).}
Ganti $u$ pada suku linear dengan $-\lambda^\star(u)$.
\label{orig03:ps03:hint2:0009}
% ORIG03-STABLE-ID: ps03.answer.0009
\par\noindent\textbf{Jawaban ringkas (ps03.prompt.0009).}
$p(u+\delta)=p(u)-\lambda^\star(u)\delta+\delta^2/2$.
\label{orig03:ps03:answer:0009}
% ORIG03-STABLE-ID: ps03.solution.0009
\par\noindent\textbf{Solusi lengkap (ps03.prompt.0009).}
Perhitungan eksak memberi
\[
 p(u+\delta)=\frac12(u+\delta)^2
 =\frac12u^2+u\delta+\frac12\delta^2
 =p(u)-\lambda^\star(u)\delta+\frac12\delta^2.
\]
Jadi $-\lambda^\star(u)\delta$ adalah perubahan linear yang diprediksi
pengali, sedangkan $\delta^2/2$ adalah galat orde kedua yang eksak.
\emph{Rubrik bukti: \ref{orig03:rubric:proof:0007}.}
\label{orig03:ps03:solution:0009}

% ORIG03-STABLE-ID: ps03.problem.0005
\begin{exercise}[Proyeksi pada bola Euclid melalui KKT]
\label{orig03:ps03:problem:0005}
Ambil $c\in\mathbb{R}^n$ dan $r$ dengan $0<r<\lVert c\rVert_2$.
Pertimbangkan
\[
 \min_x\ \frac12\lVert x-c\rVert_2^2
 \qquad\text{dengan kendala}\qquad \lVert x\rVert_2\leq r.
\]
Gunakan representasi mulus
$g(x)=\tfrac12(\lVert x\rVert_2^2-r^2)\leq0$.
\begin{enumerate}
% ORIG03-STABLE-ID: ps03.prompt.0010
\item\label{orig03:ps03:prompt:0010}
Selesaikan kondisi KKT untuk menentukan $x^\star$ dan pengali
$\mu^\star$.
% ORIG03-STABLE-ID: ps03.prompt.0011
\item\label{orig03:ps03:prompt:0011}
Verifikasi kondisi Slater dan buktikan bahwa pasangan KKT tersebut merupakan
solusi primal-dual global yang unik pada sisi primal.
\end{enumerate}
\end{exercise}

% ORIG03-STABLE-ID: ps03.hint1.0010
\par\noindent\textbf{Petunjuk tahap 1 (ps03.prompt.0010).}
Stasioneritas adalah $x-c+\mu x=0$.
\label{orig03:ps03:hint1:0010}
% ORIG03-STABLE-ID: ps03.hint2.0010
\par\noindent\textbf{Petunjuk tahap 2 (ps03.prompt.0010).}
Karena $c$ berada di luar bola, kendala optimal aktif; gunakan
$x=c/(1+\mu)$ dan $\lVert x\rVert_2=r$.
\label{orig03:ps03:hint2:0010}
% ORIG03-STABLE-ID: ps03.answer.0010
\par\noindent\textbf{Jawaban ringkas (ps03.prompt.0010).}
\[
 x^\star=r\frac{c}{\lVert c\rVert_2},
 \qquad
 \mu^\star=\frac{\lVert c\rVert_2}{r}-1>0.
\]
\label{orig03:ps03:answer:0010}
% ORIG03-STABLE-ID: ps03.solution.0010
\par\noindent\textbf{Solusi lengkap (ps03.prompt.0010).}
Lagrangiannya
$L(x,\mu)=\tfrac12\lVert x-c\rVert_2^2+
\tfrac\mu2(\lVert x\rVert_2^2-r^2)$ dengan $\mu\geq0$.
Stasioneritas memberi $(1+\mu)x=c$. Jika $\mu=0$, maka $x=c$, yang tidak
layak karena $\lVert c\rVert_2>r$. Jadi $\mu^\star>0$, dan kelengkapan
komplementer memaksa $\lVert x^\star\rVert_2=r$. Dari
$x^\star=c/(1+\mu^\star)$ diperoleh
$1+\mu^\star=\lVert c\rVert_2/r$, sehingga rumus yang dinyatakan berlaku.
\emph{Rubrik bukti: \ref{orig03:rubric:proof:0006}.}
\label{orig03:ps03:solution:0010}

% ORIG03-STABLE-ID: ps03.hint1.0011
\par\noindent\textbf{Petunjuk tahap 1 (ps03.prompt.0011).}
$x=0$ memenuhi kendala secara ketat.
\label{orig03:ps03:hint1:0011}
% ORIG03-STABLE-ID: ps03.hint2.0011
\par\noindent\textbf{Petunjuk tahap 2 (ps03.prompt.0011).}
Gabungkan Slater, konveksitas masalah, dan konveksitas kuat objektif.
\label{orig03:ps03:hint2:0011}
% ORIG03-STABLE-ID: ps03.answer.0011
\par\noindent\textbf{Jawaban ringkas (ps03.prompt.0011).}
Slater berlaku di $x=0$; KKT perlu dan cukup, dualitas kuat berlaku, dan
$x^\star$ unik karena objektif konveks kuat.
\label{orig03:ps03:answer:0011}
% ORIG03-STABLE-ID: ps03.solution.0011
\par\noindent\textbf{Solusi lengkap (ps03.prompt.0011).}
Karena $r>0$, titik $x=0$ memenuhi
$g(0)=-r^2/2<0$, jadi kondisi Slater berlaku. Objektif dan fungsi kendala
konveks serta terdiferensialkan; karena itu, di bawah Slater, kondisi KKT
bersifat perlu dan cukup dan tidak ada kesenjangan dualitas. Pasangan pada
subbutir sebelumnya memenuhi kelayakan primal, $\mu^\star\geq0$,
stasioneritas, dan kelengkapan komplementer, sehingga optimal global.
Objektif $\tfrac12\lVert x-c\rVert_2^2$ bersifat konveks kuat, maka
peminimum primal pada himpunan bola yang konveks adalah unik.
\emph{Rubrik bukti: \ref{orig03:rubric:proof:0006}.}
\label{orig03:ps03:solution:0011}

% segment-local-id: OR03-S006
% segment-id: d90.orig.v1.tr03.seg0006
% license: CC BY-SA 4.0
% provenance: Materi asesmen asli berbahasa Indonesia, disusun secara independen.

\section{Set soal IV: metode stokastik}
\label{orig03:ps04}

% stable-id: d90.orig.v1.tr03.ps04.exercise.0001
\begin{exercise}[Pencuplikan tak seragam dan koreksi kepentingan]
\label{orig03:ps04:exercise:0001}
Misalkan
\[
  F(x)=\frac1N\sum_{i=1}^N f_i(x),
  \qquad g_i=\nabla f_i(x),
\]
dan, bersyarat pada informasi masa lalu $\mathcal F$, indeks $I$ dipilih dengan
peluang $p_i>0$, $\sum_i p_i=1$. Peluang dan titik $x$ boleh
$\mathcal F$-terukur.

\begin{enumerate}
  \item Hitung bias penduga naif $H=g_I$ terhadap $\nabla F(x)$.
  \item Tunjukkan bahwa $G=g_I/(Np_I)$ tidak bias secara bersyarat.
  \item Turunkan rumus eksak untuk
  $\mathbb E[\lVert G-\nabla F(x)\rVert^2\mid\mathcal F]$.
\end{enumerate}
\end{exercise}

% stable-id: d90.orig.v1.tr03.ps04.hint.0001
\paragraph{Petunjuk bertahap.}
\label{orig03:ps04:hint:0001}
\textbf{Tahap 1.} Tulis ekspektasi bersyarat sebagai jumlah atas semua nilai
yang mungkin bagi $I$.
\textbf{Tahap 2.} Setelah memperoleh ketakbiasan, gunakan identitas
$\mathbb E\lVert G-m\rVert^2=\mathbb E\lVert G\rVert^2-\lVert m\rVert^2$
dengan $m=\mathbb E[G\mid\mathcal F]$.

% stable-id: d90.orig.v1.tr03.ps04.answer.0001
\paragraph{Jawaban singkat.}
\label{orig03:ps04:answer:0001}
Penduga naif mempunyai bias
$\sum_i p_i g_i-N^{-1}\sum_i g_i$. Sementara itu,
\[
 \mathbb E[G\mid\mathcal F]=\nabla F(x),\qquad
 \mathbb E[\lVert G-\nabla F(x)\rVert^2\mid\mathcal F]
 =\frac1{N^2}\sum_{i=1}^N\frac{\lVert g_i\rVert^2}{p_i}
  -\lVert\nabla F(x)\rVert^2.
\]

% stable-id: d90.orig.v1.tr03.ps04.solution.0001
\paragraph{Solusi lengkap.}
\label{orig03:ps04:solution:0001}
Karena $p_i$ dan $g_i$ diketahui setelah mengondisikan pada $\mathcal F$,
\[
 \mathbb E[H\mid\mathcal F]=\sum_i p_i g_i.
\]
Jadi selisihnya dari gradien penuh adalah
$\sum_i(p_i-1/N)g_i$, yang umumnya tidak nol. Untuk penduga terkoreksi,
\[
 \mathbb E[G\mid\mathcal F]
 =\sum_i p_i\frac{g_i}{Np_i}
 =\frac1N\sum_i g_i=\nabla F(x).
\]
Selanjutnya,
\[
 \mathbb E[\lVert G\rVert^2\mid\mathcal F]
 =\sum_i p_i\frac{\lVert g_i\rVert^2}{N^2p_i^2}
 =\frac1{N^2}\sum_i\frac{\lVert g_i\rVert^2}{p_i}.
\]
Mengurangkan kuadrat norma mean bersyarat memberi rumus varians yang
dinyatakan pada jawaban singkat. Positivitas setiap $p_i$ penting agar koreksi
terdefinisi dan tidak menghilangkan komponen mana pun.

% stable-id: d90.orig.v1.tr03.ps04.exercise.0002
\begin{exercise}[Ketaksamaan tiga titik untuk langkah Bregman komposit]
\label{orig03:ps04:exercise:0002}
Pada himpunan konveks tertutup $C$, misalkan $h$ terdiferensialkan dan konveks,
$r$ konveks proper semikontinu-bawah, dan
\[
 x^+\in\arg\min_{x\in C}
 \left\{\eta\langle g,x\rangle+\eta r(x)+D_h(x,x_k)\right\},
 \qquad \eta>0.
\]
Buktikan bahwa untuk setiap $u\in C$,
\[
 \eta\bigl(\langle g,x^+-u\rangle+r(x^+)-r(u)\bigr)
 \le D_h(u,x_k)-D_h(u,x^+)-D_h(x^+,x_k).
\]
Nyatakan dengan jelas kondisi optimalitas yang Anda gunakan.
\end{exercise}

% stable-id: d90.orig.v1.tr03.ps04.hint.0002
\paragraph{Petunjuk bertahap.}
\label{orig03:ps04:hint:0002}
\textbf{Tahap 1.} Pilih $s^+\in\partial r(x^+)$ dan
$n^+\in N_C(x^+)$ dari inklusi optimalitas.
\textbf{Tahap 2.} Gunakan ketaksamaan subgradien dan identitas tiga titik
untuk divergensi Bregman.

% stable-id: d90.orig.v1.tr03.ps04.answer.0002
\paragraph{Jawaban singkat.}
\label{orig03:ps04:answer:0002}
Inklusi
\[
 0\in\eta g+\eta\partial r(x^+)+\nabla h(x^+)-\nabla h(x_k)+N_C(x^+)
\]
bersama dengan
\[
 \langle\nabla h(x^+)-\nabla h(x_k),u-x^+\rangle
 =D_h(u,x_k)-D_h(u,x^+)-D_h(x^+,x_k)
\]
memberikan hasilnya.

% stable-id: d90.orig.v1.tr03.ps04.solution.0002
\paragraph{Solusi lengkap.}
\label{orig03:ps04:solution:0002}
Ambil $s^+\in\partial r(x^+)$ dan $n^+\in N_C(x^+)$ sehingga
\[
 0=\eta g+\eta s^++\nabla h(x^+)-\nabla h(x_k)+n^+.
\]
Karena $u\in C$, definisi kerucut normal memberi
$\langle n^+,u-x^+\rangle\le0$. Mengalikan persamaan optimalitas dengan
$x^+-u$ menghasilkan
\[
 \eta\langle g,x^+-u\rangle+\eta\langle s^+,x^+-u\rangle
 \le \langle\nabla h(x^+)-\nabla h(x_k),u-x^+\rangle.
\]
Ketaksamaan subgradien menyatakan
$\langle s^+,x^+-u\rangle\ge r(x^+)-r(u)$. Ruas kanan sama dengan selisih
tiga divergensi Bregman pada jawaban singkat. Substitusi kedua fakta itu
menyelesaikan bukti.

% stable-id: d90.orig.v1.tr03.ps04.exercise.0003
\begin{exercise}[Varians terbatas bukan momen kedua seragam]
\label{orig03:ps04:exercise:0003}
Pertimbangkan $f(x)=x^2/2$ dan orakel stokastik
$G(x,\xi)=x+\xi$, dengan $\mathbb E\xi=0$ dan
$\mathbb E\xi^2=\sigma^2<\infty$.

\begin{enumerate}
  \item Verifikasi ketakbiasan dan varians bersyarat yang seragam.
  \item Tunjukkan bahwa tidak ada batas seragam untuk
  $\mathbb E|G(x,\xi)|^2$ ketika $x$ menjelajahi $\mathbb R$.
  \item Buktikan arah sebaliknya: jika
  $\mathbb E\lVert G(x,\xi)\rVert^2\le M^2$ dan $G$ tidak bias, maka
  $\lVert\nabla f(x)\rVert\le M$.
\end{enumerate}
\end{exercise}

% stable-id: d90.orig.v1.tr03.ps04.hint.0003
\paragraph{Petunjuk bertahap.}
\label{orig03:ps04:hint:0003}
\textbf{Tahap 1.} Pisahkan $G-\nabla f$ dari $G$ sendiri.
\textbf{Tahap 2.} Untuk arah sebaliknya, terapkan ketaksamaan Jensen pada
fungsi kuadrat norma.

% stable-id: d90.orig.v1.tr03.ps04.answer.0003
\paragraph{Jawaban singkat.}
\label{orig03:ps04:answer:0003}
$G$ tidak bias dan
$\mathbb E|G-\nabla f(x)|^2=\sigma^2$, tetapi
$\mathbb E|G|^2=x^2+\sigma^2$, sehingga tak terbatas seragam. Sebaliknya,
$\lVert\nabla f(x)\rVert^2=\lVert\mathbb E G\rVert^2
\le\mathbb E\lVert G\rVert^2\le M^2$.

% stable-id: d90.orig.v1.tr03.ps04.solution.0003
\paragraph{Solusi lengkap.}
\label{orig03:ps04:solution:0003}
Gradien eksak ialah $\nabla f(x)=x$. Karena mean $\xi$ nol,
$\mathbb E[G(x,\xi)]=x$, sedangkan
$G(x,\xi)-\nabla f(x)=\xi$. Jadi varians galat orakel selalu
$\sigma^2$, tidak bergantung pada $x$. Namun,
\[
 \mathbb E|x+\xi|^2=x^2+2x\mathbb E\xi+\mathbb E\xi^2
 =x^2+\sigma^2,
\]
yang menuju tak hingga saat $|x|\to\infty$. Dengan demikian, asumsi varians
terbatas tidak boleh diam-diam diganti dengan asumsi momen kedua seragam.
Untuk arah terakhir, ketakbiasan dan konveksitas kuadrat norma memberi
\[
 \lVert\nabla f(x)\rVert^2
 =\lVert\mathbb E G(x,\xi)\rVert^2
 \le\mathbb E\lVert G(x,\xi)\rVert^2\le M^2.
\]

% stable-id: d90.orig.v1.tr03.ps04.exercise.0004
\begin{exercise}[Bias bersyarat pada pengacakan ulang]
\label{orig03:ps04:exercise:0004}
Dalam satu epok, ambil permutasi seragam $\pi$ dari
$\{1,\ldots,N\}$. Setelah indeks
$\pi_1,\ldots,\pi_t$ terungkap, misalkan $R_t$ adalah himpunan indeks yang
tersisa dan $\mathcal F_t$ memuat riwayat tersebut.

\begin{enumerate}
  \item Untuk titik beku $x$, hitung
  $\mathbb E[\nabla f_{\pi_{t+1}}(x)\mid\mathcal F_t]$.
  \item Tunjukkan bahwa jumlah gradien selama satu epok pada titik beku tepat
  sama dengan $N\nabla F(x)$.
  \item Jelaskan mengapa bukti SGD yang mengandalkan selisih martingal tidak
  dapat langsung dipakai ketika titik iterasi berubah di dalam epok.
\end{enumerate}
\end{exercise}

% stable-id: d90.orig.v1.tr03.ps04.hint.0004
\paragraph{Petunjuk bertahap.}
\label{orig03:ps04:hint:0004}
\textbf{Tahap 1.} Bersyarat pada riwayat, indeks berikutnya seragam hanya pada
$R_t$, bukan pada seluruh himpunan indeks.
\textbf{Tahap 2.} Bedakan identitas kombinatorial pada satu titik $x$ dari
jumlah $\nabla f_{\pi_t}(x_t)$ pada titik-titik yang berlainan.

% stable-id: d90.orig.v1.tr03.ps04.answer.0004
\paragraph{Jawaban singkat.}
\label{orig03:ps04:answer:0004}
\[
 \mathbb E[\nabla f_{\pi_{t+1}}(x)\mid\mathcal F_t]
 =\frac1{|R_t|}\sum_{i\in R_t}\nabla f_i(x),
 \qquad
 \sum_{t=1}^N\nabla f_{\pi_t}(x)=N\nabla F(x).
\]
Mean bersyarat pertama umumnya bukan gradien penuh; bila $x_t$ berubah,
identitas jumlah pada titik beku juga tidak berlaku.

% stable-id: d90.orig.v1.tr03.ps04.solution.0004
\paragraph{Solusi lengkap.}
\label{orig03:ps04:solution:0004}
Setelah $t$ indeks terungkap, setiap elemen $R_t$ mempunyai peluang
$1/|R_t|$ untuk muncul berikutnya. Karena itu mean bersyarat adalah rata-rata
gradien komponen yang tersisa. Rata-rata ini hanya sama dengan gradien penuh
dalam keadaan khusus. Di sisi lain, sebuah permutasi memuat setiap indeks
tepat sekali, sehingga untuk satu titik yang sama,
\[
 \sum_{t=1}^N\nabla f_{\pi_t}(x)=\sum_{i=1}^N\nabla f_i(x)=N\nabla F(x).
\]
Pada algoritme sebenarnya, suku ke-$t$ dinilai di $x_t$, sehingga jumlahnya
adalah $\sum_t\nabla f_{\pi_t}(x_t)$ dan tidak lagi dapat diganti oleh jumlah
gradien pada satu titik. Selain itu, galat gradien terhadap gradien penuh
umumnya mempunyai mean bersyarat tak nol. Maka argumen selisih martingal SGD
dengan pengambilan sampel independen memerlukan analisis baru, bukan sekadar
penggantian notasi.

% stable-id: d90.orig.v1.tr03.ps04.exercise.0005
\begin{exercise}[Lantai galat akibat derau persisten]
\label{orig03:ps04:exercise:0005}
Misalkan barisan tak negatif memenuhi
\[
 a_{k+1}\le(1-\mu\eta)a_k+\eta^2\sigma^2,
 \qquad 0<\mu\eta<1.
\]
Turunkan batas tertutup untuk $a_k$, hitung limit superiornya, dan jelaskan
apa yang harus diubah agar jaminan konvergensi eksak tetap masuk akal ketika
$\sigma>0$.
\end{exercise}

% stable-id: d90.orig.v1.tr03.ps04.hint.0005
\paragraph{Petunjuk bertahap.}
\label{orig03:ps04:hint:0005}
\textbf{Tahap 1.} Tetapkan $q=1-\mu\eta$ dan buka rekursi sebagai deret
geometrik.
\textbf{Tahap 2.} Bandingkan faktor kontraksi $q^k$ dengan suku paksa yang
tidak lenyap.

% stable-id: d90.orig.v1.tr03.ps04.answer.0005
\paragraph{Jawaban singkat.}
\label{orig03:ps04:answer:0005}
Dengan $q=1-\mu\eta$,
\[
 a_k\le q^k a_0+\frac{\eta\sigma^2}{\mu}(1-q^k),
 \qquad
 \limsup_{k\to\infty}a_k\le\frac{\eta\sigma^2}{\mu}.
\]
Langkah tetap dengan derau persisten hanya menjamin lingkungan solusi;
langkah menurun atau reduksi varians diperlukan untuk meniadakan lantai ini.

% stable-id: d90.orig.v1.tr03.ps04.solution.0005
\paragraph{Solusi lengkap.}
\label{orig03:ps04:solution:0005}
Iterasi rekursi memberikan
\[
 a_k\le q^k a_0+\eta^2\sigma^2\sum_{j=0}^{k-1}q^j
 =q^k a_0+\eta^2\sigma^2\frac{1-q^k}{1-q}.
\]
Karena $1-q=\mu\eta$, suku kedua sama dengan
$\eta\sigma^2(1-q^k)/\mu$. Hipotesis memastikan $0<q<1$, sehingga
$q^k\to0$ dan batas limit superior mengikuti. Bila $\eta$ tetap dan
$\sigma>0$, suku paksa tidak hilang; rekursi itu sendiri tidak dapat
membuktikan $a_k\to0$. Jadwal langkah yang menuju nol dengan syarat penjumlahan
yang sesuai, atau penduga dengan varians yang ikut menuju nol, mengubah
rekursi sehingga lantai galat dapat lenyap.

% segment-local-id: OR03-S007
% segment-id: d90.orig.v1.tr03.seg0007
% license: CC BY-SA 4.0
% provenance: Materi asesmen asli berbahasa Indonesia, disusun secara independen.

\section{Set soal V: operator monoton dan metode pemisahan}
\label{orig03:ps05}

% stable-id: d90.orig.v1.tr03.ps05.exercise.0001
\begin{exercise}[Keberadaan, ketunggalan, dan titik tetap maju--mundur]
\label{orig03:ps05:exercise:0001}
Pada ruang Euklides berdimensi hingga, misalkan $f$ proper, semikontinu-bawah,
dan $\mu$-konveks kuat dengan $\mu>0$. Misalkan $g$ konveks dan
terdiferensialkan pada seluruh ruang, dengan $\nabla g$ bersifat
$\beta$-kokersif. Buktikan bahwa $f+g$ mempunyai tepat satu peminimum $x^\star$
dan bahwa, untuk setiap $\gamma>0$,
\[
 x^\star=J_{\gamma\partial f}
 \bigl(x^\star-\gamma\nabla g(x^\star)\bigr).
\]
Jelaskan peran aturan penjumlahan subdiferensial dalam argumen Anda.
\end{exercise}

% stable-id: d90.orig.v1.tr03.ps05.hint.0001
\paragraph{Petunjuk bertahap.}
\label{orig03:ps05:hint:0001}
\textbf{Tahap 1.} Gunakan batas bawah kuadratik dari konveksitas kuat untuk
mendapatkan kekompakan himpunan sublevel.
\textbf{Tahap 2.} Ubah inklusi
$0\in\partial f(x^\star)+\nabla g(x^\star)$ menjadi definisi resolven.

% stable-id: d90.orig.v1.tr03.ps05.answer.0001
\paragraph{Jawaban singkat.}
\label{orig03:ps05:answer:0001}
$f+g$ proper, semikontinu-bawah, koersif, dan $\mu$-konveks kuat; jadi ia
mempunyai peminimum tunggal. Kontinuitas $g$ memberi
$\partial(f+g)=\partial f+\nabla g$, dan syarat Fermat ekuivalen dengan
\[
 x^\star-\gamma\nabla g(x^\star)
 \in x^\star+\gamma\partial f(x^\star),
\]
yakni persamaan titik tetap yang diminta.

% stable-id: d90.orig.v1.tr03.ps05.solution.0001
\paragraph{Solusi lengkap.}
\label{orig03:ps05:solution:0001}
Konveksitas kuat memberikan batas bawah kuadratik bagi $f$ di sekitar suatu
titik dalam domainnya. Fungsi konveks terdiferensialkan $g$ mempunyai batas
bawah afin dari bidang singgungnya. Karena itu suku kuadratik mendominasi
batas afin ketika norma membesar, sehingga $f+g$ koersif. Fungsi tersebut
proper dan semikontinu-bawah, maka dalam dimensi hingga ia mencapai minimum.
Jumlahnya tetap $\mu$-konveks kuat, sehingga dua peminimum yang berbeda tidak
mungkin ada.

Karena $g$ bernilai hingga dan kontinu di seluruh ruang, aturan penjumlahan
berlaku tanpa syarat kualifikasi tambahan:
\[
 \partial(f+g)(x)=\partial f(x)+\nabla g(x).
\]
Syarat Fermat untuk peminimum tunggal adalah
$-\nabla g(x^\star)\in\partial f(x^\star)$. Mengalikan dengan $\gamma$ dan
menata ulang memberi
$x^\star-\gamma\nabla g(x^\star)\in(I+\gamma\partial f)(x^\star)$.
Menerapkan invers $J_{\gamma\partial f}=(I+\gamma\partial f)^{-1}$ memberi
persamaan titik tetap.

% stable-id: d90.orig.v1.tr03.ps05.exercise.0002
\begin{exercise}[Keaveragan operator maju--mundur dan laju residu]
\label{orig03:ps05:exercise:0002}
Misalkan $A$ maksimal monoton dan $B$ bersifat $\beta$-kokersif. Untuk
$0<\gamma<2\beta$, definisikan
\[
 T=J_{\gamma A}(I-\gamma B),
 \qquad x_{k+1}=Tx_k,
\]
dan anggap $\operatorname{Fix}T\ne\varnothing$.

\begin{enumerate}
  \item Tunjukkan bahwa $I-\gamma B$ bersifat
  $\gamma/(2\beta)$-averaged dan $J_{\gamma A}$ bersifat $1/2$-averaged.
  \item Gunakan aturan komposisi untuk memperoleh
  $\alpha=2\beta/(4\beta-\gamma)$ sebagai parameter averaged bagi $T$.
  \item Untuk $z\in\operatorname{Fix}T$, turunkan ketaksamaan penurunan dan
  batas untuk residu terkecil di antara $k$ iterasi pertama.
\end{enumerate}
\end{exercise}

% stable-id: d90.orig.v1.tr03.ps05.hint.0002
\paragraph{Petunjuk bertahap.}
\label{orig03:ps05:hint:0002}
\textbf{Tahap 1.} Tulis $I-\gamma B=(1-a)I+a(I-2\beta B)$ dengan
$a=\gamma/(2\beta)$.
\textbf{Tahap 2.} Terapkan karakterisasi kuadrat operator averaged, lalu
jumlahkan ketaksamaan itu dari $0$ sampai $k-1$.

% stable-id: d90.orig.v1.tr03.ps05.answer.0002
\paragraph{Jawaban singkat.}
\label{orig03:ps05:answer:0002}
Komposisi kedua operator tersebut bersifat
$\alpha$-averaged dengan
$\alpha=2\beta/(4\beta-\gamma)$. Untuk setiap titik tetap $z$,
\[
 \lVert Tx_k-z\rVert^2
 \le\lVert x_k-z\rVert^2
 -\frac{1-\alpha}{\alpha}\lVert(I-T)x_k\rVert^2,
\]
dan karenanya
\[
 \min_{0\le j<k}\lVert(I-T)x_j\rVert
 \le\sqrt{\frac{\alpha}{(1-\alpha)k}}\,\lVert x_0-z\rVert.
\]

% stable-id: d90.orig.v1.tr03.ps05.solution.0002
\paragraph{Solusi lengkap.}
\label{orig03:ps05:solution:0002}
Kokersivitas $B$ menyiratkan bahwa $I-2\beta B$ non-ekspansif. Dengan
$a=\gamma/(2\beta)\in(0,1)$,
\[
 I-\gamma B=(1-a)I+a(I-2\beta B),
\]
sehingga operator maju bersifat $a$-averaged. Resolven operator maksimal
monoton firmly non-ekspansif, yang sama dengan $1/2$-averaged. Jika dua
operator masing-masing $a_1$- dan $a_2$-averaged, komposisinya dapat diambil
averaged dengan parameter
\[
 \frac{a_1+a_2-2a_1a_2}{1-a_1a_2}.
\]
Memasukkan $a_1=1/2$ dan $a_2=\gamma/(2\beta)$ menghasilkan
$\alpha=2\beta/(4\beta-\gamma)\in(0,1)$.

Karakterisasi kuadrat operator $\alpha$-averaged memberi, karena $Tz=z$,
\[
 \lVert Tx_k-z\rVert^2
 \le\lVert x_k-z\rVert^2
 -\frac{1-\alpha}{\alpha}\lVert(I-T)x_k\rVert^2.
\]
Penjumlahan meneleskopkan ruas norma dan menunjukkan
\[
 \sum_{j=0}^{k-1}\lVert(I-T)x_j\rVert^2
 \le\frac{\alpha}{1-\alpha}\lVert x_0-z\rVert^2.
\]
Nilai minimum tidak melebihi rata-rata; mengambil akar memberikan batas pada
jawaban singkat dan menunjukkan residu terbaik berorde $O(k^{-1/2})$.

% stable-id: d90.orig.v1.tr03.ps05.exercise.0003
\begin{exercise}[Jendela langkah untuk operator monoton tak simetris]
\label{orig03:ps05:exercise:0003}
Untuk $a>0$, pertimbangkan operator linear
\[
 B_a=\begin{pmatrix}1&-a\\a&1\end{pmatrix}
 \quad\text{pada }\mathbb R^2.
\]
Buktikan bahwa $B_a$ bersifat $1$-monoton kuat, mempunyai konstanta Lipschitz
$\sqrt{1+a^2}$, dan bersifat $1/(1+a^2)$-kokersif. Tentukan secara eksak
kapan iterasi maju $x_{k+1}=(I-\gamma B_a)x_k$ merupakan kontraksi. Gunakan
hasil itu untuk mendiagnosis klaim bahwa monotoni kuat saja selalu membolehkan
semua $0<\gamma<2$.
\end{exercise}

% stable-id: d90.orig.v1.tr03.ps05.hint.0003
\paragraph{Petunjuk bertahap.}
\label{orig03:ps05:hint:0003}
\textbf{Tahap 1.} Untuk $d\in\mathbb R^2$, hitung
$\langle B_ad,d\rangle$ dan $\lVert B_ad\rVert^2$.
\textbf{Tahap 2.} Hitung norma faktor
$(1-\gamma)I-\gamma aJ$, dengan $J$ rotasi seperempat putaran.

% stable-id: d90.orig.v1.tr03.ps05.answer.0003
\paragraph{Jawaban singkat.}
\label{orig03:ps05:answer:0003}
Untuk setiap $d$,
$\langle B_ad,d\rangle=\lVert d\rVert^2$ dan
$\lVert B_ad\rVert^2=(1+a^2)\lVert d\rVert^2$. Faktor kontraksinya adalah
\[
 \rho(\gamma)=\sqrt{(1-\gamma)^2+a^2\gamma^2},
\]
sehingga kontraksi terjadi tepat ketika
$0<\gamma<2/(1+a^2)$.

% stable-id: d90.orig.v1.tr03.ps05.solution.0003
\paragraph{Solusi lengkap.}
\label{orig03:ps05:solution:0003}
Bagian simetris $B_a$ adalah identitas. Karena itu
$\langle B_ad,d\rangle=\lVert d\rVert^2$, yang membuktikan monotoni kuat
dengan modulus satu. Perkalian langsung memberi
$B_a^\top B_a=(1+a^2)I$; jadi norma operatornya
$\sqrt{1+a^2}$ dan
\[
 \langle B_ad,d\rangle
 =\frac1{1+a^2}\lVert B_ad\rVert^2,
\]
yakni kokersivitas yang dinyatakan.

Matriks iterasi mempunyai bagian skalar $1-\gamma$ dan bagian rotasi
$-\gamma aJ$, sehingga norma setiap vektor dikalikan dengan
$\rho(\gamma)$. Syarat $\rho(\gamma)<1$ ekuivalen dengan
\[
 (1-\gamma)^2+a^2\gamma^2<1
 \quad\Longleftrightarrow\quad
 0<\gamma<\frac{2}{1+a^2}.
\]
Jadi interval $(0,2)$ tidak seragam terhadap bagian antisimetri. Monotoni kuat
tanpa kontrol Lipschitz atau kokersivitas tidak cukup untuk membenarkan klaim
langkah tersebut.

% stable-id: d90.orig.v1.tr03.ps05.exercise.0004
\begin{exercise}[Sertifikat titik tetap Douglas--Rachford]
\label{orig03:ps05:exercise:0004}
Pada $\mathbb R$, minimalkan $f(x)+g(x)$ dengan
$f(x)=|x|$ dan $g=\delta_{[1,\infty)}$. Ambil $\gamma=1$ dan konvensi
\[
 T=\tfrac12(I+R_gR_f),
 \qquad R_q=2\operatorname{prox}_q-I.
\]
Tentukan peminimum, lalu verifikasi secara langsung bahwa $z^\star=2$ adalah
titik tetap $T$ dan bahwa bayangannya
$\operatorname{prox}_f(z^\star)$ sama dengan peminimum.
\end{exercise}

% stable-id: d90.orig.v1.tr03.ps05.hint.0004
\paragraph{Petunjuk bertahap.}
\label{orig03:ps05:hint:0004}
\textbf{Tahap 1.} Gunakan penyusutan lunak untuk $\operatorname{prox}_{|\cdot|}$
dan proyeksi untuk $\operatorname{prox}_{\delta_{[1,\infty)}}$.
\textbf{Tahap 2.} Ikuti urutan $2\mapsto R_f(2)\mapsto R_gR_f(2)$ tanpa
menukar konvensi komposisi.

% stable-id: d90.orig.v1.tr03.ps05.answer.0004
\paragraph{Jawaban singkat.}
\label{orig03:ps05:answer:0004}
Peminimumnya $x^\star=1$. Selain itu,
$\operatorname{prox}_f(2)=1$, $R_f(2)=0$,
$\operatorname{prox}_g(0)=1$, dan $R_g(0)=2$. Maka $T(2)=2$ dan bayangan
$\operatorname{prox}_f(2)=x^\star$.

% stable-id: d90.orig.v1.tr03.ps05.solution.0004
\paragraph{Solusi lengkap.}
\label{orig03:ps05:solution:0004}
Kendala efektif $g$ memaksa $x\ge1$. Pada himpunan itu $|x|=x$, sehingga
nilai terkecil dicapai secara tunggal di $x^\star=1$. Proksimal $f$ dengan
parameter satu adalah penyusutan lunak; akibatnya
$\operatorname{prox}_f(2)=2-1=1$ dan $R_f(2)=2(1)-2=0$. Proksimal indikator
adalah proyeksi ke $[1,\infty)$, sehingga
$\operatorname{prox}_g(0)=1$ dan $R_g(0)=2(1)-0=2$. Oleh sebab itu
\[
 T(2)=\tfrac12(2+2)=2.
\]
Variabel titik tetap $z^\star$ tidak harus sama dengan solusi primal; pada
konvensi ini solusi diperoleh dari bayangan proksimalnya, yang memang bernilai
satu.

% stable-id: d90.orig.v1.tr03.ps05.exercise.0005
\begin{exercise}[Metode titik proksimal dengan galat yang dapat dijumlahkan]
\label{orig03:ps05:exercise:0005}
Misalkan $A$ maksimal monoton pada ruang Euklides berdimensi hingga,
$\operatorname{zer}A\ne\varnothing$, dan $\gamma>0$. Pertimbangkan
\[
 x_{k+1}=J_{\gamma A}x_k+e_k,
 \qquad \sum_{k=0}^{\infty}\lVert e_k\rVert<\infty.
\]
Buktikan bahwa $(x_k)$ konvergen ke suatu titik di $\operatorname{zer}A$.
Dalam bukti Anda, tunjukkan juga bahwa
$x_k-J_{\gamma A}x_k\to0$.
\end{exercise}

% stable-id: d90.orig.v1.tr03.ps05.hint.0005
\paragraph{Petunjuk bertahap.}
\label{orig03:ps05:hint:0005}
\textbf{Tahap 1.} Tetapkan $y_k=J_{\gamma A}x_k$ dan gunakan
firm non-ekspansivitas terhadap $z\in\operatorname{zer}A$.
\textbf{Tahap 2.} Jumlahkan galat kuadrat yang dihasilkan; lalu gunakan
ketertutupan graf operator maksimal monoton pada subsekuens konvergen.

% stable-id: d90.orig.v1.tr03.ps05.answer.0005
\paragraph{Jawaban singkat.}
\label{orig03:ps05:answer:0005}
Firm non-ekspansivitas memberi
\[
 \lVert y_k-z\rVert^2+\lVert x_k-y_k\rVert^2
 \le\lVert x_k-z\rVert^2.
\]
Perturbasi yang normanya dapat dijumlahkan membuat barisan kuasi-Fej\'{e}r,
dan penjumlahan ketaksamaan memperlihatkan
$\sum_k\lVert x_k-y_k\rVert^2<\infty$. Semua titik klaster adalah nol $A$;
dalam dimensi hingga sifat kuasi-Fej\'{e}r lalu memberi konvergensi seluruh
barisan.

% stable-id: d90.orig.v1.tr03.ps05.solution.0005
\paragraph{Solusi lengkap.}
\label{orig03:ps05:solution:0005}
Ambil $z\in\operatorname{zer}A$. Maka $J_{\gamma A}z=z$. Dengan
$y_k=J_{\gamma A}x_k$, firm non-ekspansivitas menghasilkan
\[
 \lVert y_k-z\rVert^2+\lVert x_k-y_k\rVert^2
 \le\lVert x_k-z\rVert^2.                                      \tag{1}
\]
Karena $x_{k+1}=y_k+e_k$,
$\lVert x_{k+1}-z\rVert\le\lVert x_k-z\rVert+\lVert e_k\rVert$.
Jumlah norma galat hingga, sehingga $(x_k)$ dan $(y_k)$ terbatas dan jarak
terhadap setiap nol mempunyai limit.

Mengembangkan kuadrat $\lVert y_k+e_k-z\rVert^2$ lalu memakai (1) memberi
\[
 \lVert x_{k+1}-z\rVert^2
 \le\lVert x_k-z\rVert^2-\lVert x_k-y_k\rVert^2
   +2\langle y_k-z,e_k\rangle+\lVert e_k\rVert^2.
\]
Karena $(y_k)$ terbatas dan galat dapat dijumlahkan, kedua suku galat di ruas
kanan dapat dijumlahkan. Penjumlahan ketaksamaan menunjukkan
$\sum_k\lVert x_k-y_k\rVert^2<\infty$, jadi residu itu menuju nol.

Setiap subsekuens konvergen $x_{k_j}\to\bar x$ juga memenuhi
$y_{k_j}\to\bar x$. Dari definisi resolven,
$(x_{k_j}-y_{k_j})/\gamma\in Ay_{k_j}$, dan ruas kiri menuju nol. Graf
operator maksimal monoton tertutup, sehingga $0\in A\bar x$. Jadi semua titik
klaster berada di $\operatorname{zer}A$. Barisan terbatas dalam dimensi hingga
mempunyai titik klaster, dan sifat kuasi-Fej\'{e}r beserta fakta bahwa semua
titik klasternya adalah nol memaksa seluruh barisan konvergen ke satu nol
$A$.

% segment-local-id: OR03-S008
% segment-id: d90.orig.v1.tr03.seg0008
% license: CC BY-SA 4.0
% provenance: Materi asesmen asli berbahasa Indonesia, disusun secara independen.

\section{Set soal VI: transportasi optimal dan sintesis}
\label{orig03:ps06}

% stable-id: d90.orig.v1.tr03.ps06.exercise.0001
\begin{exercise}[Stabilitas terhadap perturbasi biaya seragam]
\label{orig03:ps06:exercise:0001}
Misalkan $\Pi(\alpha,\beta)$ adalah himpunan kopling probabilitas dan dua
fungsi biaya memenuhi
$\lVert c-\widetilde c\rVert_\infty\le\varepsilon$. Tuliskan $v$ dan
$\widetilde v$ untuk nilai optimal dengan biaya $c$ dan $\widetilde c$.

\begin{enumerate}
  \item Buktikan $|v-\widetilde v|\le\varepsilon$.
  \item Jika $P_c$ optimal untuk biaya $c$, buktikan bahwa ia paling jauh
  $2\varepsilon$-suboptimal untuk biaya $\widetilde c$.
\end{enumerate}
\end{exercise}

% stable-id: d90.orig.v1.tr03.ps06.hint.0001
\paragraph{Petunjuk bertahap.}
\label{orig03:ps06:hint:0001}
\textbf{Tahap 1.} Integralkan batas titik demi titik terhadap sembarang kopling
probabilitas.
\textbf{Tahap 2.} Sisipkan nilai $v$ di antara biaya
$\widetilde c$ pada $P_c$ dan nilai $\widetilde v$.

% stable-id: d90.orig.v1.tr03.ps06.answer.0001
\paragraph{Jawaban singkat.}
\label{orig03:ps06:answer:0001}
Untuk setiap $P\in\Pi(\alpha,\beta)$,
$|\langle c,P\rangle-\langle\widetilde c,P\rangle|\le\varepsilon$.
Mengambil infimum memberi $|v-\widetilde v|\le\varepsilon$, dan
\[
 \langle\widetilde c,P_c\rangle-\widetilde v\le2\varepsilon.
\]

% stable-id: d90.orig.v1.tr03.ps06.solution.0001
\paragraph{Solusi lengkap.}
\label{orig03:ps06:solution:0001}
Karena massa total setiap kopling adalah satu,
\[
 \left|\int(c-\widetilde c)\,dP\right|
 \le\int|c-\widetilde c|\,dP\le\varepsilon.
\]
Jadi $\langle\widetilde c,P\rangle\le\langle c,P\rangle+\varepsilon$
untuk setiap $P$. Mengambil infimum memberi
$\widetilde v\le v+\varepsilon$. Menukar peran kedua biaya memberi
$v\le\widetilde v+\varepsilon$, sehingga batas nilai optimal terbukti.
Selanjutnya,
\[
 \langle\widetilde c,P_c\rangle-\widetilde v
 \le\langle c,P_c\rangle+\varepsilon-(v-\varepsilon)
 =2\varepsilon,
\]
karena $\langle c,P_c\rangle=v$. Batas kedua memuat dua sumber galat: evaluasi
rencana dan pergeseran nilai optimal.

% stable-id: d90.orig.v1.tr03.ps06.exercise.0002
\begin{exercise}[Sertifikat dual untuk translasi]
\label{orig03:ps06:exercise:0002}
Pada $\mathbb R$, ambil biaya
$c(x,y)=\tfrac12(y-x)^2$. Misalkan $\alpha$ mempunyai momen pertama hingga dan
$\beta=(x\mapsto x+a)_\#\alpha$ untuk suatu $a\in\mathbb R$.

\begin{enumerate}
  \item Hitung biaya kopling deterministik $y=x+a$.
  \item Buktikan bahwa
  $\phi(x)=-ax$ dan $\psi(y)=ay-a^2/2$ layak untuk masalah dual Kantorovich.
  \item Gunakan kesetaraan primal--dual untuk menyimpulkan optimalitas kopling
  translasi tersebut.
\end{enumerate}
\end{exercise}

% stable-id: d90.orig.v1.tr03.ps06.hint.0002
\paragraph{Petunjuk bertahap.}
\label{orig03:ps06:hint:0002}
\textbf{Tahap 1.} Kurangi $\phi(x)+\psi(y)$ dari $c(x,y)$ dan lengkapi
kuadrat dalam $y-x$.
\textbf{Tahap 2.} Gunakan identitas
$\int y\,d\beta(y)=\int(x+a)\,d\alpha(x)$.

% stable-id: d90.orig.v1.tr03.ps06.answer.0002
\paragraph{Jawaban singkat.}
\label{orig03:ps06:answer:0002}
Biaya kopling translasi adalah $a^2/2$, dan
\[
 c(x,y)-\phi(x)-\psi(y)=\tfrac12(y-x-a)^2\ge0.
\]
Nilai dual kedua potensial juga $a^2/2$, sehingga kopling dan potensial itu
optimal.

% stable-id: d90.orig.v1.tr03.ps06.solution.0002
\paragraph{Solusi lengkap.}
\label{orig03:ps06:solution:0002}
Pada kopling $y=x+a$, biaya bernilai konstan
$\tfrac12a^2$, sehingga integralnya $a^2/2$. Untuk semua $x,y$,
\[
 \tfrac12(y-x)^2-\bigl(-ax+ay-a^2/2\bigr)
 =\tfrac12(y-x-a)^2\ge0.
\]
Maka $\phi(x)+\psi(y)\le c(x,y)$, dengan kesetaraan tepat pada graf translasi.
Momen pertama menjamin integral potensial terdefinisi, dan
\[
 \int\phi\,d\alpha+\int\psi\,d\beta
 =-a\int x\,d\alpha
  +a\int(x+a)\,d\alpha-a^2/2
 =a^2/2.
\]
Dual lemah menyatakan bahwa nilai ini tidak dapat melebihi biaya kopling mana
pun. Karena kopling translasi mencapai nilai yang sama, celah dualnya nol dan
kopling itu optimal.

% stable-id: d90.orig.v1.tr03.ps06.exercise.0003
\begin{exercise}[Penskalaan Sinkhorn dalam domain log dan kebebasan gauge]
\label{orig03:ps06:exercise:0003}
Untuk marginal positif $a\in\mathbb R^n$ dan $b\in\mathbb R^m$, definisikan
\[
 K_{ij}=\exp(-C_{ij}/\varepsilon),
 \qquad P=\operatorname{diag}(u)K\operatorname{diag}(v),
 \]
dengan $u_i,v_j>0$. Tuliskan $\rho_i=\log u_i$ dan $\tau_j=\log v_j$.

\begin{enumerate}
  \item Turunkan pembaruan baris dan kolom dalam bentuk log-sum-exp.
  \item Buktikan bahwa $u\mapsto tu$, $v\mapsto v/t$ tidak mengubah $P$.
  \item Jelaskan secara matematis mengapa pemusatan $\rho$ dan $\tau$ dengan
  pergeseran yang berlawanan dapat mencegah luapan numerik tanpa mengubah
  kopling.
\end{enumerate}
\end{exercise}

% stable-id: d90.orig.v1.tr03.ps06.hint.0003
\paragraph{Petunjuk bertahap.}
\label{orig03:ps06:hint:0003}
\textbf{Tahap 1.} Selesaikan persamaan $P\1=a$ terhadap setiap $u_i$,
lalu ambil logaritma; lakukan hal yang sama untuk kolom.
\textbf{Tahap 2.} Tambahkan konstanta $s$ pada semua $\rho_i$ dan kurangkan
$s$ dari semua $\tau_j$.

% stable-id: d90.orig.v1.tr03.ps06.answer.0003
\paragraph{Jawaban singkat.}
\label{orig03:ps06:answer:0003}
\[
 \rho_i=\log a_i-\operatorname{LSE}_j(-C_{ij}/\varepsilon+\tau_j),
 \quad
 \tau_j=\log b_j-\operatorname{LSE}_i(-C_{ij}/\varepsilon+\rho_i).
\]
Transformasi $(\rho,\tau)\mapsto(\rho+s\1,\tau-s\1)$ membiarkan
setiap jumlah $\rho_i+\tau_j$ tetap, sehingga $P$ tidak berubah.

% stable-id: d90.orig.v1.tr03.ps06.solution.0003
\paragraph{Solusi lengkap.}
\label{orig03:ps06:solution:0003}
Kendala baris adalah
$u_i\sum_jK_{ij}v_j=a_i$, jadi
\[
 \rho_i=\log a_i-
 \log\sum_j\exp(-C_{ij}/\varepsilon+\tau_j).
\]
Logaritma jumlah eksponensial adalah fungsi $\operatorname{LSE}$. Dengan cara
yang sama, kendala kolom menghasilkan rumus untuk $\tau_j$ pada jawaban
singkat. Jika $t>0$, maka
\[
 \operatorname{diag}(tu)K\operatorname{diag}(v/t)
 =\operatorname{diag}(u)K\operatorname{diag}(v).
\]
Dalam koordinat log, ini tepat sama dengan menambahkan $s=\log t$ pada
$\rho$ dan mengurangkannya dari $\tau$. Karena entri $P$ bergantung pada
$\rho_i+\tau_j-C_{ij}/\varepsilon$, semua entri tetap sama. Kita boleh memilih
$s$, misalnya untuk memusatkan rentang nilai log, agar eksponensial perantara
tidak terlalu besar atau kecil. Operasi itu mengubah representasi, bukan
kopling atau marginalnya.

% stable-id: d90.orig.v1.tr03.ps06.exercise.0004
\begin{exercise}[Batas bias regularisasi entropi]
\label{orig03:ps06:exercise:0004}
Pada tabel kopling probabilitas berukuran $n\times m$, definisikan
$H(P)=-\sum_{ij}P_{ij}\log P_{ij}$ dengan konvensi $0\log0=0$.
Misalkan $P^\star$ meminimalkan $\langle C,P\rangle$ dan $P_\varepsilon$
meminimalkan
\[
 \langle C,P\rangle-\varepsilon H(P)
\]
pada himpunan kopling yang sama. Buktikan
\[
 0\le\langle C,P_\varepsilon\rangle-
       \langle C,P^\star\rangle
 \le\varepsilon\log(nm).
\]
\end{exercise}

% stable-id: d90.orig.v1.tr03.ps06.hint.0004
\paragraph{Petunjuk bertahap.}
\label{orig03:ps06:hint:0004}
\textbf{Tahap 1.} Bandingkan nilai objektif teratur pada $P_\varepsilon$ dan
$P^\star$.
\textbf{Tahap 2.} Gunakan $0\le H(P)\le\log(nm)$ untuk setiap distribusi pada
$nm$ sel.

% stable-id: d90.orig.v1.tr03.ps06.answer.0004
\paragraph{Jawaban singkat.}
\label{orig03:ps06:answer:0004}
Optimalitas kedua rencana memberi batas bawah nol dan
\[
 \langle C,P_\varepsilon\rangle-langle C,P^\star\rangle
 \le\varepsilon\bigl(H(P_\varepsilon)-H(P^\star)\bigr)
 \le\varepsilon\log(nm).
\]

% stable-id: d90.orig.v1.tr03.ps06.solution.0004
\paragraph{Solusi lengkap.}
\label{orig03:ps06:solution:0004}
Karena $P^\star$ optimal untuk biaya tanpa regularisasi,
$\langle C,P_\varepsilon\rangle\ge\langle C,P^\star\rangle$. Di sisi lain,
optimalitas $P_\varepsilon$ untuk objektif teratur menyatakan
\[
 \langle C,P_\varepsilon\rangle-\varepsilon H(P_\varepsilon)
 \le\langle C,P^\star\rangle-\varepsilon H(P^\star).
\]
Menata ulang memberi ketaksamaan pertama pada jawaban singkat. Entropi
distribusi dengan paling banyak $nm$ atom berada di antara nol dan
$\log(nm)$. Maka
$H(P_\varepsilon)-H(P^\star)\le\log(nm)$, yang menyelesaikan batas atas.
Batas ini mengukur bias biaya tak teratur, bukan jarak antara kedua rencana.

% stable-id: d90.orig.v1.tr03.ps06.exercise.0005
\begin{exercise}[Kelengkungan dual entropik dan arah gauge]
\label{orig03:ps06:exercise:0005}
Pertimbangkan fungsi dual
\[
 D(f,g)=a^\top f+b^\top g
 -\varepsilon\sum_{i,j}
 \exp\!\left(\frac{f_i+g_j-C_{ij}}{\varepsilon}\right),
\]
dan tuliskan
$P_{ij}=\exp((f_i+g_j-C_{ij})/\varepsilon)$.

\begin{enumerate}
  \item Hitung gradien terhadap $f$ dan $g$.
  \item Hitung bentuk kuadratik Hessian pada arah $(s,t)$ dan buktikan
  kekonkafan.
  \item Identifikasi arah gauge. Jika graf bipartit dukungan positif $P$
  terhubung, buktikan bahwa hanya arah gauge yang mempunyai kelengkungan nol.
\end{enumerate}
\end{exercise}

% stable-id: d90.orig.v1.tr03.ps06.hint.0005
\paragraph{Petunjuk bertahap.}
\label{orig03:ps06:hint:0005}
\textbf{Tahap 1.} Diferensiasikan satu eksponensial dan kelompokkan jumlah
baris serta kolom.
\textbf{Tahap 2.} Bentuk kuadratik dapat ditulis sebagai jumlah berbobot dari
$(s_i+t_j)^2$.

% stable-id: d90.orig.v1.tr03.ps06.answer.0005
\paragraph{Jawaban singkat.}
\label{orig03:ps06:answer:0005}
\[
 \nabla_fD=a-P\1,
 \qquad \nabla_gD=b-P^\top\1,
\]
dan
\[
 \begin{pmatrix}s\\t\end{pmatrix}^{\!\top}
 \nabla^2D
 \begin{pmatrix}s\\t\end{pmatrix}
 =-\frac1\varepsilon\sum_{i,j}P_{ij}(s_i+t_j)^2\le0.
\]
Arah $(s,t)=(c\1,-c\1)$ adalah arah gauge; bila dukungan
terhubung, semua arah nol berbentuk demikian.

% stable-id: d90.orig.v1.tr03.ps06.solution.0005
\paragraph{Solusi lengkap.}
\label{orig03:ps06:solution:0005}
Diferensiasi terhadap $f_i$ menghasilkan
$a_i-\sum_jP_{ij}$, sedangkan diferensiasi terhadap $g_j$ menghasilkan
$b_j-\sum_iP_{ij}$. Diferensiasi kedua pada lintasan
$(f+\theta s,g+\theta t)$ memberi
\[
 \left.\frac{d^2}{d\theta^2}D(f+\theta s,g+\theta t)\right|_{\theta=0}
 =-\frac1\varepsilon\sum_{i,j}P_{ij}(s_i+t_j)^2.
\]
Bobotnya nonnegatif, sehingga $D$ konkaf. Bentuk kuadratik nol tepat bila
$s_i+t_j=0$ pada setiap sisi dukungan positif. Pada graf bipartit yang
terhubung, kondisi ini memaksa semua $s_i$ sama dengan satu konstanta $c$ dan
semua $t_j=-c$. Itulah transformasi
$(f,g)\mapsto(f+c\1,g-c\1)$ yang tidak mengubah $P$. Jadi dual
bersifat konkaf ketat setelah satu gauge dipatok, meskipun tidak konkaf ketat
pada representasi yang belum dinormalisasi.

% segment-local-id: OR03-S009
% segment-id: d90.orig.v1.tr03.seg0009
% license: CC BY-SA 4.0
% provenance: Materi asesmen asli berbahasa Indonesia, disusun secara independen.

\section{Rubrik analitik untuk pembuktian}
\label{orig03:rubric}

Setiap kriteria di bawah dinilai dengan skala yang sama:
\textbf{0} tidak ada atau tidak sah;
\textbf{1} fragmen relevan dengan galat mayor;
\textbf{2} arah sebagian benar tetapi masih ada celah substantif;
\textbf{3} benar dengan satu kekurangan minor yang tidak mengubah kesimpulan;
dan \textbf{4} lengkap, tepat, dan seluruh dependensi logis dinyatakan.
Setiap rubrik mempunyai empat kriteria, jadi skor maksimumnya adalah 16.
Keanggunan gaya tidak boleh menutupi galat matematika, dan notasi alternatif
yang ekuivalen harus dinilai berdasarkan isinya.

% stable-id: d90.orig.v1.tr03.rubric.proof.0001
\subsection{Rubrik proof.0001: hipotesis, tipe, dan kuantor}
\label{orig03:rubric:proof:0001}

\paragraph{Kriteria (masing-masing 0--4).}
\begin{enumerate}
  \item Menyatakan objek, ruang, domain, dan kodomain secara konsisten.
  \item Menempatkan semua kuantor dan dependensi parameter pada lingkup yang
  benar, termasuk perbedaan antara pernyataan titik demi titik dan seragam.
  \item Menggunakan tepat hipotesis ketertutupan, konveksitas, keterbatasan,
  atau keterukuran pada langkah yang memerlukannya.
  \item Menutup semua kasus tepi dan memastikan bahwa kesimpulan mempunyai
  tipe logis yang sama dengan klaim semula.
\end{enumerate}

\paragraph{Kegagalan umum.}
Menganggap domain sama dengan seluruh ruang; menukar ``untuk setiap'' dengan
``terdapat''; memakai konstanta yang diam-diam bergantung pada iterasi; membagi
dengan besaran yang belum dibuktikan positif; atau menyatakan limit tanpa
menentukan topologi.

% stable-id: d90.orig.v1.tr03.rubric.proof.0002
\subsection{Rubrik proof.0002: keberadaan dan argumen topologis}
\label{orig03:rubric:proof:0002}

\paragraph{Kriteria (masing-masing 0--4).}
\begin{enumerate}
  \item Mengidentifikasi himpunan sublevel atau barisan peminimum yang tepat
  dan membuktikan keterbatasan atau kekompakannya.
  \item Menggunakan semikontinuitas-bawah, ketertutupan, atau kekompakan dengan
  arah ketaksamaan dan topologi yang benar.
  \item Membedakan dengan tegas keberadaan, ketunggalan, dan kestabilan, serta
  memberi hipotesis tersendiri bagi masing-masing klaim.
  \item Menangani fungsi bernilai-diperluas, domain efektif, dan nilai tak
  hingga tanpa operasi aljabar yang tidak sah.
\end{enumerate}

\paragraph{Kegagalan umum.}
Mengatakan ``kontinu maka minimum ada'' pada himpunan tak kompak; menganggap
konveksitas memberi ketunggalan; menyimpulkan koersivitas hanya dari
keterbatasan-bawah; mengabaikan domain kosong; atau mengambil limit dalam
domain yang belum dibuktikan tertutup.

% stable-id: d90.orig.v1.tr03.rubric.proof.0003
\subsection{Rubrik proof.0003: penurunan, teleskop, dan laju}
\label{orig03:rubric:proof:0003}

\paragraph{Kriteria (masing-masing 0--4).}
\begin{enumerate}
  \item Menurunkan ketaksamaan satu-langkah dari asumsi yang dinyatakan, dengan
  tanda, faktor, dan norma yang benar.
  \item Menentukan besaran potensial nonnegatif dan meneleskopkan indeks tanpa
  kehilangan suku batas.
  \item Memilih langkah atau parameter dalam rentang yang benar dan melacak
  semua konstanta hingga batas akhir.
  \item Membedakan laju nilai fungsi, jarak iterasi, dan residu, termasuk apakah
  klaim berlaku untuk iterasi terakhir, terbaik, atau rata-rata.
\end{enumerate}

\paragraph{Kegagalan umum.}
Menghapus suku bertanda tak diketahui; kesalahan indeks awal/akhir; memilih
langkah setelah melihat kuantitas acak masa depan; mengubah batas
$O(k^{-1/2})$ menjadi $O(k^{-1})$ dengan menguadratkan sepihak; atau
menyembunyikan konstanta yang bergantung pada dimensi.

% stable-id: d90.orig.v1.tr03.rubric.proof.0004
\subsection{Rubrik proof.0004: dualitas, kualifikasi, dan KKT}
\label{orig03:rubric:proof:0004}

\paragraph{Kriteria (masing-masing 0--4).}
\begin{enumerate}
  \item Membentuk Lagrangian dengan tanda, ruang dual, dan domain yang benar.
  \item Menyatakan kondisi kualifikasi yang benar-benar cukup sebelum
  menggunakan dualitas kuat atau keberadaan pengali.
  \item Menurunkan kelayakan primal, kelayakan dual, stasioneritas, dan
  kelonggaran komplementer tanpa melewatkan inklusi subdiferensial.
  \item Membuktikan kebutuhan atau kecukupan kondisi KKT sesuai kelas masalah,
  serta menghubungkannya dengan titik pelana atau nilai dual.
\end{enumerate}

\paragraph{Kegagalan umum.}
Menyatakan KKT tanpa kualifikasi; memberi pengali bertanda salah; mengganti
inklusi dengan persamaan gradien saat fungsi tak mulus; menyamakan dualitas
lemah dengan kuat; atau memakai kelonggaran komplementer untuk menyimpulkan
kelayakan yang belum dibuktikan.

% stable-id: d90.orig.v1.tr03.rubric.proof.0005
\subsection{Rubrik proof.0005: filtrasi dan argumen stokastik}
\label{orig03:rubric:proof:0005}

\paragraph{Kriteria (masing-masing 0--4).}
\begin{enumerate}
  \item Mendefinisikan filtrasi dan menyatakan dengan benar besaran yang
  terukur sebelum pengambilan sampel berikutnya.
  \item Menghitung ekspektasi bersyarat, ketakbiasan, dan bias tanpa
  mengasumsikan independensi yang tidak tersedia.
  \item Membedakan batas varians dari batas momen kedua dan menggunakan hukum
  varians total atau Jensen secara sah.
  \item Menukar ekspektasi, penjumlahan, limit, atau subdiferensial hanya jika
  syarat integrabilitas dan teorema yang dibutuhkan telah dipenuhi.
\end{enumerate}

\paragraph{Kegagalan umum.}
Mengondisikan pada informasi yang sudah memuat sampel baru; menyebut galat
sebagai selisih martingal padahal mean bersyaratnya tidak nol; mengganti
$\mathbb E\lVert G\rVert^2$ dengan varians; memakai sampel tanpa pengembalian
seolah-olah independen; atau mengambil ekspektasi dari kuantitas tak
terintegralkan.

% stable-id: d90.orig.v1.tr03.rubric.proof.0006
\subsection{Rubrik proof.0006: operator monoton dan resolven}
\label{orig03:rubric:proof:0006}

\paragraph{Kriteria (masing-masing 0--4).}
\begin{enumerate}
  \item Menyatakan domain operator bernilai-himpunan dan membedakan monotoni,
  monotoni maksimal, monotoni kuat, serta kokersivitas.
  \item Mengubah inklusi menjadi persamaan resolven atau titik tetap dengan
  urutan komposisi dan parameter yang benar.
  \item Menggunakan firm non-ekspansivitas, keaveragan, atau kontraksi dengan
  konstanta yang dibuktikan, bukan diasumsikan.
  \item Menghubungkan residu atau titik klaster dengan nol operator melalui
  ketertutupan graf dan argumen konvergensi yang lengkap.
\end{enumerate}

\paragraph{Kegagalan umum.}
Menganggap semua operator monoton maksimal; memperlakukan resolven sebagai
invers biasa; menukar $R_AR_B$ dengan $R_BR_A$; mengabaikan domain pada aturan
jumlah; atau menyimpulkan konvergensi iterasi hanya dari residu yang menuju nol
tanpa kontrol titik klaster.

% stable-id: d90.orig.v1.tr03.rubric.proof.0007
\subsection{Rubrik proof.0007: contoh penyangkal dan sertifikat komputasional}
\label{orig03:rubric:proof:0007}

\paragraph{Kriteria (masing-masing 0--4).}
\begin{enumerate}
  \item Menargetkan tepat satu hipotesis atau kesimpulan dan memastikan contoh
  memenuhi semua asumsi lain yang dipertahankan.
  \item Menghitung objek kunci secara eksak atau dengan batas galat yang
  tersertifikasi, termasuk kasus tepi yang relevan.
  \item Memisahkan bukti matematis dari ilustrasi numerik serta menyatakan
  presisi, toleransi, dan residu bila komputasi dipakai.
  \item Menjelaskan secara eksplisit bagaimana hasil contoh membatalkan klaim
  universal atau bagaimana sertifikat memverifikasi klaim yang terbatas.
\end{enumerate}

\paragraph{Kegagalan umum.}
Contoh melanggar asumsi lain; satu plot diperlakukan sebagai bukti universal;
residu kecil disamakan dengan galat solusi tanpa batas kondisi; pembulatan
mengubah tanda; atau program menguji hanya satu parameter sementara klaimnya
untuk semua parameter.


% segment-local-id: OR03-S010
% segment-id: d90.orig.v1.tr03.seg0010
% license: CC BY-SA 4.0
% provenance: Materi asesmen asli berbahasa Indonesia, disusun secara independen.

\section{Ujian tengah kumulatif}
\label{orig03:midterm}

% stable-id: d90.orig.v1.tr03.midterm.exercise.0001
\begin{exercise}[Keberadaan dan ketunggalan model kuadrat--$\ell^1$]
\label{orig03:midterm:exercise:0001}
Untuk $A\in\mathbb R^{m\times n}$, $b\in\mathbb R^m$, dan $\lambda>0$,
pertimbangkan
\[
 F(x)=\lVert Ax-b\rVert_2^2+\lambda\lVert x\rVert_1.
\]
Buktikan bahwa $F$ selalu mempunyai peminimum. Apakah peminimum itu selalu
tunggal? Berikan satu kondisi sederhana pada $A$ yang menjamin ketunggalan.
\end{exercise}

% stable-id: d90.orig.v1.tr03.midterm.hint.0001
\paragraph{Petunjuk bertahap.}
\label{orig03:midterm:hint:0001}
\textbf{Tahap 1.} Bandingkan $F(x)$ dengan $\lambda\lVert x\rVert_1$ ketika
$\lVert x\rVert_2\to\infty$.
\textbf{Tahap 2.} Untuk ketunggalan, cari kondisi yang membuat suku kuadrat
konveks ketat; untuk menyangkal ketunggalan umum, izinkan $A$ memiliki
kernel.

% stable-id: d90.orig.v1.tr03.midterm.answer.0001
\paragraph{Jawaban singkat.}
\label{orig03:midterm:answer:0001}
$F$ kontinu dan koersif karena
$F(x)\ge\lambda\lVert x\rVert_1$, sehingga minimum dicapai. Ketunggalan tidak
selalu berlaku; misalnya $A=0$ dan $b=0$ sebenarnya memberi minimum tunggal,
tetapi matriks dengan kolom identik dapat menghasilkan banyak representasi
$\ell^1$ dengan nilai sama. Kondisi sederhana yang cukup adalah $A$ berperingkat
kolom penuh.

% stable-id: d90.orig.v1.tr03.midterm.solution.0001
\paragraph{Solusi lengkap.}
\label{orig03:midterm:solution:0001}
Kedua suku $F$ kontinu dan tak negatif. Karena
$\lVert x\rVert_1\ge\lVert x\rVert_2$,
\[
 F(x)\ge\lambda\lVert x\rVert_2\longrightarrow\infty
 \quad\text{ketika }\lVert x\rVert_2\to\infty.
\]
Jadi semua himpunan sublevel tak kosong yang relevan kompak, dan teorema
Weierstrass memberi keberadaan peminimum.

Ketunggalan dapat gagal. Sebagai contoh, ambil
$A=(1\;1)$, $b=1$, dan $0<\lambda<2$. Untuk $x_1,x_2\ge0$ dengan
$x_1+x_2=s$, objektif hanya bergantung pada $s$:
$(s-1)^2+\lambda s$. Nilai minimumnya dicapai pada
$s=1-\lambda/2>0$, dan setiap pembagian nonnegatif dari $s$ memberi
peminimum berbeda. Jika $A$ berperingkat kolom penuh, maka
$x\mapsto\lVert Ax-b\rVert^2$ konveks ketat. Jumlah fungsi konveks ketat
dengan fungsi konveks tetap konveks ketat, sehingga peminimum yang sudah
terbukti ada harus tunggal.

% stable-id: d90.orig.v1.tr03.midterm.exercise.0002
\begin{exercise}[Konjugat kuadrat bergeser]
\label{orig03:midterm:exercise:0002}
Untuk $\mu>0$ dan $c\in\mathbb R^n$, definisikan
\[
 f(x)=\frac\mu2\lVert x\rVert^2+\langle c,x\rangle.
\]
Hitung $f^*$ secara langsung dari definisi, verifikasi $f^{**}=f$, dan
tentukan konstanta Lipschitz terkecil dari $\nabla f^*$.
\end{exercise}

% stable-id: d90.orig.v1.tr03.midterm.hint.0002
\paragraph{Petunjuk bertahap.}
\label{orig03:midterm:hint:0002}
\textbf{Tahap 1.} Maksimalkan kuadrat konkaf dalam definisi
$f^*(y)=\sup_x\{\langle y,x\rangle-f(x)\}$.
\textbf{Tahap 2.} Terapkan perhitungan yang sama pada $f^*$ atau substitusikan
gradien inversnya.

% stable-id: d90.orig.v1.tr03.midterm.answer.0002
\paragraph{Jawaban singkat.}
\label{orig03:midterm:answer:0002}
\[
 f^*(y)=\frac1{2\mu}\lVert y-c\rVert^2,
 \qquad
 \nabla f^*(y)=\frac{y-c}{\mu}.
\]
Karena itu $f^{**}=f$, dan konstanta Lipschitz terkecil gradien $f^*$ adalah
$1/\mu$.

% stable-id: d90.orig.v1.tr03.midterm.solution.0002
\paragraph{Solusi lengkap.}
\label{orig03:midterm:solution:0002}
Definisi konjugat memberi
\[
 f^*(y)=\sup_x\left\{\langle y-c,x\rangle-\frac\mu2\lVert x\rVert^2\right\}.
\]
Titik stasionernya $x=(y-c)/\mu$, dan substitusi menghasilkan
$\lVert y-c\rVert^2/(2\mu)$. Selanjutnya,
\[
 f^{**}(x)=\sup_y\left\{\langle x,y\rangle
 -\frac1{2\mu}\lVert y-c\rVert^2\right\}.
\]
Peminimum negatif kuadrat, atau kondisi stasioner, memberi
$y=c+\mu x$. Nilainya
$\langle c,x\rangle+\mu\lVert x\rVert^2/2=f(x)$. Akhirnya,
\[
 \lVert\nabla f^*(y)-\nabla f^*(z)\rVert
 =\frac1\mu\lVert y-z\rVert.
\]
Kesetaraan berlaku untuk setiap pasangan berbeda, jadi $1/\mu$ bukan hanya
batas yang cukup, melainkan konstanta terkecil.

% stable-id: d90.orig.v1.tr03.midterm.exercise.0003
\begin{exercise}[Ketaksamaan satu langkah gradien proksimal]
\label{orig03:midterm:exercise:0003}
Misalkan $f$ konveks dengan gradien $L$-Lipschitz dan $r$ konveks proper
semikontinu-bawah. Untuk $0<\eta\le1/L$, definisikan
\[
 x^+=\operatorname{prox}_{\eta r}(x-\eta\nabla f(x)),
 \qquad F=f+r.
\]
Buktikan bahwa untuk setiap $u\in\operatorname{dom}r$,
\[
 F(x^+)-F(u)
 \le\frac{\lVert x-u\rVert^2-\lVert x^+-u\rVert^2}{2\eta}.
\]
Kemudian turunkan batas $O(1/k)$ untuk nilai fungsi iterasi gradien proksimal
bila peminimum $x^\star$ ada.
\end{exercise}

% stable-id: d90.orig.v1.tr03.midterm.hint.0003
\paragraph{Petunjuk bertahap.}
\label{orig03:midterm:hint:0003}
\textbf{Tahap 1.} Gabungkan lemma penurunan untuk $f$ dengan inklusi
optimalitas proksimal dan ketaksamaan subgradien untuk $r$.
\textbf{Tahap 2.} Gunakan identitas tiga norma, tetapkan $u=x^\star$, lalu
jumlahkan dan manfaatkan monotoni nilai objektif.

% stable-id: d90.orig.v1.tr03.midterm.answer.0003
\paragraph{Jawaban singkat.}
\label{orig03:midterm:answer:0003}
Suku tambahan yang muncul adalah
$-(1/(2\eta)-L/2)\lVert x^+-x\rVert^2\le0$, sehingga ketaksamaan yang diminta
berlaku. Untuk iterasi $x_{j+1}$ yang sama,
\[
 F(x_k)-F(x^\star)
 \le\frac{\lVert x_0-x^\star\rVert^2}{2\eta k}.
\]

% stable-id: d90.orig.v1.tr03.midterm.solution.0003
\paragraph{Solusi lengkap.}
\label{orig03:midterm:solution:0003}
Inklusi proksimal menyediakan $s^+\in\partial r(x^+)$ dengan
\[
 s^+=-\nabla f(x)-\frac1\eta(x^+-x).
\]
Lemma penurunan, konveksitas $f$, dan ketaksamaan subgradien memberi
\begin{align*}
 F(x^+)-F(u)
 &\le \langle\nabla f(x),x^+-u\rangle
      +\frac L2\lVert x^+-x\rVert^2
      +\langle s^+,x^+-u\rangle\\
 &=-\frac1\eta\langle x^+-x,x^+-u\rangle
      +\frac L2\lVert x^+-x\rVert^2\\
 &=\frac{\lVert x-u\rVert^2-\lVert x^+-u\rVert^2}{2\eta}
   -\left(\frac1{2\eta}-\frac L2\right)\lVert x^+-x\rVert^2.
\end{align*}
Koefisien terakhir nonnegatif karena $\eta\le1/L$, sehingga suku negatif itu
dapat dibuang. Dengan $u=x^\star$, penjumlahan untuk $j=0,\ldots,k-1$
meneleskopkan jarak dan memberi
\[
 \sum_{j=0}^{k-1}\bigl(F(x_{j+1})-F(x^\star)\bigr)
 \le\frac{\lVert x_0-x^\star\rVert^2}{2\eta}.
\]
Ketaksamaan yang sama dengan $u=x_j$ menunjukkan nilai $F(x_j)$ tidak naik.
Karena itu suku terakhir tidak melebihi rata-rata $k$ suku, yang menghasilkan
batas pada jawaban singkat.

% stable-id: d90.orig.v1.tr03.midterm.exercise.0004
\begin{exercise}[Newton dan penerimaan langkah penuh Armijo]
\label{orig03:midterm:exercise:0004}
Misalkan
\[
 f(x)=\tfrac12x^\top Qx-b^\top x,
 \qquad Q\succ0.
\]
Dari titik sembarang $x$, tentukan arah Newton $p$, buktikan bahwa $x+p$
adalah peminimum eksak, dan tunjukkan bahwa uji Armijo
\[
 f(x+p)\le f(x)+c\nabla f(x)^\top p
\]
menerima langkah penuh untuk setiap $0<c\le1/2$.
\end{exercise}

% stable-id: d90.orig.v1.tr03.midterm.hint.0004
\paragraph{Petunjuk bertahap.}
\label{orig03:midterm:hint:0004}
\textbf{Tahap 1.} Selesaikan $Qp=-\nabla f(x)$ dan hitung gradien di $x+p$.
\textbf{Tahap 2.} Gunakan ekspansi kuadrat eksak sepanjang $p$ serta
$\nabla f(x)^\top p=-p^\top Qp$.

% stable-id: d90.orig.v1.tr03.midterm.answer.0004
\paragraph{Jawaban singkat.}
\label{orig03:midterm:answer:0004}
$p=Q^{-1}(b-Qx)$ dan $x+p=Q^{-1}b=x^\star$. Selain itu,
\[
 f(x+p)=f(x)-\tfrac12p^\top Qp,
 \qquad \nabla f(x)^\top p=-p^\top Qp,
\]
sehingga Armijo berlaku tepat bila $c\le1/2$ (atau trivial bila $p=0$).

% stable-id: d90.orig.v1.tr03.midterm.solution.0004
\paragraph{Solusi lengkap.}
\label{orig03:midterm:solution:0004}
Gradiennya $Qx-b$ dan Hessiannya $Q$. Persamaan Newton
$Qp=-(Qx-b)$ memberi $p=Q^{-1}b-x$, sehingga titik baru adalah
$Q^{-1}b$. Karena $Q$ definit positif, titik stasioner ini peminimum tunggal.
Ekspansi eksak memberi
\[
 f(x+p)=f(x)+\nabla f(x)^\top p+\tfrac12p^\top Qp
 =f(x)-\tfrac12p^\top Qp.
\]
Ruas kanan uji Armijo adalah
$f(x)-c\,p^\top Qp$. Untuk $p\ne0$, definit positifnya $Q$ membuat uji
ekuivalen dengan $1/2\ge c$. Jika $p=0$, titik sudah optimal dan kedua ruas
sama untuk semua $c$.

% stable-id: d90.orig.v1.tr03.midterm.exercise.0005
\begin{exercise}[KKT untuk proyeksi pada bola]
\label{orig03:midterm:exercise:0005}
Ambil $R>0$ dan $c\in\mathbb R^n$ dengan $\lVert c\rVert>R$. Selesaikan
\[
 \min_x\ \tfrac12\lVert x-c\rVert^2
 \quad\text{dengan syarat}\quad
 h(x)=\tfrac12(\lVert x\rVert^2-R^2)\le0.
\]
Verifikasi kondisi Slater, tulis seluruh kondisi KKT, dan hitung pasangan
primal--dual optimal secara eksplisit.
\end{exercise}

% stable-id: d90.orig.v1.tr03.midterm.hint.0005
\paragraph{Petunjuk bertahap.}
\label{orig03:midterm:hint:0005}
\textbf{Tahap 1.} Titik nol adalah titik Slater; karena $c$ di luar bola,
kendala pada solusi harus aktif.
\textbf{Tahap 2.} Stasioneritas memberi $(1+\lambda)x=c$.

% stable-id: d90.orig.v1.tr03.midterm.answer.0005
\paragraph{Jawaban singkat.}
\label{orig03:midterm:answer:0005}
KKT terdiri dari
\[
 x-c+\lambda x=0,quad h(x)\le0,quad\lambda\ge0,quad\lambda h(x)=0.
\]
Solusinya
\[
 x^\star=R\frac{c}{\lVert c\rVert},
 \qquad \lambda^\star=\frac{\lVert c\rVert}{R}-1.
\]

% stable-id: d90.orig.v1.tr03.midterm.solution.0005
\paragraph{Solusi lengkap.}
\label{orig03:midterm:solution:0005}
Karena $h(0)=-R^2/2<0$, Slater berlaku. Masalah konveks dan objektifnya
konveks ketat, sehingga KKT perlu dan cukup serta solusi primal tunggal.
Lagrangiannya
\[
 L(x,\lambda)=\tfrac12\lVert x-c\rVert^2
 +\frac\lambda2(\lVert x\rVert^2-R^2).
\]
Stasioneritas dan tiga kondisi lainnya tercantum pada jawaban singkat. Solusi
tanpa kendala $x=c$ tidak layak, maka kendala aktif dan
$\lVert x\rVert=R$. Dari $(1+\lambda)x=c$ diperoleh bahwa $x$ searah dengan
$c$. Menormakan persamaan itu menghasilkan
$1+\lambda=\lVert c\rVert/R$, sehingga pasangan yang dinyatakan mengikuti.
Pengalinya positif, maka semua kondisi KKT terpenuhi.

% stable-id: d90.orig.v1.tr03.midterm.exercise.0006
\begin{exercise}[Ketika kualifikasi gagal dan pengali tidak ada]
\label{orig03:midterm:exercise:0006}
Pertimbangkan masalah konveks skalar
\[
 \min_{x\in\mathbb R} x
 \quad\text{dengan syarat}\quad x^2\le0.
\]
Tentukan solusi, tunjukkan bahwa Slater gagal, dan buktikan bahwa tidak ada
pengali $\lambda\ge0$ yang memenuhi kondisi KKT pada solusi. Jelaskan mengapa
contoh ini tidak menyangkal teorema KKT dengan kualifikasi.
\end{exercise}

% stable-id: d90.orig.v1.tr03.midterm.hint.0006
\paragraph{Petunjuk bertahap.}
\label{orig03:midterm:hint:0006}
\textbf{Tahap 1.} Tentukan seluruh himpunan layak dari ketaksamaan kuadrat.
\textbf{Tahap 2.} Diferensiasikan $L(x,\lambda)=x+\lambda x^2$ di titik layak.

% stable-id: d90.orig.v1.tr03.midterm.answer.0006
\paragraph{Jawaban singkat.}
\label{orig03:midterm:answer:0006}
Satu-satunya titik layak dan optimal ialah $x^\star=0$. Tidak ada titik dengan
$x^2<0$, jadi Slater gagal. Stasioneritas KKT menuntut
$1+2\lambda x^\star=0$, yakni $1=0$, mustahil untuk setiap $\lambda$.

% stable-id: d90.orig.v1.tr03.midterm.solution.0006
\paragraph{Solusi lengkap.}
\label{orig03:midterm:solution:0006}
Karena kuadrat selalu nonnegatif, syarat $x^2\le0$ ekuivalen dengan $x=0$.
Maka $x^\star=0$ adalah satu-satunya titik layak dan otomatis peminimum.
Slater memerlukan titik dengan $x^2<0$, yang tidak ada. Untuk pengali
$\lambda\ge0$, Lagrangian adalah $L(x,\lambda)=x+\lambda x^2$ dan kondisi
stasioneritas pada $x^\star$ ialah
\[
 0=\frac{\partial L}{\partial x}(0,\lambda)=1+2\lambda\cdot0=1,
\]
sebuah kontradiksi. Teorema KKT yang menjamin keberadaan pengali memakai
kondisi kualifikasi seperti Slater. Hipotesis itu gagal di sini, sehingga
ketiadaan pengali justru menunjukkan mengapa kualifikasi tidak boleh
dihilangkan.

% segment-local-id: OR03-S011
% segment-id: d90.orig.v1.tr03.seg0011
% license: CC BY-SA 4.0
% provenance: Materi asesmen asli berbahasa Indonesia, disusun secara independen.

\section{Ujian akhir kumulatif}
\label{orig03:final}

% stable-id: d90.orig.v1.tr03.final.exercise.0001
\begin{exercise}[Eksistensi komposit dengan regularisasi kuadratik]
\label{orig03:final:exercise:0001}
Misalkan $h:\mathbb R^m\to\mathbb R$ konveks, kontinu, dan tak negatif;
$r:\mathbb R^n\to(-\infty,+\infty]$ proper, semikontinu-bawah, konveks, dan
tak negatif; $A:\mathbb R^n\to\mathbb R^m$ linear; serta $\lambda>0$.
Buktikan bahwa
\[
 F(x)=h(Ax)+r(x)+\frac\lambda2\lVert x\rVert^2
\]
mempunyai tepat satu peminimum. Nyatakan bagian argumen yang membuktikan
keberadaan dan bagian yang membuktikan ketunggalan.
\end{exercise}

% stable-id: d90.orig.v1.tr03.final.hint.0001
\paragraph{Petunjuk bertahap.}
\label{orig03:final:hint:0001}
\textbf{Tahap 1.} Gunakan ketaknegatifan dua suku pertama untuk membatasi
himpunan sublevel.
\textbf{Tahap 2.} Periksa apa yang terjadi pada modulus konveksitas kuat ketika
fungsi konveks ditambahkan pada kuadrat $\lambda\lVert x\rVert^2/2$.

% stable-id: d90.orig.v1.tr03.final.answer.0001
\paragraph{Jawaban singkat.}
\label{orig03:final:answer:0001}
$F$ proper, semikontinu-bawah, dan
$F(x)\ge\lambda\lVert x\rVert^2/2$, sehingga ia koersif dan mencapai minimum.
Selain itu, $F$ bersifat $\lambda$-konveks kuat, sehingga peminimumnya tunggal.

% stable-id: d90.orig.v1.tr03.final.solution.0001
\paragraph{Solusi lengkap.}
\label{orig03:final:solution:0001}
Karena $r$ proper, terdapat $x_0$ dengan $r(x_0)<\infty$; fungsi $h$ bernilai
hingga, jadi $F(x_0)<\infty$ dan $F$ proper. Komposisi $h\circ A$ kontinu dan
konveks, sementara $r$ semikontinu-bawah. Maka $F$ semikontinu-bawah. Dari
ketaknegatifan,
\[
 F(x)\ge\frac\lambda2\lVert x\rVert^2\longrightarrow+\infty
 \quad\text{ketika }\lVert x\rVert\to\infty.
\]
Jadi sebuah barisan peminimum berada dalam himpunan terbatas; dalam dimensi
hingga kita dapat mengambil subsekuens konvergen, dan semikontinuitas-bawah
menunjukkan bahwa limitnya mencapai infimum. Ini membuktikan keberadaan.

Fungsi kuadrat adalah $\lambda$-konveks kuat. Penambahan dua fungsi konveks
mempertahankan modulus itu, jadi $F$ juga $\lambda$-konveks kuat. Jika ada dua
peminimum berbeda, ketaksamaan konveksitas kuat pada titik tengah akan memberi
nilai yang lebih kecil secara ketat daripada nilai minimum, kontradiksi. Maka
peminimum tunggal.

% stable-id: d90.orig.v1.tr03.final.exercise.0002
\begin{exercise}[Pencarian balik untuk model proksimal]
\label{orig03:final:exercise:0002}
Misalkan $f$ konveks dengan gradien $L$-Lipschitz, $r$ proper
semikontinu-bawah dan konveks, serta $x\in\operatorname{dom}r$. Untuk $M>0$,
definisikan
\[
 x_M=\arg\min_y\left\{
 r(y)+f(x)+\langle\nabla f(x),y-x\rangle
 +\frac M2\lVert y-x\rVert^2\right\}.
\]
Algoritme menggandakan $M$ sampai uji majorisasi
\[
 f(x_M)\le f(x)+\langle\nabla f(x),x_M-x\rangle
 +\frac M2\lVert x_M-x\rVert^2                         \tag{A}
\]
terpenuhi. Buktikan bahwa pencarian berhenti secara hingga, lalu buktikan
sertifikat penurunan
\[
 (f+r)(x_M)\le(f+r)(x)-\frac M2\lVert x_M-x\rVert^2.
\]
Nyatakan residu gradien proksimal yang dapat dihitung dari langkah ini.
\end{exercise}

% stable-id: d90.orig.v1.tr03.final.hint.0002
\paragraph{Petunjuk bertahap.}
\label{orig03:final:hint:0002}
\textbf{Tahap 1.} Lemma penurunan menjamin (A) untuk semua $M\ge L$.
\textbf{Tahap 2.} Model yang diminimalkan bersifat $M$-konveks kuat dan
nilainya di $y=x$ sama dengan $(f+r)(x)$.

% stable-id: d90.orig.v1.tr03.final.answer.0002
\paragraph{Jawaban singkat.}
\label{orig03:final:answer:0002}
Penggandaan akhirnya menghasilkan $M\ge L$, sehingga (A) berlaku. Konveksitas
kuat model memberi nilai model di $x_M$ paling besar
$(f+r)(x)-M\lVert x_M-x\rVert^2/2$; (A) memindahkan batas itu ke objektif
sebenarnya. Residu yang lazim adalah
$G_M(x)=M(x-x_M)$.

% stable-id: d90.orig.v1.tr03.final.solution.0002
\paragraph{Solusi lengkap.}
\label{orig03:final:solution:0002}
Lemma penurunan menyatakan untuk semua $y$,
\[
 f(y)\le f(x)+\langle\nabla f(x),y-x\rangle
 +\frac L2\lVert y-x\rVert^2.
\]
Jadi begitu urutan penggandaan mencapai $M\ge L$, uji (A) pasti diterima;
jumlah penggandaan hingga.

Tuliskan model di dalam kurung sebagai $Q_M(y)$. Ia $M$-konveks kuat dan
$x_M$ adalah peminimumnya. Dengan memakai $0\in\partial Q_M(x_M)$ dalam
ketaksamaan konveksitas kuat,
\[
 Q_M(x)\ge Q_M(x_M)+\frac M2\lVert x-x_M\rVert^2.
\]
Karena $Q_M(x)=(f+r)(x)$ dan uji (A) memberi
$(f+r)(x_M)\le Q_M(x_M)$, kedua ketaksamaan menghasilkan sertifikat yang
diminta. Vektor $G_M(x)=M(x-x_M)$ dapat dihitung dari langkah. Ia nol tepat
bila $x=x_M$, yang melalui kondisi optimalitas ekuivalen dengan
$0\in\nabla f(x)+\partial r(x)$.

% stable-id: d90.orig.v1.tr03.final.exercise.0003
\begin{exercise}[KKT untuk regularisasi analisis dan kendala norma]
\label{orig03:final:exercise:0003}
Untuk $\tau>0$ dan $R>0$, pertimbangkan
\[
 \min_x\ \frac12\lVert Ax-b\rVert^2+\tau\lVert x\rVert_1
 \quad\text{dengan syarat}\quad \lVert Bx-d\rVert_2\le R.
\]
Anggap terdapat $\bar x$ dengan $\lVert B\bar x-d\rVert_2<R$. Turunkan sistem
KKT lengkap. Tuliskan subgradien norma Euclidean baik ketika
$Bx-d\ne0$ maupun ketika $Bx-d=0$, dan jelaskan mengapa sistem KKT cukup untuk
optimalitas global.
\end{exercise}

% stable-id: d90.orig.v1.tr03.final.hint.0003
\paragraph{Petunjuk bertahap.}
\label{orig03:final:hint:0003}
\textbf{Tahap 1.} Gunakan pengali $\nu\ge0$ untuk fungsi kendala
$q(x)=\lVert Bx-d\rVert_2-R$.
\textbf{Tahap 2.} Terapkan aturan rantai subdiferensial pada norma dan pisahkan
kasus argumen nol.

% stable-id: d90.orig.v1.tr03.final.answer.0003
\paragraph{Jawaban singkat.}
\label{orig03:final:answer:0003}
Terdapat $s\in\partial\lVert x\rVert_1$, $q\in\partial\lVert Bx-d\rVert_2$,
dan $\nu\ge0$ sehingga
\[
 0=A^\top(Ax-b)+\tau s+\nu B^\top q,
\]
\[
 \lVert Bx-d\rVert_2\le R,qquad
 \nu(\lVert Bx-d\rVert_2-R)=0.
\]
Jika $z\ne0$, $\partial\lVert z\rVert_2=\{z/\lVert z\rVert_2\}$; jika
$z=0$, subdiferensialnya adalah bola satuan. Slater membuat KKT perlu dan,
karena masalah konveks, juga cukup.

% stable-id: d90.orig.v1.tr03.final.solution.0003
\paragraph{Solusi lengkap.}
\label{orig03:final:solution:0003}
Lagrangian adalah
\[
 \mathcal L(x,\nu)=\tfrac12\lVert Ax-b\rVert^2
 +\tau\lVert x\rVert_1+\nu(\lVert Bx-d\rVert_2-R).
\]
Kondisi Slater yang diberikan menjamin keberadaan pengali optimal dan aturan
KKT tanpa celah dual. Stasioneritas adalah inklusi
\[
 0\in A^\top(Ax-b)+\tau\partial\lVert x\rVert_1
 +\nu B^\top\partial\lVert Bx-d\rVert_2.
\]
Memilih $s$ dan $q$ dari kedua subdiferensial mengubahnya menjadi persamaan
pada jawaban singkat. Secara koordinat,
$s_i=\operatorname{sign}(x_i)$ jika $x_i\ne0$ dan $s_i\in[-1,1]$ jika
$x_i=0$. Untuk norma Euclidean, $q=z/\lVert z\rVert$ saat $z=Bx-d\ne0$,
sedangkan sembarang $q$ dengan $\lVert q\rVert\le1$ sah saat $z=0$.

Sistem dilengkapi oleh kelayakan primal, kelayakan dual $\nu\ge0$, dan
kelonggaran komplementer. Objektif dan fungsi kendala konveks; maka setiap
pasangan yang memenuhi sistem membentuk sertifikat titik pelana Lagrangian
dan memberi peminimum global.

% stable-id: d90.orig.v1.tr03.final.exercise.0004
\begin{exercise}[Distribusi kepentingan dengan momen kedua minimum]
\label{orig03:final:exercise:0004}
Misalkan $g_1,\ldots,g_N$ adalah vektor tak nol yang diketahui. Untuk peluang
$p_i>0$, $\sum_i p_i=1$, ambil $I\sim p$ dan
$G=g_I/(Np_I)$. Di antara semua distribusi $p$, tentukan distribusi yang
meminimalkan $\mathbb E\lVert G\rVert^2$, hitung nilai minimumnya, dan tulis
varians minimum terhadap mean $\bar g=N^{-1}\sum_i g_i$.
\end{exercise}

% stable-id: d90.orig.v1.tr03.final.hint.0004
\paragraph{Petunjuk bertahap.}
\label{orig03:final:hint:0004}
\textbf{Tahap 1.} Masalah skalar yang harus diminimalkan adalah
$\sum_i\lVert g_i\rVert^2/p_i$ pada simpleks.
\textbf{Tahap 2.} Gunakan pengali Lagrange atau Cauchy--Schwarz, lalu kurangi
$\lVert\bar g\rVert^2$ dari momen kedua.

% stable-id: d90.orig.v1.tr03.final.answer.0004
\paragraph{Jawaban singkat.}
\label{orig03:final:answer:0004}
\[
 p_i^\star=\frac{\lVert g_i\rVert}{\sum_j\lVert g_j\rVert},
 \qquad
 \min_p\mathbb E\lVert G\rVert^2
 =\frac1{N^2}\left(\sum_i\lVert g_i\rVert\right)^2.
\]
Varians minimumnya adalah nilai terakhir dikurangi $\lVert\bar g\rVert^2$.

% stable-id: d90.orig.v1.tr03.final.solution.0004
\paragraph{Solusi lengkap.}
\label{orig03:final:solution:0004}
Ketakbiasan mengikuti dari
$\mathbb E G=\sum_i p_i g_i/(Np_i)=\bar g$. Momen keduanya ialah
\[
 \mathbb E\lVert G\rVert^2
 =\frac1{N^2}\sum_i\frac{\lVert g_i\rVert^2}{p_i}.
\]
Cauchy--Schwarz memberi
\[
 \left(\sum_i\lVert g_i\rVert\right)^2
 =\left(\sum_i\frac{\lVert g_i\rVert}{\sqrt{p_i}}\sqrt{p_i}\right)^2
 \le\left(\sum_i\frac{\lVert g_i\rVert^2}{p_i}\right)
      \left(\sum_i p_i\right).
\]
Kesetaraan terjadi tepat ketika
$\lVert g_i\rVert/\sqrt{p_i}$ sebanding dengan $\sqrt{p_i}$, yakni
$p_i\propto\lVert g_i\rVert$. Karena semua vektor tak nol, distribusi itu
strictly positif dan layak. Substitusi memberi nilai minimum. Identitas
$\mathbb E\lVert G-\bar g\rVert^2=
\mathbb E\lVert G\rVert^2-\lVert\bar g\rVert^2$ memberi varians yang diminta.

% stable-id: d90.orig.v1.tr03.final.exercise.0005
\begin{exercise}[Laju rata-rata descent cermin stokastik]
\label{orig03:final:exercise:0005}
Misalkan $C$ konveks, $h$ terdiferensialkan dan $1$-konveks kuat terhadap
norma $\lVert\cdot\rVert$, serta $f$ konveks pada $C$. Dengan filtrasi
$\mathcal F_k$, anggap
\[
 \mathbb E[G_k\mid\mathcal F_k]=g_k\in\partial f(x_k),
 \qquad
 \mathbb E[\lVert G_k\rVert_*^2\mid\mathcal F_k]\le G^2.
\]
Iterasi diberikan oleh
\[
 x_{k+1}\in\arg\min_{x\in C}
 \{\eta\langle G_k,x\rangle+D_h(x,x_k)\}.
\]
Jika $x^\star$ meminimalkan $f$ dan
$D_h(x^\star,x_0)\le D^2/2$, buktikan untuk
$\bar x_K=K^{-1}\sum_{k=0}^{K-1}x_k$ dan
$\eta=D/(G\sqrt K)$ bahwa
\[
 \mathbb E[f(\bar x_K)-f(x^\star)]\le\frac{DG}{\sqrt K}.
\]
\end{exercise}

% stable-id: d90.orig.v1.tr03.final.hint.0005
\paragraph{Petunjuk bertahap.}
\label{orig03:final:hint:0005}
\textbf{Tahap 1.} Dari ketaksamaan tiga titik dan konveksitas kuat $h$,
turunkan batas satu langkah dengan suku $\eta^2\lVert G_k\rVert_*^2/2$.
\textbf{Tahap 2.} Ambil ekspektasi bersyarat, jumlahkan, gunakan Jensen pada
$f(\bar x_K)$, lalu optimalkan langkah konstan.

% stable-id: d90.orig.v1.tr03.final.answer.0005
\paragraph{Jawaban singkat.}
\label{orig03:final:answer:0005}
Ketaksamaan dasar memberi
\[
 \eta\,\mathbb E[f(x_k)-f(x^\star)]
 \le\mathbb E[D_h(x^\star,x_k)-D_h(x^\star,x_{k+1})]
 +\frac{\eta^2G^2}{2}.
\]
Setelah dijumlahkan dan dibagi $\eta K$, batasnya
$D^2/(2\eta K)+\eta G^2/2$; pilihan langkah yang diberikan membuatnya
$DG/\sqrt K$.

% stable-id: d90.orig.v1.tr03.final.solution.0005
\paragraph{Solusi lengkap.}
\label{orig03:final:solution:0005}
Optimalitas langkah cermin dan identitas tiga titik menghasilkan, untuk
$u\in C$,
\[
 \eta\langle G_k,x_{k+1}-u\rangle
 \le D_h(u,x_k)-D_h(u,x_{k+1})-D_h(x_{k+1},x_k).
\]
Tambahkan $\eta\langle G_k,x_k-x_{k+1}\rangle$ pada kedua ruas. Karena
$D_h(x_{k+1},x_k)\ge\lVert x_{k+1}-x_k\rVert^2/2$, Young memberi
\[
 \eta\langle G_k,x_k-u\rangle
 \le D_h(u,x_k)-D_h(u,x_{k+1})
 +\frac{\eta^2}{2}\lVert G_k\rVert_*^2.
\]
Ambil $u=x^\star$ dan ekspektasi bersyarat. Karena $x_k$ terukur terhadap
$\mathcal F_k$, ketakbiasan memberi
$\mathbb E[\langle G_k,x_k-x^\star\rangle\mid\mathcal F_k]
=\langle g_k,x_k-x^\star\rangle\ge f(x_k)-f(x^\star)$.
Maka ketaksamaan pada jawaban singkat berlaku.

Menjumlahkan untuk $k=0,\ldots,K-1$ meneleskopkan divergensi dan memberi
\[
 \frac1K\sum_{k=0}^{K-1}\mathbb E[f(x_k)-f(x^\star)]
 \le\frac{D^2}{2\eta K}+\frac{\eta G^2}{2}.
\]
Konveksitas dan Jensen membatasi ruas kiri dari bawah oleh
$\mathbb E[f(\bar x_K)-f(x^\star)]$. Substitusi
$\eta=D/(G\sqrt K)$ membuat kedua suku kanan sama dengan
$DG/(2\sqrt K)$, sehingga hasil mengikuti.

% stable-id: d90.orig.v1.tr03.final.exercise.0006
\begin{exercise}[Kontraksi resolven operator monoton kuat]
\label{orig03:final:exercise:0006}
Misalkan $A$ maksimal monoton dan $\mu$-monoton kuat dengan $\mu>0$.
Buktikan bahwa untuk setiap $\gamma>0$,
\[
 \lVert J_{\gamma A}x-J_{\gamma A}y\rVert
 \le\frac1{1+\gamma\mu}\lVert x-y\rVert.
\]
Simpulkan bahwa $A$ mempunyai paling banyak satu nol dan, jika nol
$x^\star$ ada, iterasi $x_{k+1}=J_{\gamma A}x_k$ konvergen linear kepadanya.
\end{exercise}

% stable-id: d90.orig.v1.tr03.final.hint.0006
\paragraph{Petunjuk bertahap.}
\label{orig03:final:hint:0006}
\textbf{Tahap 1.} Tetapkan $p=J_{\gamma A}x$, $q=J_{\gamma A}y$, lalu tulis
anggota $Ap$ dan $Aq$ dari definisi resolven.
\textbf{Tahap 2.} Terapkan monotoni kuat dan Cauchy--Schwarz; kemudian ambil
$y=x^\star$.

% stable-id: d90.orig.v1.tr03.final.answer.0006
\paragraph{Jawaban singkat.}
\label{orig03:final:answer:0006}
Monotoni kuat memberi
\[
 \langle p-q,x-y\rangle\ge(1+\gamma\mu)\lVert p-q\rVert^2,
\]
yang bersama Cauchy--Schwarz menghasilkan kontraksi. Karena
$J_{\gamma A}x^\star=x^\star$,
\[
 \lVert x_k-x^\star\rVert
 \le(1+\gamma\mu)^{-k}\lVert x_0-x^\star\rVert.
\]

% stable-id: d90.orig.v1.tr03.final.solution.0006
\paragraph{Solusi lengkap.}
\label{orig03:final:solution:0006}
Dari definisi resolven,
\[
 \frac{x-p}{\gamma}\in Ap,
 \qquad
 \frac{y-q}{\gamma}\in Aq.
\]
Monotoni kuat memberikan
\[
 \left\langle p-q,\frac{x-p}{\gamma}-\frac{y-q}{\gamma}\right\rangle
 \ge\mu\lVert p-q\rVert^2.
\]
Setelah dikalikan $\gamma$ dan ditata ulang, ini menjadi ketaksamaan pada
jawaban singkat. Jika $p\ne q$, Cauchy--Schwarz dan pembagian oleh
$\lVert p-q\rVert$ memberi batas Lipschitz; bila $p=q$ batasnya trivial.

Jika $u$ dan $v$ keduanya nol $A$, monotoni kuat memberi
$0\ge\mu\lVert u-v\rVert^2$, sehingga $u=v$. Untuk nol yang ada,
$J_{\gamma A}x^\star=x^\star$. Penerapan berulang faktor kontraksi
$q=(1+\gamma\mu)^{-1}<1$ memberi batas linear yang dinyatakan.

% stable-id: d90.orig.v1.tr03.final.exercise.0007
\begin{exercise}[Sinkhorn sebagai maksimisasi blok dual]
\label{orig03:final:exercise:0007}
Misalkan $a_i,b_j>0$ dan total massanya sama. Untuk $\varepsilon>0$, tinjau
\[
 D(f,g)=a^\top f+b^\top g
 -\varepsilon\sum_{i,j}
 \exp\!\left(\frac{f_i+g_j-C_{ij}}{\varepsilon}\right).
\]
Buktikan bahwa, dengan $g$ tetap, pemaksimal unik terhadap setiap $f_i$ adalah
\[
 f_i=\varepsilon\log a_i-\varepsilon
 \log\sum_j\exp\!\left(\frac{g_j-C_{ij}}{\varepsilon}\right),
\]
dan turunkan pembaruan analog untuk $g_j$. Simpulkan bahwa penskalaan Sinkhorn
adalah maksimisasi koordinat blok eksak dan nilai dual tidak menurun pada
setiap setengah langkah. Jelaskan ketidaktunggalan bersama akibat gauge.
\end{exercise}

% stable-id: d90.orig.v1.tr03.final.hint.0007
\paragraph{Petunjuk bertahap.}
\label{orig03:final:hint:0007}
\textbf{Tahap 1.} Nolkan turunan parsial terhadap satu $f_i$ dan periksa tanda
turunan keduanya.
\textbf{Tahap 2.} Ulangi untuk $g_j$ dan amati transformasi
$(f,g)\mapsto(f+c\1,g-c\1)$.

% stable-id: d90.orig.v1.tr03.final.answer.0007
\paragraph{Jawaban singkat.}
\label{orig03:final:answer:0007}
Persamaan stasioner blok $f$ ialah
\[
 a_i=\sum_j\exp((f_i+g_j-C_{ij})/\varepsilon),
\]
yang memberi rumus pada soal; rumus blok $g$ diperoleh dengan menukar baris
dan kolom. Setiap turunan kedua blok negatif, jadi pembaruan adalah maksimum
unik dan tidak menurunkan $D$. Pasangan bersama hanya unik hingga penambahan
konstanta pada $f$ dan pengurangan konstanta yang sama pada $g$.

% stable-id: d90.orig.v1.tr03.final.solution.0007
\paragraph{Solusi lengkap.}
\label{orig03:final:solution:0007}
Untuk $g$ tetap,
\[
 \frac{\partial D}{\partial f_i}
 =a_i-\sum_j\exp((f_i+g_j-C_{ij})/\varepsilon).
\]
Menetapkannya nol dan mengeluarkan faktor $\exp(f_i/\varepsilon)$ memberi
rumus eksplisit pada soal. Turunan kedua adalah
\[
 -\frac1\varepsilon\sum_j
 \exp((f_i+g_j-C_{ij})/\varepsilon)<0,
\]
dan tidak ada suku silang antara $f_i$ yang berlainan ketika $g$ tetap. Jadi
rumus itu adalah pemaksimal blok unik. Secara simetris,
\[
 g_j=\varepsilon\log b_j-\varepsilon
 \log\sum_i\exp\!\left(\frac{f_i-C_{ij}}{\varepsilon}\right).
\]
Mengganti satu blok dengan pemaksimal eksaknya tidak dapat menurunkan nilai
fungsi dual. Dalam variabel skala eksponensial, kedua rumus adalah pembaruan
baris dan kolom Sinkhorn.

Jika $f$ ditambah $c\1$ dan $g$ dikurangi $c\1$, semua jumlah
$f_i+g_j$ tetap. Suku linear juga tetap karena total massa $a$ dan $b$ sama.
Jadi nilai dual dan kopling tidak berubah; inilah kebebasan gauge yang membuat
pemaksimal pasangan tidak tunggal sebelum normalisasi.

% stable-id: d90.orig.v1.tr03.final.exercise.0008
\begin{exercise}[Pemilihan pemisahan untuk model stokastik terstruktur]
\label{orig03:final:exercise:0008}
Pertimbangkan
\[
 \min_{x\in C}\quad
 s(x)+\tau\lVert Wx\rVert_1,
 \qquad
 s(x)=\mathbb E_\xi[\ell_\xi(A_\xi x-b_\xi)],
\]
dengan $C=\{x:\lVert x\rVert_2\le R\}$, $\tau>0$, dan $s$ konveks
terdiferensialkan. Tuliskan inklusi optimalitas primal dan sistem primal--dual
dengan variabel $y$ untuk suku $\ell^1$. Bandingkan kapan gradien proksimal
stokastik langsung layak dipakai dan kapan pemisahan primal--dual lebih
sesuai. Berikan sepasang residu KKT dan asumsi orakel minimum yang diperlukan
untuk klaim konvergensi stokastik.
\end{exercise}

% stable-id: d90.orig.v1.tr03.final.hint.0008
\paragraph{Petunjuk bertahap.}
\label{orig03:final:hint:0008}
\textbf{Tahap 1.} Gabungkan $\partial(\tau\lVert W\cdot\rVert_1)$ dengan
kerucut normal $N_C$, lalu gunakan aturan rantai melalui $W^\top$.
\textbf{Tahap 2.} Bandingkan proksimal gabungan dengan proyeksi ke $C$ dan
proksimal konjugat norma $\ell^1$ yang dapat dihitung secara terpisah.

% stable-id: d90.orig.v1.tr03.final.answer.0008
\paragraph{Jawaban singkat.}
\label{orig03:final:answer:0008}
Optimalitas ekuivalen dengan keberadaan $y$ sehingga
\[
 0\in\nabla s(x)+W^\top y+N_C(x),
 \qquad y\in\tau\partial\lVert Wx\rVert_1.
\]
Gradien proksimal langsung sesuai bila
$\operatorname{prox}_{\eta(\tau\lVert W\cdot\rVert_1+\delta_C)}$ mudah;
pemisahan primal--dual sesuai bila hanya proyeksi ke $C$ dan proyeksi dual ke
bola $\ell^\infty$ yang mudah. Residu dapat berupa jarak kedua inklusi ke nol.
Orakel harus tidak bias secara bersyarat, bermomen kedua terkendali, dan
langkah harus memenuhi syarat penjumlahan yang sesuai.

% stable-id: d90.orig.v1.tr03.final.solution.0008
\paragraph{Solusi lengkap.}
\label{orig03:final:solution:0008}
Dengan aturan rantai konveks,
\[
 \partial\bigl(\tau\lVert W\cdot\rVert_1\bigr)(x)
 =W^\top\bigl(\tau\partial\lVert Wx\rVert_1\bigr).
\]
Karena kendala ditulis sebagai indikator $\delta_C$, syarat Fermat adalah
\[
 0\in\nabla s(x)+W^\top
       \bigl(\tau\partial\lVert Wx\rVert_1\bigr)+N_C(x).
\]
Memperkenalkan $y\in\tau\partial\lVert Wx\rVert_1$ menghasilkan sistem pada
jawaban singkat. Secara ekuivalen,
$Wx\in\partial(\tau\lVert\cdot\rVert_1)^*(y)$, dan fungsi konjugat itu adalah
indikator bola $\{y:\lVert y\rVert_\infty\le\tau\}$.

Jika proksimal gabungan
$\tau\lVert W\cdot\rVert_1+\delta_C$ tersedia dalam bentuk tertutup atau dapat
dihitung murah, langkah gradien proksimal stokastik memisahkan bagian halus
$s$ secara langsung. Untuk $W$ umum, proksimal gabungan biasanya tidak
terurai menjadi penyusutan lunak lalu proyeksi. Metode primal--dual menjaga
$Wx$ sebagai operator linear dan hanya membutuhkan proyeksi ke bola $C$ serta
proyeksi $y$ ke bola $\ell^\infty$ berjari-jari $\tau$; karena itu ia lebih
sesuai pada keadaan tersebut.

Satu pasangan residu matematis adalah
\[
 r_{\mathrm p}(x,y)=
 \operatorname{dist}\bigl(0,\nabla s(x)+W^\top y+N_C(x)\bigr),
\]
\[
 r_{\mathrm d}(x,y)=
 \operatorname{dist}\bigl(Wx,\partial(\tau\lVert\cdot\rVert_1)^*(y)\bigr),
\]
dengan nilai tak hingga bila $y$ berada di luar domain konjugat. Keduanya nol
tepat pada pasangan KKT\@. Untuk analisis stokastik, penduga gradien
$G_k$ sekurang-kurangnya harus memenuhi
$\mathbb E[G_k\mid\mathcal F_k]=\nabla s(x_k)$ dan batas momen kedua atau
varians yang sesuai. Pada skema Robbins--Monro dasar, langkah positif lazimnya
memenuhi $\sum_k\eta_k=\infty$ dan $\sum_k\eta_k^2<\infty$; klaim laju dengan
langkah tetap memerlukan asumsi tambahan seperti reduksi varians atau menerima
lantai galat. Norma operator $W$ juga harus masuk ke batas langkah
primal--dual, bukan diabaikan.

% OR03-S012 | lapisan asli: laboratorium 3 globalisasi Newton
% segment-id: d90.orig.v1.tr03.seg0012
% lab-id: d90.orig.v1.tr03.lab.0003
% lab-hint-id: d90.orig.v1.tr03.lab.hint.0003
% lab-answer-id: d90.orig.v1.tr03.lab.answer.0003
% lab-solution-id: d90.orig.v1.tr03.lab.solution.0003

\section{Laboratorium 3: kegagalan lokal dan globalisasi Newton}
\label{orig03:lab:globalisasi-newton}

Laboratorium ini memperlihatkan bahwa arah yang tampak ``orde dua'' belum
tentu aman secara global. Kita memakai fungsi konveks mulus
\begin{equation}
  F(x,y)=\log(\cosh x)+\frac12y^2,
  \qquad
  \nabla F(x,y)=\begin{pmatrix}\tanh x\\y\end{pmatrix},
  \qquad
  \nabla^2F(x,y)=
  \begin{pmatrix}\operatorname{sech}^2x&0\\0&1\end{pmatrix}.
  \label{orig03:eq:lab3-model}
\end{equation}
Minimum uniknya adalah $(0,0)$ dengan nilai nol. Walaupun Hessian selalu
positif, kelengkungan pada arah $x$ sangat kecil ketika $|x|$ besar. Karena
itu, langkah Newton penuh dari titik awal yang disediakan dapat melompat ke
daerah dengan nilai objektif jauh lebih tinggi. Contoh ini memisahkan dua
persoalan: pemilihan arah dan pemilihan panjang langkah.

Kode tetap berada di
\texttt{labs/original-03/globalisasi-newton.py}. Dengan seed dan konfigurasi
yang dibekukan, kode membandingkan empat metode:
\begin{enumerate}
  \item gradien dengan langkah tetap yang sengaja terlalu besar;
  \item gradien dengan pencarian balik Armijo;
  \item Newton murni dengan langkah penuh;
  \item Newton teredam dengan Hessian terkoreksi
    $B_k=\nabla^2F(z_k)+\tau_kI$, dengan
    $\lambda_{\min}(B_k)\geq \underline\mu$.
\end{enumerate}
Jejak lengkap ditulis ke CSV, ringkasan dan sertifikat ditulis ke JSON, dan
SVG hanya memvisualkan data yang sama. Tidak ada kesimpulan numerik yang
bergantung pada grafik.

\subsection*{Prompt}
\begin{exercise}[Tugas laboratorium]
Jalankan skrip sekali dan kerjakan hal-hal berikut.
\begin{enumerate}
  \item Turunkan langkah Newton murni pada koordinat $x$ dan jelaskan mengapa
    $x_0\approx3$ menghasilkan lompatan besar walaupun fungsi konveks.
  \item Dari jejak CSV, temukan kenaikan objektif pertama bagi Newton murni
    dan kegagalan metode gradien berlangkah tetap.
  \item Buktikan bahwa syarat Armijo
    \[
      F(z_k+\alpha_kd_k)\leq F(z_k)+c_1\alpha_k
      \nabla F(z_k)^\top d_k
    \]
    menghasilkan penurunan ketat bila $d_k$ adalah arah turun.
  \item Buktikan bahwa koreksi
    $\tau_k=\max\{0,\underline\mu-\lambda_{\min}(\nabla^2F(z_k))\}$
    membuat $d_k=-B_k^{-1}\nabla F(z_k)$ menjadi arah turun selama
    gradien tidak nol.
  \item Cocokkan bukti Anda dengan seluruh boolean sertifikat dalam JSON dan
    jelaskan mengapa ``gagal'' di sini merupakan bukti terkontrol, bukan
    galat program.
\end{enumerate}
\end{exercise}

\subsection*{Petunjuk bertahap}
\label{orig03:lab:hint:0003}
\begin{enumerate}
  \item Gunakan $\operatorname{sech}^2x=1-\tanh^2x$ dan tulis
    $x^+=x-\tanh(x)/\operatorname{sech}^2(x)$.
  \item Bandingkan kolom \texttt{objective} pada dua baris berturutan; jangan
    menyimpulkan kegagalan hanya dari gambar.
  \item Pada ruas kanan Armijo, tanda
    $\nabla F(z_k)^\top d_k<0$ adalah fakta penentu.
  \item Gunakan $B_k\succeq\underline\mu I$ untuk menunjukkan
    $\nabla F(z_k)^\top d_k=-g_k^\top B_k^{-1}g_k<0$.
  \item Periksa juga bahwa semua langkah metode teredam diterima oleh Armijo
    dan nilai objektifnya monoton menurun.
\end{enumerate}

\subsection*{Jawaban singkat}
\label{orig03:lab:answer:0003}
Newton murni gagal karena invers kelengkungan
$1/\operatorname{sech}^2(x_0)$ sangat besar. Langkah gradien tetap gagal
karena faktor koordinat kuadratiknya adalah $1-\alpha=-1.5$. Pencarian balik
memendekkan langkah sampai memenuhi penurunan cukup. Koreksi Hessian menjamin
arah turun, sedangkan Armijo mengglobalisasi panjang langkah; kombinasi
keduanya mencapai toleransi gradien yang dibekukan.

\subsection*{Solusi acuan lengkap}
\label{orig03:lab:solution:0003}
Pada koordinat pertama, Newton murni memberi
\[
 x^+=x-\frac{\tanh x}{\operatorname{sech}^2x}
     =x-\sinh x\cosh x.
\]
Untuk $x\approx3$, suku $\sinh x\cosh x$ berorde $10^2$, sehingga iterasi
melompat melintasi minimizer dan menaikkan $\log(\cosh x)$. Pada iterasi
berikutnya $\operatorname{sech}^2x$ mendekati nol pada presisi mesin; skrip
berhenti sebelum pembagian hampir singular merusak jejak. Untuk gradien tetap
pada koordinat kedua berlaku $y^+=(1-2.5)y=-1.5y$, jadi magnitudonya tumbuh
geometrik.

Untuk arah gradien $d=-g$, kita memiliki $g^\top d=-\|g\|^2<0$. Karena $F$
diferensiabel, terdapat $\bar\alpha>0$ sehingga setiap
$0<\alpha\leq\bar\alpha$ memenuhi Armijo; pencarian balik dengan faktor dalam
$(0,1)$ karena itu berhenti dalam banyak langkah berhingga. Untuk arah Newton
terkoreksi, konstruksi $B_k\succeq\underline\mu I$ memberi
\[
 g_k^\top d_k=-g_k^\top B_k^{-1}g_k
 \leq-\frac{\|g_k\|^2}{\lambda_{\max}(B_k)}<0.
\]
Maka argumen pencarian balik yang sama berlaku. Sertifikat program memeriksa
langsung batas eigenvalue, tanda turunan arah, semua penerimaan Armijo,
monotonisitas objektif, dan norma gradien akhir. Sebaliknya, metode yang tidak
diglobalisasi diwajibkan memperlihatkan kenaikan atau ledakan yang telah
ditentukan; dengan demikian jalur gagal merupakan bagian dari eksperimen.

\paragraph{Hak dan firewall O018.}
Teks, soal, solusi, dan kode laboratorium ini ditulis mandiri untuk lapisan
Original-03 dan tersedia berdasarkan CC BY-SA 4.0. Tidak ada prosa, soal,
solusi, kode, atau susunan pedagogis O018 yang disalin atau diadaptasi; O018
tetap merupakan korpus terpisah.

% OR03-S013 | lapisan asli: laboratorium 4 transportasi entropik
% segment-id: d90.orig.v1.tr03.seg0013
% lab-id: d90.orig.v1.tr03.lab.0004
% lab-hint-id: d90.orig.v1.tr03.lab.hint.0004
% lab-answer-id: d90.orig.v1.tr03.lab.answer.0004
% lab-solution-id: d90.orig.v1.tr03.lab.solution.0004

\section{Laboratorium 4: Sinkhorn log-domain dan sertifikat transportasi}
\label{orig03:lab:transportasi-entropik}

Diberikan histogram positif $a\in\mathbb{R}^m$, $b\in\mathbb{R}^n$ dengan
$\1^\top a=\1^\top b=1$ dan matriks biaya $C$, masalah
transportasi teratur-entropi yang dipakai di sini adalah
\begin{equation}
 \min_{P\1=a,\;P^\top\1=b,\;P\geq0}
 \langle C,P\rangle+\varepsilon\sum_{ij}P_{ij}(\log P_{ij}-1).
 \label{orig03:eq:sinkhorn-primal}
\end{equation}
Representasi langsung $K_{ij}=\exp(-C_{ij}/\varepsilon)$ rentan underflow.
Skrip \texttt{labs/original-03/transportasi-entropik.py} menyimpan potensial
dual $f,g$ dan melakukan pembaruan dengan \texttt{logsumexp}. Kopling dibentuk
hanya dari
\[
 \log P_{ij}=\frac{f_i+g_j-C_{ij}}{\varepsilon}.
\]

Eksperimen utama memakai titik dan massa yang dihasilkan oleh seed tetap.
Dua eksperimen ekuivalen kemudian dijalankan: (i) $C$ dan $\varepsilon$
dikalikan faktor positif yang sama; (ii) konstanta besar ditambahkan ke semua
entri $C$. Transformasi pertama mempertahankan rasio biaya-entropi, sedangkan
transformasi kedua menambah konstanta pada nilai objektif tanpa mengubah
minimizer. Kasus kedua memaksa seluruh kernel naif menjadi nol, tetapi
pembaruan log-domain tetap menghasilkan kopling yang sama.

\subsection*{Prompt}
\begin{exercise}[Tugas laboratorium]
\begin{enumerate}
  \item Turunkan pembaruan bergantian bagi $f$ dan $g$ dari kendala marginal.
  \item Buktikan bentuk dual dan identitas celah primal-dual yang diperiksa
    skrip.
  \item Verifikasi dari CSV bahwa setiap entri kopling utama positif dan
    bahwa jumlah baris serta kolom cocok dengan $a$ dan $b$ dalam toleransi.
  \item Buktikan invariansi kopling ketika $(C,\varepsilon)$ diganti dengan
    $(sC,s\varepsilon)$, $s>0$.
  \item Jelaskan mengapa $C\mapsto C+\gamma\1\1^\top$ tidak
    mengubah kopling, lalu cocokkan dengan uji underflow terkontrol.
\end{enumerate}
\end{exercise}

\subsection*{Petunjuk bertahap}
\label{orig03:lab:hint:0004}
\begin{enumerate}
  \item Ambil log dari persamaan $P_{ij}=\exp((f_i+g_j-C_{ij})/\varepsilon)$
    dan selesaikan satu potensial ketika potensial lain dibekukan.
  \item Gunakan konjugat dari $p\mapsto\varepsilon p(\log p-1)$, yaitu
    $q\mapsto\varepsilon\exp(q/\varepsilon)$.
  \item Jangan menjumlahkan kernel naif pada uji stres; periksa penghitung nol
    dan residual hasil log-domain secara terpisah.
  \item Untuk penskalaan, pilih potensial baru $(sf,sg)$.
  \item Untuk penambahan konstanta, serap $\gamma$ ke salah satu potensial.
\end{enumerate}

\subsection*{Jawaban singkat}
\label{orig03:lab:answer:0004}
Pembaruan log-domain menormalkan satu marginal pada satu waktu tanpa pernah
membentuk kernel kecil. Dualnya adalah
\[
 \max_{f,g}\ a^\top f+b^\top g-
 \varepsilon\sum_{ij}\exp((f_i+g_j-C_{ij})/\varepsilon).
\]
CSV dan JSON membuktikan residual marginal, positivitas numerik, kesepakatan
primal-dual, invariansi penskalaan, invariansi pergeseran biaya, dan pemulihan
dari kasus ketika kernel naif seluruhnya underflow.

\subsection*{Solusi acuan lengkap}
\label{orig03:lab:solution:0004}
Untuk $g$ tetap, kendala baris memberi
\[
 a_i=\sum_j\exp((f_i+g_j-C_{ij})/\varepsilon),
 \quad
 f_i=\varepsilon\left[\log a_i-
 \operatorname{LSE}_j((g_j-C_{ij})/\varepsilon)\right].
\]
Rumus bagi $g_j$ analog dari kendala kolom. Konjugasi entropi menghasilkan
dual pada jawaban singkat. Jika $P$ dibentuk dari potensial, nilai primal
\[
 p=\langle C,P\rangle+\varepsilon\sum_{ij}P_{ij}(\log P_{ij}-1)
\]
dan nilai dual
\[
 d=a^\top f+b^\top g-\varepsilon\sum_{ij}P_{ij}
\]
berimpit pada optimum; program memeriksa $|p-d|$.

Jika $C'=sC$ dan $\varepsilon'=s\varepsilon$, potensial
$f'=sf$, $g'=sg$ memberi log-kopling yang persis sama. Jika
$C'=C+\gamma\1\1^\top$, pilih $f'=f+\gamma\1$ dan $g'=g$;
sekali lagi log-kopling tidak berubah. Namun kernel naif memuat faktor
$\exp(-\gamma/\varepsilon)$ dan dapat menjadi nol pada aritmetika floating
point. Potensial log-domain menyerap konstanta sebelum eksponensiasi akhir,
sehingga residual marginal, positivitas kopling utama, dan kesepakatan dengan
dua masalah ekuivalen tetap dapat diperiksa.

\paragraph{Hak dan firewall O018.}
Teks, soal, solusi, dan kode ini merupakan material baru CC BY-SA 4.0. Tidak
ada byte, prosa, latihan, solusi, kode, atau struktur instruksional O018 yang
digunakan; peran laboratorium ini terbatas pada penutupan komputasional
Original-03.

% OR03-S014 | lapisan asli: proyek kapstone masalah invers komposit
% segment-id: d90.orig.v1.tr03.seg0014
% capstone-id: d90.orig.v1.tr03.capstone.0001
% capstone-unit-id: d90.orig.v1.tr03.capstone.unit
% capstone-milestone-id: d90.orig.v1.tr03.capstone.milestone.0001
% capstone-milestone-id: d90.orig.v1.tr03.capstone.milestone.0002
% capstone-milestone-id: d90.orig.v1.tr03.capstone.milestone.0003
% capstone-milestone-id: d90.orig.v1.tr03.capstone.milestone.0004
% capstone-milestone-id: d90.orig.v1.tr03.capstone.milestone.0005
% capstone-milestone-id: d90.orig.v1.tr03.capstone.milestone.0006
% capstone-milestone-id: d90.orig.v1.tr03.capstone.milestone.0007
% capstone-hint-id: d90.orig.v1.tr03.capstone.hint.0001
% capstone-answer-id: d90.orig.v1.tr03.capstone.answer.0001
% capstone-solution-id: d90.orig.v1.tr03.capstone.solution.0001

\section{Proyek kapstone: masalah invers komposit yang tahan pencilan}
\label{orig03:capstone:invers-komposit}
\label{orig03:capstone:unit}

Proyek ini menyatukan pemodelan konveks, gradien Lipschitz, operator
proksimal, formulasi dual, pemisahan primal-dual, anggaran oracle, dan
sertifikat a posteriori. Instans tetap dibentuk oleh
\texttt{labs/original-03/kapstone-invers-komposit.py}: matriks
$A\in\mathbb{R}^{72\times48}$, sinyal renggang $x^\natural$, derau kecil,
dan delapan pencilan besar. Modelnya
\begin{equation}
 \min_x\; \Phi(x):=
 \sum_{i=1}^{72}h_\delta((Ax-b)_i)+\lambda\|x\|_1,
 \qquad
 h_\delta(r)=
 \begin{cases}
   \frac12r^2,&|r|\leq\delta,\\
   \delta(|r|-\frac12\delta),&|r|>\delta.
 \end{cases}
 \label{orig03:eq:capstone-model}
\end{equation}
FISTA dengan restart gradien dan PDHG dijalankan untuk jumlah iterasi yang sama.
Setiap iterasi inti masing-masing memakai satu aplikasi $A$ dan satu aplikasi
$A^\top$, sehingga perbandingan memakai anggaran operator yang sama. Jejak
CSV mencatat objektif, batas bawah dual layak, celah tersertifikasi, dan galat
rekonstruksi. SVG hanya merupakan visualisasi jejak tersebut.

\subsection*{Prompt}
\begin{exercise}[Tugas proyek dan tujuh tonggak]
Selesaikan tujuh tonggak berikut dan serahkan uraian matematis bersama JSON,
CSV, serta SVG yang dihasilkan tanpa perubahan konfigurasi.
\begin{enumerate}
  \item \label{orig03:capstone:milestone:0001}\textbf{Instans dan pencilan.} Rekonstruksi instans dari seed,
    verifikasi ukuran, support sinyal, dan delapan indeks pencilan.
  \item \label{orig03:capstone:milestone:0002}\textbf{Model konveks.} Buktikan kekonveksan \eqref{orig03:eq:capstone-model}
    dan jelaskan mengapa Huber lebih tahan pencilan daripada kuadrat murni.
  \item \label{orig03:capstone:milestone:0003}\textbf{Bagian mulus.} Turunkan gradien
    $A^\top\operatorname{clip}(Ax-b,-\delta,\delta)$ dan batas Lipschitz
    $L=\|A\|_2^2$.
  \item \label{orig03:capstone:milestone:0004}\textbf{Langkah proksimal.} Turunkan soft-threshold bagi
    $\lambda\|x\|_1$ dan tulis iterasi FISTA yang dipakai.
  \item \label{orig03:capstone:milestone:0005}\textbf{Dual dan PDHG.} Turunkan konjugat Huber, kendala dual, serta
    pembaruan PDHG; buktikan $\tau\sigma\|A\|_2^2<1$.
  \item \label{orig03:capstone:milestone:0006}\textbf{Anggaran sepadan.} Bandingkan kedua metode pada jumlah
    perkalian matriks-vektor yang sama, bukan pada waktu dinding.
  \item \label{orig03:capstone:milestone:0007}\textbf{Sertifikat dan interpretasi.} Verifikasi kelayakan dual,
    celah primal-dual, pemetaan gradien proksimal, dan perbaikan terhadap
    least squares; nyatakan juga apa yang tidak dibuktikan eksperimen ini.
\end{enumerate}
\end{exercise}

\subsection*{Petunjuk bertahap}
\label{orig03:capstone:hint:0001}
\begin{enumerate}
  \item Hash matriks dan data dalam JSON adalah identitas instans, bukan
    pengganti pemeriksaan ukuran dan seed.
  \item Turunan Huber adalah proyeksi skalar ke $[-\delta,\delta]$ dan
    bersifat satu-Lipschitz.
  \item Komposisi dengan $A$ mengalikan konstanta Lipschitz oleh
    $\|A\|_2^2$.
  \item Gunakan
    $\operatorname{prox}_{t\lambda\|\cdot\|_1}(v)
     =\operatorname{sign}(v)(|v|-t\lambda)_+$.
  \item Konjugat Huber adalah $\frac12u^2+\iota_{[-\delta,\delta]}(u)$;
    konjugat norma satu menghasilkan kendala
    $\|A^\top u\|_\infty\leq\lambda$.
  \item Kolom \texttt{matrix\_vector\_products} adalah sumbu perbandingan
    yang benar.
  \item Skala kandidat dual bila perlu untuk memperoleh batas bawah yang
    benar-benar layak sebelum menghitung celah.
\end{enumerate}

\subsection*{Jawaban singkat}
\label{orig03:capstone:answer:0001}
Model adalah jumlah fungsi konveks. Bagian Huber mulus memiliki gradien
$L$-Lipschitz dengan $L=\|A\|_2^2$, dan bagian norma satu memiliki proksimal
soft-threshold. Dualnya memaksimalkan
\[
 -b^\top u-\frac12\|u\|^2
 \quad\text{dengan}\quad
 \|u\|_\infty\leq\delta,
 \quad \|A^\top u\|_\infty\leq\lambda.
\]
FISTA dan PDHG mencapai objektif yang cocok pada anggaran operator sama;
batas dual layak memberi celah terukur, dan kedua rekonstruksi mengungguli
least squares pada instans berpencilan yang dibekukan.

\subsection*{Solusi acuan lengkap}
\label{orig03:capstone:solution:0001}
Karena $h_\delta$ dan norma satu konveks, komposisi afin dan penjumlahan
mempertahankan kekonveksan. Turunan Huber adalah
$h_\delta'(r)=\operatorname{clip}(r,-\delta,\delta)$. Proyeksi ini
satu-Lipschitz, sehingga
\[
 \|\nabla f(x)-\nabla f(y)\|
 \leq\|A\|_2^2\|x-y\|.
\]
Langkah $t<1/L$ yang dibekukan karena itu sah. Optimalisasi proksimal skalar
memberi soft-threshold, lalu ekstrapolasi Nesterov menghasilkan FISTA; restart
gradien memakai vektor yang sudah tersedia tanpa mengubah dua aplikasi operator inti per
iterasi.

Dari
$h_\delta^*(u)=\frac12u^2+\iota_{[-\delta,\delta]}(u)$ dan
$(\lambda\|\cdot\|_1)^*(v)=\iota_{\{\|v\|_\infty\leq\lambda\}}(v)$,
dual yang tertulis pada jawaban singkat mengikuti melalui Fenchel--Rockafellar.
Proksimal $h_\delta^*$ adalah pembagian oleh $1+\sigma$ diikuti kliping ke
$[-\delta,\delta]$; proksimal primal tetap soft-threshold. Pilihan
$\tau=\sigma=0.99/\|A\|_2$ memberi
$\tau\sigma\|A\|_2^2=0.99^2<1$.

Pada iterasi berhingga, kandidat dual dapat melanggar kendala kedua sedikit.
Penskalaan
\[
 \widehat u=\rho u,
 \qquad
 \rho=\min\left\{1,
 \frac{\lambda}{\|A^\top u\|_\infty}\right\}
\]
mempertahankan kendala kotak dan memastikan kelayakan. Nilai dual
$-b^\top\widehat u-\|\widehat u\|^2/2$ menjadi batas bawah rigor bagi
objektif primal. Program menguji kelayakan, celah terhadap batas terbaik,
norma pemetaan gradien proksimal FISTA, kesamaan anggaran, kedekatan kedua
objektif, dan galat rekonstruksi terhadap baseline least squares. Bukti ini
berlaku untuk instans dan toleransi yang dinyatakan; ia tidak mengklaim bahwa
parameter yang sama optimal bagi semua masalah invers atau semua pola
pencilan.

\paragraph{Hak dan firewall O018.}
Proyek, bukti, soal, solusi, kode, dan instans sintetik ini ditulis mandiri dan
tersedia berdasarkan CC BY-SA 4.0. Tidak ada materi O018 yang disalin,
diadaptasi, atau dijadikan templat tersembunyi; O018 tetap terpisah dari
lapisan penutupan Original-03.

